% Leçon 26 — Grammaire : L’addition et soustraction en chinois 加法和减法 ; les différentes monnaies ; Les nombres ; 除了……以外，…… — Chinois Tle A — Cours signé OnBuch+
% Programme officiel MINESEC. Compiler avec : tectonic ZH26-grammaire-laddition-et-soustraction-en-chinois-les.tex
\newcommand{\DOCMATIERE}{Chinois}
\newcommand{\DOCNIVEAU}{Tle A}
\newcommand{\DOCMODULE}{Module 4 : Activités économiques et monde du travail}
\newcommand{\DOCLECON}{ZH26}
\newcommand{\DOCTITRE}{Leçon 26 — Grammaire : L’addition et soustraction en chinois 加法和减法 ; les différentes monnaies ; Les nombres ; 除了……以外，……}
% OnBuch+ — préambule « cours signés » ENRICHI (compilé avec tectonic / XeLaTeX)
% Système visuel « Pop » d'OnBuch+ : fond crème (jamais blanc), bordures encre
% épaisses, ombres « dures » décalées, titres Archivo Black en capitales.
% Version MATHÉMATIQUES : graphes (pgfplots/tikz), boîtes pédagogiques
% (définition, propriété, méthode, attention, le savais-tu, exercice résolu,
% exercice, TP, bilan), page de garde et signature OnBuch+ ancrée en bas.
%
% Macros attendues dans main.tex AVANT \input{preamble} :
%   \DOCMATIERE  \DOCNIVEAU  \DOCMODULE  \DOCLECON (numéro)  \DOCTITRE
\documentclass[11pt]{article}

\usepackage{fontspec}
\usepackage[french]{babel} % français = langue principale ; le chinois via \zh{..} / textebox (xeCJK)
\usepackage{xeCJK}
\usepackage{pifont} % \ding{51} \ding{55} utilisés par les modèles
\usepackage[a4paper,margin=2cm,top=3.4cm,bottom=2.8cm,headheight=1.4cm,footskip=1.3cm]{geometry}
\usepackage{amsmath,amssymb,amsfonts}
\usepackage{mathtools} % \xmapsto, \coloneqq... souvent utilisés par les modèles
\usepackage{cancel} % simplifications barrées (bilans de forces)
% \abs{z} (module, valeur absolue) et \norm{u} (norme), utilisés par les modèles
\providecommand{\abs}{}\let\abs\relax\DeclarePairedDelimiter{\abs}{\lvert}{\rvert}
\providecommand{\norm}{}\let\norm\relax\DeclarePairedDelimiter{\norm}{\lVert}{\rVert}
\usepackage[table]{xcolor} % [table] : fournit aussi \rowcolors (alternance de lignes)
\usepackage{pagecolor}
\usepackage{tikz}
\usetikzlibrary{arrows.meta,positioning,calc,shapes.geometric,shapes.misc,shapes.arrows,decorations.pathmorphing,decorations.pathreplacing,decorations.markings,patterns,shadows,mindmap,trees,fit,backgrounds,matrix,intersections,angles,quotes}
\usepackage{pgfplots}
\usepackage[version=4]{mhchem} % équations nucléaires \ce{^{235}_{92}U}
\usepackage[european,siunitx]{circuitikz} % schémas électriques
\pgfplotsset{compat=1.18}
\usepgfplotslibrary{fillbetween,groupplots} % aires entre courbes (calcul intégral)
\usepackage{enumitem}
\usepackage{fancyhdr}
% [most] charge toutes les bibliothèques tcolorbox (listings, minted, raster,
% documentation...) et gonfle inutilement la mémoire TeX sur un document long
% (~300 boîtes) : on ne charge que ce qui est réellement utilisé.
\usepackage[skins,breakable]{tcolorbox}
\usepackage{graphicx}
\usepackage{colortbl}
\usepackage{booktabs}
\usepackage{tabularx}
\usepackage{diagbox} % case d'en-tête diagonale des tableaux à double entrée
% Colonnes extensibles centrée / gauche / droite (C, L, R), souvent utilisées par les modèles
\newcolumntype{C}{>{\centering\arraybackslash}X}
\newcolumntype{L}{>{\raggedright\arraybackslash}X}
\newcolumntype{R}{>{\raggedleft\arraybackslash}X}
\usepackage{array}
\usepackage{multicol}
\usepackage{multirow}
\usepackage{siunitx}
% parse-numbers=false : accepte aussi \SI{2,5 \times 10^{-3}}{...}
\sisetup{locale=FR,output-decimal-marker={,},group-minimum-digits=4,parse-numbers=false}
\DeclareSIUnit\ohm{\ensuremath{\symup{\Omega}}} % U+2126 absent de la police
\DeclareSIUnit\division{div} % réglages d'oscilloscope (ms/div, V/div)
\DeclareSIUnit\tr{tr} % tours (vitesse de rotation, ex. \SI{1500}{\tr\per\minute})
\DeclareSIUnit\molar{M} % concentration molaire (ex. \SI{3}{\molar}, \SI{1}{\micro\molar})
\DeclareSIUnit\an{an} % année (le modèle confond parfois avec \year, un compteur TeX) — ex. \SI{5}{\kilo\meter\per\an}
\DeclareSIUnit\FCFA{FCFA} % franc CFA (ex. \SI{500}{\FCFA\per\kilo\gram})
\DeclareSIUnit\degreeBrix{\textdegree Brix} % teneur en sucre d'un sirop/jus (réfractométrie)
\DeclareSIUnit\month{mois} % le modèle utilise parfois \month, un compteur TeX, comme unité de durée
\usepackage{qrcode}
% \degree : commande fréquemment utilisée par les modèles pour un angle ou une
% température en dehors de \SI{}{} (non fournie par les paquets chargés).
\providecommand{\degree}{^{\circ}}
% \qed (fin de démonstration), utilisé par les modèles sans amsthm
\providecommand{\qed}{\hfill$\square$}
% \barre : double barre des valeurs interdites dans un tableau de variations (tkz-tab)
\providecommand{\barre}{\ensuremath{\|}}
% environnement proof (amsthm non chargé) utilisé par les modèles
\ifdefined\proof\else\newenvironment{proof}[1][Démonstration]{\par\noindent\textit{#1.}\ }{\qed\par}\fi
\usepackage{lastpage}
\usepackage{hyperref}
\hypersetup{hidelinks}

% ============================================================
% Palette « Pop » OnBuch+ — mêmes hex que la classe P (plus_ui.dart)
% ============================================================
\definecolor{poporange}{HTML}{F59321}
\definecolor{popdark}{HTML}{E07A0C}
\definecolor{popcream}{HTML}{FFF6E8}
\definecolor{popink}{HTML}{1C1714}
\definecolor{poppurple}{HTML}{7A5AE0}
\definecolor{popgreen}{HTML}{2E9E6A}
\definecolor{poppink}{HTML}{E0457B}
\definecolor{popblue}{HTML}{2F7DE1}
\definecolor{popgold}{HTML}{C98A12}
\definecolor{popmuted}{HTML}{8A7F76}
\definecolor{popline}{HTML}{EADFD2}
\definecolor{popgray}{HTML}{8A7F76}
\colorlet{popgrey}{popgray}
% Alias tolérants (le modèle nomme parfois les couleurs différemment)
\colorlet{popwhite}{white}
\colorlet{popblack}{popink}
\colorlet{popred}{poppink}
\colorlet{popyellow}{popgold}
% Teintes claires (fonds de tableaux/encadrés) — le modèle nomme parfois
% les couleurs avec un suffixe L (ex. popblueL) sans les avoir définies.
\colorlet{poporangeL}{poporange!20}
\colorlet{popdarkL}{popdark!20}
\colorlet{poppurpleL}{poppurple!20}
\colorlet{popgreenL}{popgreen!20}
\colorlet{poppinkL}{poppink!20}
\colorlet{popblueL}{popblue!20}
\colorlet{popgoldL}{popgold!20}
\colorlet{popredL}{popred!20}
\colorlet{popgrayL}{popgray!20}
% Teintes claires pour fonds de tableaux / zones
\colorlet{popblueL}{popblue!10!white}
\colorlet{poporangeL}{poporange!14!white}
\colorlet{popgreenL}{popgreen!12!white}
\colorlet{poppurpleL}{poppurple!12!white}
\colorlet{poppinkL}{poppink!12!white}
\colorlet{popgoldL}{popgold!14!white}

\pagecolor{popcream}

% ============================================================
% Typographie
% ============================================================
\defaultfontfeatures{Ligatures=TeX}
\setmainfont{PlusJakartaSans-Medium.ttf}[Path=../../../fonts/,BoldFont=PlusJakartaSans-Bold.ttf]
% Symboles, API et emojis absents de la police principale : repli sur DejaVu Sans (caractères actifs)
\newfontface\symfont{DejaVuSans.ttf}[Path=../../../fonts/]
\makeatletter
\newcommand\symactive[1]{\catcode#1=13 \begingroup\lccode`\~=#1 \lowercase{\endgroup\def~}{{\symfont\char#1}}}
\newcommand\symblank[1]{\catcode#1=13 \begingroup\lccode`\~=#1 \lowercase{\endgroup\def~}{}}
\newcommand\symhyph[1]{\catcode#1=13 \begingroup\lccode`\~=#1 \lowercase{\endgroup\def~}{\mbox{-}}}
\newcount\symc
\newcommand\symrange[2]{\symc=#1 \@whilenum\symc<#2 \do{\expandafter\symactive\expandafter{\the\symc}\advance\symc by 1 }}
\newcommand\symblankrange[2]{\symc=#1 \@whilenum\symc<#2 \do{\expandafter\symblank\expandafter{\the\symc}\advance\symc by 1 }}
\makeatother
\symrange{"0250}{"02FF}\symactive{"2713}\symactive{"2714}\symactive{"2717}\symactive{"2718}\symactive{"2610}\symactive{"2611}\symactive{"2612}
\symrange{"2600}{"27BF}
\catcode"2705=13 \begingroup\lccode`\~="2705 \lowercase{\endgroup\def~}{{\symfont\char"2713}}\symblank{"26BD}\symblank{"26A1}
\symrange{"2460}{"257F}\symactive{"223C}\symactive{"2248}\symactive{"2260}
\symactive{"01F8}\symactive{"01F9}\symactive{"1E3E}\symactive{"1E3F}
\symhyph{"2011}\symhyph{"2E1A}\symblankrange{"1F300}{"1FAFF}\symblank{"FE0F}
% Caractères chinois : WenQuanYi Zen Hei (sous-ensemble GB2312 + pinyin), gras/italique simulés
\setCJKmainfont{WenQuanYiZenHei-subset.ttf}[Path=../../../fonts/,AutoFakeBold=3,AutoFakeSlant=0.2]
\xeCJKsetup{CJKspace=false}
% Pinyin : voyelles à caron (ǎ ǐ ǒ ǔ ǖ ǘ ǚ ǜ) absentes de la police principale -> caractères actifs
\newfontface\pyfont{WenQuanYiZenHei-subset.ttf}[Path=../../../fonts/]
\newcommand\pyactive[1]{\catcode#1=13 \begingroup\lccode`\~=#1 \lowercase{\endgroup\def~}{{\pyfont\char#1}}}
\pyactive{"01CD}\pyactive{"01CE}\pyactive{"01CF}\pyactive{"01D0}\pyactive{"01D1}\pyactive{"01D2}\pyactive{"01D3}\pyactive{"01D4}\pyactive{"01D5}\pyactive{"01D6}\pyactive{"01D7}\pyactive{"01D8}\pyactive{"01D9}\pyactive{"01DA}\pyactive{"01DB}\pyactive{"01DC}
% Maths en sans-serif (Fira Math) avec chiffres et lettres droites de Plus
% Jakarta, pour que les formules se fondent dans le texte.
\usepackage[math-style=ISO,bold-style=ISO]{unicode-math}
\setmathfont{FiraMath-Regular.otf}[Path=../../../fonts/]
\setmathfont{PlusJakartaSans-Medium.ttf}[Path=../../../fonts/,range={up/num,up/{Latin,latin},\mathup}]
\usepackage{icomma} % virgule décimale sans espace en mode math
% Caractères mathématiques Unicode absents des polices de texte (Plus Jakarta,
% Archivo Black) : rendus par leur équivalent mathématique, en texte comme en
% formule — sinon ils s'affichent en carré vide (ex. « Arithmétique dans ℤ »).
\usepackage{newunicodechar}
\newunicodechar{ℝ}{\ensuremath{\mathbb{R}}}
\newunicodechar{ℤ}{\ensuremath{\mathbb{Z}}}
\newunicodechar{ℂ}{\ensuremath{\mathbb{C}}}
\newunicodechar{ℕ}{\ensuremath{\mathbb{N}}}
\newunicodechar{ℚ}{\ensuremath{\mathbb{Q}}}
\newunicodechar{𝔻}{\ensuremath{\mathbb{D}}}
\newunicodechar{α}{\ensuremath{\alpha}}
\newunicodechar{β}{\ensuremath{\beta}}
\newunicodechar{γ}{\ensuremath{\gamma}}
\newunicodechar{δ}{\ensuremath{\delta}}
\newunicodechar{ε}{\ensuremath{\varepsilon}}
\newunicodechar{θ}{\ensuremath{\theta}}
\newunicodechar{λ}{\ensuremath{\lambda}}
\newunicodechar{μ}{\ensuremath{\mu}}
\newunicodechar{σ}{\ensuremath{\sigma}}
\newunicodechar{τ}{\ensuremath{\tau}}
\newunicodechar{φ}{\ensuremath{\varphi}}
\newunicodechar{ψ}{\ensuremath{\psi}}
\newunicodechar{ω}{\ensuremath{\omega}}
\newunicodechar{Σ}{\ensuremath{\Sigma}}
% majuscules produites par \MakeUppercase dans les titres (α-aminés -> Α)
\newunicodechar{Α}{\ensuremath{\alpha}}
\newunicodechar{Β}{\ensuremath{\beta}}
\newunicodechar{Γ}{\ensuremath{\Gamma}}
\newunicodechar{Δ}{\ensuremath{\Delta}}
\newunicodechar{Ε}{\ensuremath{\varepsilon}}
\newunicodechar{Θ}{\ensuremath{\Theta}}
\newunicodechar{Λ}{\ensuremath{\Lambda}}
\newunicodechar{Μ}{\ensuremath{\mu}}
\newunicodechar{Τ}{\ensuremath{\tau}}
\newunicodechar{Φ}{\ensuremath{\Phi}}
\newunicodechar{Ψ}{\ensuremath{\Psi}}
\newunicodechar{Ω}{\ensuremath{\Omega}}
\newunicodechar{↦}{\ensuremath{\mapsto}}
\newunicodechar{∈}{\ensuremath{\in}}
\newunicodechar{∉}{\ensuremath{\notin}}
\newunicodechar{∧}{\ensuremath{\wedge}}
\newunicodechar{⇔}{\ensuremath{\Leftrightarrow}}
\newunicodechar{⟺}{\ensuremath{\Longleftrightarrow}}
\newunicodechar{⇒}{\ensuremath{\Rightarrow}}
\newunicodechar{⟹}{\ensuremath{\Longrightarrow}}
\newunicodechar{∘}{\ensuremath{\circ}}
\newunicodechar{∪}{\ensuremath{\cup}}
\newunicodechar{∩}{\ensuremath{\cap}}
\newunicodechar{≡}{\ensuremath{\equiv}}
\newunicodechar{⊂}{\ensuremath{\subset}}
\newunicodechar{⊥}{\ensuremath{\perp}}
\newunicodechar{∃}{\ensuremath{\exists}}
\newunicodechar{∀}{\ensuremath{\forall}}
\newunicodechar{∛}{\ensuremath{\sqrt[3]{\,}}}
\newunicodechar{∜}{\ensuremath{\sqrt[4]{\,}}}
\newunicodechar{⃗}{\ensuremath{{}^{\to}}}
\newunicodechar{ⁿ}{\ifmmode^{n}\else\textsuperscript{\ensuremath{n}}\fi}
\newunicodechar{ˣ}{\ifmmode^{x}\else\textsuperscript{\ensuremath{x}}\fi}
\newunicodechar{ᵇ}{\ifmmode^{b}\else\textsuperscript{\ensuremath{b}}\fi}
\newunicodechar{ʳ}{\ifmmode^{r}\else\textsuperscript{\ensuremath{r}}\fi}
\newunicodechar{ᵘ}{\ifmmode^{u}\else\textsuperscript{\ensuremath{u}}\fi}
\newunicodechar{ʸ}{\ifmmode^{y}\else\textsuperscript{\ensuremath{y}}\fi}
\newunicodechar{ᵃ}{\ifmmode^{a}\else\textsuperscript{\ensuremath{a}}\fi}
\newunicodechar{ᵉ}{\ifmmode^{e}\else\textsuperscript{\ensuremath{e}}\fi}
\newunicodechar{ᵏ}{\ifmmode^{k}\else\textsuperscript{\ensuremath{k}}\fi}
\newunicodechar{ᵖ}{\ifmmode^{p}\else\textsuperscript{\ensuremath{p}}\fi}
\newunicodechar{ᴺ}{\ifmmode^{N}\else\textsuperscript{\ensuremath{N}}\fi}
\newunicodechar{ᶜ}{\ifmmode^{c}\else\textsuperscript{\ensuremath{c}}\fi}
\newunicodechar{ʰ}{\ifmmode^{h}\else\textsuperscript{\ensuremath{h}}\fi}
\newunicodechar{ᵠ}{\ifmmode^{q}\else\textsuperscript{\ensuremath{q}}\fi}
\newunicodechar{ₙ}{\ifmmode_{n}\else\textsubscript{\ensuremath{n}}\fi}
\newunicodechar{ₐ}{\ifmmode_{a}\else\textsubscript{\ensuremath{a}}\fi}
\newunicodechar{ᵢ}{\ifmmode_{i}\else\textsubscript{\ensuremath{i}}\fi}
\newunicodechar{ᵣ}{\ifmmode_{r}\else\textsubscript{\ensuremath{r}}\fi}
\newunicodechar{ₖ}{\ifmmode_{k}\else\textsubscript{\ensuremath{k}}\fi}
\newunicodechar{ᵦ}{\ifmmode_{\beta}\else\textsubscript{\ensuremath{\beta}}\fi}
\newunicodechar{ₓ}{\ifmmode_{x}\else\textsubscript{\ensuremath{x}}\fi}
\newfontfamily\popdisplay{ArchivoBlack-Regular.ttf}[Path=../../../fonts/]
\newfontfamily\poplogo{SpaceGrotesk-Bold.ttf}[Path=../../../fonts/]
\newfontfamily\popsemi{PlusJakartaSans-SemiBold.ttf}[Path=../../../fonts/]
\newfontfamily\popxbold{PlusJakartaSans-ExtraBold.ttf}[Path=../../../fonts/]
\newfontfamily\popxboldls{PlusJakartaSans-ExtraBold.ttf}[Path=../../../fonts/,LetterSpace=3.0]

\setlist[enumerate]{leftmargin=1.7em,itemsep=5pt,topsep=4pt}
\setlist[itemize]{leftmargin=1.7em,itemsep=4pt,topsep=4pt,label={\color{poporange}\textbullet}}
\setlength{\parindent}{0pt}
\setlength{\parskip}{5pt}
\renewcommand{\arraystretch}{1.25}

% pgfplots : style Pop par défaut
\pgfplotsset{
  every axis/.append style={axis line style={popink,line width=1pt},tick style={popink},
    grid style={popline,line width=0.5pt},label style={font=\small},tick label style={font=\footnotesize}},
  popcurve/.style={poporange,line width=1.8pt,no markers,smooth},
}
% Même style disponible dans un \draw TikZ simple (hors pgfplots)
\tikzset{popcurve/.style={poporange,line width=1.8pt}}
\tikzset{popgrid/.style={popline,very thin}}
\colorlet{popcurve}{poporange} % le nom du style est parfois utilisé comme couleur

% ============================================================
% Pill / squiggle
% ============================================================
\newcommand{\poppill}[3]{%
  \tikz[baseline=-0.55ex]\node[draw=popink,line width=1.2pt,rounded corners=2.6pt,%
    fill=#2,inner xsep=6.5pt,inner ysep=3pt]{%
    \popxboldls\fontsize{7}{7}\selectfont\color{#3}\MakeUppercase{#1}};%
}
\newcommand{\popsquiggle}[1]{%
  \tikz[baseline=-0.4ex]{
    \draw[#1,line width=2.6pt,line cap=round]
      plot [smooth] coordinates {(0,0.09) (0.34,0.01) (0.68,0.21) (1.02,0.01) (1.36,0.10)};
  }%
}
% Mot-clé surligné (marqueur orange sous le texte)
\newcommand{\cle}[1]{{\popxbold\color{popdark}#1}}

% ============================================================
% En-tête / pied
% ============================================================
\pagestyle{fancy}
\fancyhf{}
\renewcommand{\headrulewidth}{0pt}
\renewcommand{\footrulewidth}{0pt}
\lhead{\raisebox{-0.15ex}{\poplogo\fontsize{13}{13}\selectfont\color{popink}OnBuch\color{poporange}+}}
\rhead{\poppill{\DOCMATIERE\ \textbullet\ \DOCNIVEAU\ \textbullet\ Leçon \DOCLECON}{white}{popink}}
\lfoot{\popxbold\fontsize{7.6}{7.6}\selectfont\color{popmuted}\MakeUppercase{Cours signé OnBuch+}\ \textbullet\ {\popsemi\DOCTITRE}}
\rfoot{\tikz[baseline=-0.6ex]\node[rounded corners=8.5pt,fill=popink,minimum height=17pt,minimum width=17pt,inner xsep=5pt,inner sep=0]{\popxbold\fontsize{8}{8}\selectfont\color{popcream}\thepage\,/\,\pageref*{LastPage}};}

% ============================================================
% Titres
% ============================================================
\newcommand{\chaptitre}[2]{%
  \begin{tcolorbox}[enhanced,colback=poporange,colframe=popink,boxrule=2.6pt,arc=11pt,
      drop shadow=popink,left=15pt,right=15pt,top=12pt,bottom=11pt,before skip=3pt,after skip=18pt]
    \poppill{#2}{popink}{popcream}\par\vspace{9pt}
    {\raggedright\hyphenpenalty=10000\exhyphenpenalty=10000\popdisplay\fontsize{22}{24}\selectfont\color{popcream}\MakeUppercase{#1}\par}
  \end{tcolorbox}
}
% Section numérotée : pastille orange avec le numéro + titre Archivo
\newcounter{popsec}
\newcommand{\coursec}[1]{%
  \stepcounter{popsec}\setcounter{popsub}{0}%
  \par\vspace{16pt}\noindent
  \begin{minipage}{\linewidth}
  \tikz[baseline=-0.7ex]\node[draw=popink,line width=1.6pt,fill=poporange,rounded corners=4pt,minimum size=22pt,inner sep=0]{\popdisplay\fontsize{12}{12}\selectfont\color{popcream}\thepopsec};%
  \hspace{8pt}{\popdisplay\fontsize{15}{17}\selectfont\color{popink}\MakeUppercase{#1}}\par
  \vspace{4pt}\noindent{\color{poporange}\rule{36pt}{3.6pt}}
  \end{minipage}\par\nopagebreak\vspace{8pt}
}
\newcounter{popsub}
\newcommand{\courssub}[1]{%
  \stepcounter{popsub}%
  \par\vspace{9pt}\noindent{\popxbold\fontsize{11.5}{13}\selectfont\color{popink}{\color{poporange}\thepopsec.\thepopsub}\hspace{6pt}#1}\par\nopagebreak\vspace{4pt}
}

% ============================================================
% Boîtes pédagogiques — même recette : fond blanc, bordure épaisse colorée,
% ombre dure encre, badge plein. Titre optionnel [#1].
% ============================================================
\tcbset{popbase/.style={breakable,enhanced,colback=white,boxrule=2.3pt,arc=9pt,
  shadow={3pt}{-3pt}{0pt}{popink},left=14pt,right=14pt,top=11pt,bottom=11pt,before skip=11pt,after skip=15pt,
  fontupper=\popsemi\fontsize{10.6}{14.5}\selectfont\color{popink},
  fontlower=\popsemi\fontsize{10.6}{14.5}\selectfont\color{popink}}}
% Coche dessinée (le glyphe ✓ n'existe pas dans les polices du document)
\newcommand{\popcheck}[1]{\tikz[baseline=-0.5ex]\draw[#1,line width=1.6pt,line cap=round,line join=round](-0.13,0.0)--(-0.04,-0.09)--(0.13,0.1);}
\providecommand{\coche}{\popcheck{popgreen}}
\providecommand{\resultat}[1]{\par\smallskip\noindent\fcolorbox{popgreen}{popgreenL}{\parbox{\dimexpr\linewidth-2\fboxsep-2\fboxrule\relax}{\textbf{Résultat.} #1}}\par\smallskip}
\newcommand{\popboxhead}[3]{\poppill{#1}{#2}{white}\ifx\relax#3\relax\else\hspace{6pt}{\popxbold\fontsize{10.6}{12}\selectfont\color{popink}#3}\fi\par\vspace{7pt}}

\newtcolorbox{aretenir}[1][]{popbase,colframe=popgreen,before upper={\popboxhead{À retenir}{popgreen}{#1}}}
% Texte en chinois (document de lecture, dialogue, extrait) : cadre
% d'étude. Un titre optionnel (source, auteur) s'affiche dans l'en-tête.
\newtcolorbox{textebox}[1][]{popbase,colframe=popblue,before upper={\popboxhead{文本 · Texte}{popblue}{#1}},after upper={}}
% Marques ✓ / ✗ (forme correcte / fautive) souvent utilisées par les modèles
\providecommand{\cross}{\textcolor{poppink}{\ensuremath{\times}}}
\providecommand{\xmark}{\cross}\providecommand{\crossmark}{\cross}\providecommand{\cmark}{\ensuremath{\checkmark}}
% Ligne de réponse / justification à compléter (inventée par les modèles)
\providecommand{\justification}[1]{\par\noindent\textit{Justification~:} \dotfill\par}
% \textipa (package tipa non chargé) : la police ne couvre pas l'API -> simple passage
\providecommand{\textipa}[1]{#1}
% Soulignement (ulem non chargé) : \uline, \ul
\providecommand{\uline}[1]{\underline{#1}}\providecommand{\ul}[1]{\underline{#1}}
% \glossaire{\item ...} inventée par les modèles : liste de glossaire
\providecommand{\glossaire}[1]{\begin{itemize}#1\end{itemize}}
% \alert (beamer) inventée par les modèles : texte mis en évidence
\providecommand{\alert}[1]{\textbf{\textcolor{poppink}{#1}}}
% variantes ASCII d'ordinaux écrites en macro
\providecommand{\iere}{\textsuperscript{re}}\providecommand{\ieme}{\textsuperscript{e}}\providecommand{\ere}{\textsuperscript{re}}\providecommand{\eme}{\textsuperscript{e}}\providecommand{\er}{\textsuperscript{er}}\providecommand{\ie}{\textsuperscript{e}}
% Ordinaux français écrits en macro par les modèles : 1\ière, 3\ième
\newcommand{\ière}{\textsuperscript{re}}\newcommand{\ième}{\textsuperscript{e}}
% \exemple (inventée par les modèles) : étiquette « Exemple : »
\providecommand{\exemple}{\textbf{Exemple~:}\ }
% \square utilisable aussi hors mode math (cases à cocher)
\AtBeginDocument{\let\squareMath\square\renewcommand{\square}{\ensuremath{\squareMath}}}
% Mot ou phrase en chinois à l'intérieur d'un paragraphe français
\newcommand{\zh}[1]{\ifmmode\text{#1}\else{#1}\fi}
\let\eng\zh \let\esp\zh \let\all\zh % alias tolérés
% flèches d'intonation inventées par les modèles
\providecommand{\textsearrow}{}\renewcommand{\textsearrow}{\ensuremath{\searrow}}\providecommand{\textnearrow}{}\renewcommand{\textnearrow}{\ensuremath{\nearrow}}
% \fr{..} : mot français (inventé par les modèles)
\providecommand{\fr}[1]{\textit{#1}}
% zhbox : encadré de texte chinois inventé par les modèles
\newenvironment{zhbox}{\begin{quote}}{\end{quote}}
% \courssubsub : titre de niveau 3 inventé par les modèles
\providecommand{\courssubsub}[1]{\par\medskip\noindent\textbf{#1}\par\smallskip}
% \vs : « versus » inventé par les modèles
\providecommand{\vs}{\textit{vs}\ }
% \blank : case à compléter inventée par les modèles
\providecommand{\blank}[1][]{\underline{\hspace{1.6em}}}
\newtcolorbox{exemplebox}[1][]{popbase,colframe=poporange,before upper={\popboxhead{Exemple}{poporange}{#1}}}
\newtcolorbox{definition}[1][]{popbase,colframe=popblue,before upper={\popboxhead{Définition}{popblue}{#1}}}
\newtcolorbox{propriete}[1][]{popbase,colframe=poppurple,before upper={\popboxhead{Propriété}{poppurple}{#1}}}
\newtcolorbox{methode}[1][]{popbase,colframe=popgold,before upper={\popboxhead{Méthode}{popgold}{#1}}}
\newtcolorbox{attention}[1][]{popbase,colframe=poppink,before upper={\popboxhead{Attention}{poppink}{#1}}}
\newtcolorbox{experience}[1][]{popbase,colframe=popblue,colback=popblueL,before upper={\popboxhead{Expérience}{popblue}{#1}}}
% « Le savais-tu ? » : carte violette pleine (contexte, culture, Cameroun)
\newtcolorbox{savaistu}[1][]{popbase,colback=poppurple,colframe=popink,
  fontupper=\popsemi\fontsize{10.4}{14.2}\selectfont\color{white},
  before upper={\poppill{Le savais-tu\,?}{white}{poppurple}\ifx\relax#1\relax\else\hspace{6pt}{\popxbold\color{white}#1}\fi\par\vspace{7pt}}}
% Exercice résolu : énoncé en haut, \tcblower puis la solution sur fond crème
\newcounter{popexo}\newcounter{popexr}
\newtcolorbox{exoresolu}[1][]{popbase,colframe=poporange,colbacklower=popcream,
  segmentation style={popink,line width=1.2pt,dashed},
  before upper={\stepcounter{popexr}\popboxhead{Exercice résolu \thepopexr}{poporange}{#1}},
  before lower={\poppill{Solution}{popink}{popcream}\par\vspace{6pt}}}
% Exercice (non résolu en place) — la correction est dans la partie Corrigés
\newtcolorbox{exercice}[1][]{popbase,colframe=popink,
  before upper={\stepcounter{popexo}\popboxhead{Exercice \thepopexo}{popink}{#1}}}
% Corrigé
\newtcolorbox{corrige}[1][]{popbase,colframe=popgreen,colback=popgreenL,
  before upper={\popboxhead{Corrigé}{popgreen}{#1}}}
% Objectifs / prérequis / bilan
\newtcolorbox{objectifs}{popbase,colframe=popink,colback=poporangeL,
  before upper={\popboxhead{Objectifs de la leçon}{popink}{}}}
\newtcolorbox{prerequis}{popbase,colframe=popink,colback=popblueL,
  before upper={\popboxhead{Prérequis}{popblue}{}}}
\newtcolorbox{bilan}{popbase,colframe=popink,colback=white,boxrule=2.8pt,
  before upper={\popboxhead{Fiche bilan}{poporange}{L'essentiel à connaître pour l'examen}}}
% Figure encadrée avec légende
\newtcolorbox{popfigure}[1][]{enhanced,breakable=false,colback=white,colframe=popline,boxrule=1.4pt,arc=8pt,
  left=10pt,right=10pt,top=10pt,bottom=8pt,before skip=10pt,after skip=12pt,halign=center,
  lower separated=false,halign lower=center,
  fontlower=\popsemi\fontsize{9}{11.5}\selectfont\color{popmuted},
  #1}
\newcommand{\legende}[1]{\par\vspace{6pt}{\popsemi\fontsize{9}{11.5}\selectfont\color{popmuted}#1}}

% ============================================================
% Signature OnBuch+
% ============================================================
\newcommand{\poplogoinline}[1]{{\poplogo\fontsize{#1}{#1}\selectfont\color{popink}OnBuch\color{poporange}+}}
\newcommand{\signatureonbuch}{%
  \begin{tcolorbox}[enhanced,colback=popink,colframe=popink,boxrule=2.2pt,arc=10pt,
      drop shadow=poporange,left=14pt,right=14pt,top=10pt,bottom=10pt,before skip=0pt,after skip=0pt]
    \tikz[baseline=-0.7ex]\node[circle,fill=popgreen,draw=popcream,line width=1.3pt,minimum size=22pt,inner sep=0]{\popcheck{white}};%
    \hspace{9pt}\begin{minipage}[c]{0.62\linewidth}
      {\popxbold\fontsize{12}{13}\selectfont\color{popcream}Signé\ \textbullet\ {\poplogo OnBuch\color{poporange}+}}\par\vspace{2pt}
      {\raggedright\popsemi\fontsize{8}{10.5}\selectfont\color{popline}\MakeUppercase{Équipe pédagogique OnBuch+}\\\MakeUppercase{Conforme au programme officiel MINESEC}\par}
    \end{minipage}\hfill
    {\poplogo\fontsize{22}{22}\selectfont\color{popcream}OnBuch\color{poporange}+}
  \end{tcolorbox}%
}

% ============================================================
% Page de garde
% #1 titre · #2 matière · #3 niveau · #4 module · #5 n° leçon · #6 durée ·
% #7 description / objectifs (texte ou itemize)
% ============================================================
\newcommand{\pagegarde}[7]{%
  \thispagestyle{empty}
  \newgeometry{a4paper,margin=1.7cm,top=1.6cm,bottom=1.4cm}%
  \begin{tikzpicture}[remember picture,overlay]
    \fill[poporange!18] ([xshift=-3.2cm,yshift=-3.6cm]current page.north east) circle (4.6cm);
    \fill[poppurple!14] ([xshift=2.4cm,yshift=7.6cm]current page.south west) circle (3.8cm);
  \end{tikzpicture}%
  {\poplogo\fontsize{30}{30}\selectfont\color{popink}OnBuch\color{poporange}+}\hfill
  \raisebox{0.6ex}{\poppill{Cours signé \textbullet\ Premium}{popink}{popcream}}\par
  \vspace{1pt}\popsquiggle{poppurple}\par
  \vspace{8pt}
  \begin{tcolorbox}[enhanced,colback=poporange,colframe=popink,boxrule=2.8pt,arc=13pt,
      drop shadow=popink,left=18pt,right=18pt,top=12pt,bottom=13pt,after skip=10pt]
    \poppill{#2 \textbullet\ #3}{popink}{popcream}\hspace{5pt}\poppill{#4}{white}{popink}\par\vspace{8pt}
    {\popdisplay\fontsize{14}{14}\selectfont\color{popink}LEÇON #5}\par\vspace{4pt}
    {\raggedright\hyphenpenalty=10000\exhyphenpenalty=10000\popdisplay\fontsize{25}{27}\selectfont\color{popcream}\MakeUppercase{#1}\par}
    \vspace{8pt}
    {\popxbold\fontsize{9.5}{9.5}\selectfont\color{popink}\MakeUppercase{Durée conseillée : #6}}
  \end{tcolorbox}
  \begin{tcolorbox}[enhanced,colback=white,colframe=popink,boxrule=2.2pt,arc=10pt,
      drop shadow=popink,left=16pt,right=16pt,top=10pt,bottom=10pt,after skip=8pt]
    \poppill{Dans cette leçon}{poppurple}{white}\par\vspace{5pt}
    \popsemi\fontsize{9.3}{12.6}\selectfont\color{popink}#7
  \end{tcolorbox}
  \noindent\begin{minipage}[t]{0.487\linewidth}
    \begin{tcolorbox}[enhanced,colback=popgreen,colframe=popink,boxrule=2.2pt,arc=10pt,
        drop shadow=popink,left=11pt,right=11pt,top=10pt,bottom=10pt,height=3.1cm,valign=center]
      \tcbox[colback=white,colframe=popink,boxrule=1.3pt,arc=5pt,boxsep=0pt,nobeforeafter,
        left=3pt,right=3pt,top=3pt,bottom=3pt,valign=center]{\qrcode[height=1.9cm]{https://play.google.com/store/apps/details?id=com.luvvix.onbuch&hl=fr}}%
      \hspace{8pt}\begin{minipage}[c]{0.5\linewidth}
        {\popdisplay\fontsize{11}{12.5}\selectfont\color{white}\MakeUppercase{Télécharge OnBuch}}\par\vspace{4pt}
        {\raggedright\popsemi\fontsize{8.4}{11}\selectfont\color{white}Gratuit \textbullet\ Google Play\\Tuteur IA, annales, concours, résultats\par}
      \end{minipage}
    \end{tcolorbox}
  \end{minipage}\hfill
  \begin{minipage}[t]{0.487\linewidth}
    \begin{tcolorbox}[enhanced,colback=poppurple,colframe=popink,boxrule=2.2pt,arc=10pt,
        drop shadow=popink,left=11pt,right=11pt,top=10pt,bottom=10pt,height=3.1cm,valign=center]
      \tikz[baseline=-0.7ex]\node[circle,fill=white,draw=popink,line width=1.3pt,minimum size=1.6cm,inner sep=0]{\popdisplay\fontsize{24}{24}\selectfont\color{poppurple}+};%
      \hspace{8pt}\begin{minipage}[c]{0.6\linewidth}
        {\popdisplay\fontsize{11}{12.5}\selectfont\color{white}\MakeUppercase{Passe à OnBuch+}}\par\vspace{4pt}
        {\raggedright\popsemi\fontsize{8.4}{11}\selectfont\color{white}Tous les cours signés, Tuteur IA illimité, fiches PDF — onglet OnBuch+ de l'app\par}
      \end{minipage}
    \end{tcolorbox}
  \end{minipage}
  \vspace{8pt}
  \popcontenu
  \vfill
  \signatureonbuch
  \restoregeometry
  \newpage
}

% Rangée « Ce cours contient » (page de garde)
\newcommand{\poptile}[3]{% #1 couleur #2 gros texte #3 légende
  \begin{tcolorbox}[enhanced,colback=#1,colframe=popink,boxrule=2pt,arc=9pt,shadow={2.5pt}{-2.5pt}{0pt}{popink},
      left=6pt,right=6pt,top=9pt,bottom=9pt,halign=center,valign=center,height=1.75cm,width=\linewidth,nobeforeafter]
    {\popdisplay\fontsize{12.5}{13}\selectfont\color{white}\MakeUppercase{#2}}\par\vspace{3pt}
    {\centering\popsemi\fontsize{7.8}{9.5}\selectfont\color{white}#3\par}
  \end{tcolorbox}}
\newcommand{\popcontenu}{%
  \par\noindent{\popxboldls\fontsize{7.5}{7.5}\selectfont\color{popmuted}CE COURS SIGNÉ CONTIENT}\par\vspace{7pt}
  \noindent\begin{minipage}[t]{0.235\linewidth}\poptile{popblue}{Cours}{complet, illustré, pas à pas}\end{minipage}\hfill
  \begin{minipage}[t]{0.235\linewidth}\poptile{poporange}{Méthodes}{et exercices résolus}\end{minipage}\hfill
  \begin{minipage}[t]{0.235\linewidth}\poptile{poppink}{Exercices}{type examen + corrigés}\end{minipage}\hfill
  \begin{minipage}[t]{0.235\linewidth}\poptile{popgreen}{Fiche bilan}{l'essentiel à retenir}\end{minipage}\par}

% Sommaire
\newtcolorbox{sommaire}{enhanced,colback=white,colframe=popink,boxrule=2.2pt,arc=10pt,
  drop shadow=popink,left=18pt,right=18pt,top=13pt,bottom=13pt,before skip=6pt,after skip=20pt,
  before upper={\poppill{Sommaire}{poppurple}{white}\par\vspace{9pt}\popsemi\fontsize{10.6}{14.5}\selectfont\color{popink}}}

% Signature de fin de cours — poussée tout en bas de la dernière page
\newcommand{\findecours}{%
  \par\vfill\nopagebreak
  \begin{center}\popsquiggle{poporange}\end{center}\vspace{4pt}
  \signatureonbuch
}
\AtBeginDocument{\def\llbracket{\mathopen{[\mkern-3.5mu[}}\def\rrbracket{\mathclose{]\mkern-3.5mu]}}} % intervalles d'entiers (glyphe ⟦ absent de la police)
% Glyphes absents de Fira Math : redéfinis à partir de glyphes disponibles
\AtBeginDocument{%
  \def\setminus{\mathbin{\backslash}}\let\smallsetminus\setminus
  \def\vdots{\mathord{\vbox{\baselineskip3.2pt\lineskiplimit0pt\kern4pt\hbox{.}\hbox{.}\hbox{.}}}}%
  \def\ddots{\mathinner{\mkern1mu\raise6.5pt\hbox{.}\mkern2mu\raise3.6pt\hbox{.}\mkern2mu\raise0.7pt\hbox{.}\mkern1mu}}%
  \def\triangle{\mathord{\scalebox{0.9}{$\symup{\Delta}$}}}%
  \def\nabla{\mathord{\raisebox{\depth}{\scalebox{1}[-1]{$\symup{\Delta}$}}}}%
  \def\checkmark{\popcheck{popdark}}%
  \def\star{\mathbin{\ast}}\def\bigstar{\mathord{\ast}}%
  \def\diamond{\mathbin{\scalebox{0.6}{$\Diamond$}}}%
  \def\bigcup{\mathop{\mathchoice{\raisebox{-0.3em}{\scalebox{1.7}{$\cup$}}}{\scalebox{1.25}{$\cup$}}{\cup}{\cup}}\displaylimits}%
  \def\bigcap{\mathop{\mathchoice{\raisebox{-0.3em}{\scalebox{1.7}{$\cap$}}}{\scalebox{1.25}{$\cap$}}{\cap}{\cap}}\displaylimits}%
  \def\lesseqqgtr{\mathrel{\lessgtr}}\def\bigoplus{\mathop{\oplus}\displaylimits}%
}
\providecommand{\ssi}{si et seulement si} % abréviation fréquente
\AtBeginDocument{\providecommand{\wideparen}{\overparen}} % arc de cercle

%% FIGFIX-BEGIN v3 — correctifs globaux des figures (docs/onbuch-generation/tools/figfix)
\usepackage{adjustbox}\usepackage{etoolbox}
% toute figure TikZ/pgfplots plus large que la ligne est réduite à la largeur disponible
\BeforeBeginEnvironment{tikzpicture}{\begin{adjustbox}{max width=\linewidth}}
\AfterEndEnvironment{tikzpicture}{\end{adjustbox}}
% légendes pgfplots lisibles par-dessus les courbes
\pgfplotsset{every axis/.append style={legend style={fill=white,fill opacity=0.92,text opacity=1,draw=black!15,inner sep=3pt}}}
% étiquettes d'arêtes (syntaxe edge["..."]) avec fond blanc
\tikzset{every edge quotes/.append style={fill=white,inner sep=1.2pt}}
% bulles de cartes mentales : texte à la ligne dans la bulle, bulle un peu plus grande
\tikzset{every concept/.append style={text width=2.0cm,align=center,minimum size=2.3cm,inner sep=2pt}}
%% FIGFIX-END
\begin{document}

\pagegarde{Leçon 26 — Grammaire : L’addition et soustraction en chinois 加法和减法 ; les différentes monnaies ; Les nombres ; 除了……以外，……}{Chinois}{Tle A}{Module 4 : Activités économiques et monde du travail}{ZH26}{2 h}{Cette leçon de grammaire enseigne comment exprimer l'addition et la soustraction en chinois (加法和减法), nommer les différentes monnaies du monde et manipuler les nombres dans des calculs simples. Elle introduit également la structure 除了……以外，…… pour exprimer l'addition ou l'exclusion, avec comparaison systématique au français.
\vspace{4pt}
\begin{multicols}{2}\small
\begin{itemize}
\item À la fin de cette leçon, je sais former une addition en chinois avec 加 et 等于 (A + B = C)
\item À la fin de cette leçon, je sais former une soustraction en chinois avec 减 et 等于 (A - B = C)
\item À la fin de cette leçon, je sais nommer les principales monnaies : 人民币/块, 美元, 欧元, 法郎, 英镑, 奈拉, 塞地
\item À la fin de cette leçon, je sais utiliser correctement les nombres (unités, dizaines, centaines, milliers, 万) dans des calculs et des prix
\item À la fin de cette leçon, je sais construire la structure 除了……以外，…… avec 也 (exclusion) et 还/都 (addition)
\item À la fin de cette leçon, je sais distinguer les deux emplois de 除了……以外 (sens inclusif et exclusif) selon le mot-outil qui suit
\end{itemize}\end{multicols}}

\begin{sommaire}
\begin{multicols}{2}
\begin{enumerate}
\item Observation : calculer au marché
\item Structure 1 : l'addition 加法
\item Structure 2 : la soustraction 减法
\item Les monnaies et les nombres
\item Structure 3 : 除了……以外，……
\item Comparaison avec le français et synthèse
\item Activité d'intégration
\item Méthodes et exercices résolus
\item Exercices
\item Corrigés des exercices
\item Fiche bilan
\end{enumerate}
\end{multicols}
\end{sommaire}

\begin{objectifs}
\begin{itemize}
\item À la fin de cette leçon, je sais former une addition en chinois avec 加 et 等于 (A + B = C)
\item À la fin de cette leçon, je sais former une soustraction en chinois avec 减 et 等于 (A - B = C)
\item À la fin de cette leçon, je sais nommer les principales monnaies : 人民币/块, 美元, 欧元, 法郎, 英镑, 奈拉, 塞地
\item À la fin de cette leçon, je sais utiliser correctement les nombres (unités, dizaines, centaines, milliers, 万) dans des calculs et des prix
\item À la fin de cette leçon, je sais construire la structure 除了……以外，…… avec 也 (exclusion) et 还/都 (addition)
\item À la fin de cette leçon, je sais distinguer les deux emplois de 除了……以外 (sens inclusif et exclusif) selon le mot-outil qui suit
\item À la fin de cette leçon, je sais éviter les erreurs fréquentes des francophones (ordre des mots, oubli de 等于, confusion 二/两, calque de « sauf »)
\item À la fin de cette leçon, je sais rédiger un court dialogue d'achat avec calcul de prix et monnaie rendue
\end{itemize}
\end{objectifs}

\begin{prerequis}
\begin{itemize}
\item Les nombres de 0 à 10 et les classificateurs de base (个, 块, 本) vus en Première
\item La formation des nombres jusqu'à 9999 (十, 百, 千) et l'introduction de 万
\item Les verbes et adverbes déjà étudiés : 是, 有, 也, 还, 都
\item Le vocabulaire de l'argent et des achats : 钱, 买, 卖, 多少, 贵, 便宜
\item La structure interrogative 多少钱 ? vue en Première
\end{itemize}
\end{prerequis}

\newpage

\coursec{Observation : calculer au marché}

\courssub{Situation d'accroche : la monnaie d'Amina}

C'est samedi matin au marché Mokolo, à Yaoundé. Amina achète 3 ananas à 500 FCFA l'unité, puis tend un billet de 2000 FCFA au vendeur. Le calcul est immédiat : 3 fois 500 font 1500, et 2000 moins 1500 font 500. Il lui reste 500 FCFA. Addition, soustraction : voilà des opérations que nous faisons chaque jour, au marché, à la boutique, au comptoir de la cantine.

En Chine, le même calcul se fait en yuans (\zh{人民币} rénmínbì), avec des structures bien précises : \zh{三加二等于五} (3 + 2 = 5), \zh{十减四等于六} (10 − 4 = 6). Et quand on voyage, on découvre qu'« à part le yuan, on peut aussi payer par téléphone » : \zh{除了人民币以外，还可以用手机支付。} Comment dire tout cela en chinois ? C'est l'objet de cette leçon.

\textbf{Problématique.} Comment exprimer en chinois une addition, une soustraction, et l'idée de « à part…, en plus de… » ? Où se placent les mots-outils par rapport aux nombres ? Observons d'abord.

\courssub{Lecture guidée : quatre opérations à voix haute}

Lisons ensemble les quatre phrases suivantes. Respecte bien les tons : \zh{加} jiā se dit avec le premier ton (haut et plat), \zh{减} jiǎn avec le troisième ton (descendant puis montant).

\begin{textebox}[Au tableau de calcul]
三加五等于八。\\
\emph{Sān jiā wǔ děngyú bā.}\\
« Trois plus cinq égale huit. » (3 + 5 = 8)\\[4pt]
十减四等于六。\\
\emph{Shí jiǎn sì děngyú liù.}\\
« Dix moins quatre égale six. » (10 − 4 = 6)\\[4pt]
一百加二百等于三百。\\
\emph{Yìbǎi jiā èrbǎi děngyú sānbǎi.}\\
« Cent plus deux cents égale trois cents. » (100 + 200 = 300)\\[4pt]
一千减五百等于五百。\\
\emph{Yìqiān jiǎn wǔbǎi děngyú wǔbǎi.}\\
« Mille moins cinq cents égale cinq cents. » (1000 − 500 = 500)
\end{textebox}

\begin{attention}[Prononciation : 加 et 减]
Ne confonds pas \zh{加} (jiā, plus, 1\ier{} ton) et \zh{减} (jiǎn, moins, 3\ieme{} ton) : même initiale \emph{j}, même finale \emph{ian}, mais des tons différents. C'est l'erreur orale la plus fréquente des francophones : à l'écrit, les signes + et − nous dispensent de prononcer ; en chinois, tout passe par la voix, donc par le ton. Entraîne-toi en alternant : jiā — jiǎn — jiā — jiǎn.
\end{attention}

\courssub{Repérage des mots-outils}

Dans les quatre phrases, trois mots reviennent à chaque fois. Ce sont les mots-outils du calcul :

\begin{center}
\begin{tabularx}{\linewidth}{|X|X|X|X|}
\hline
\rowcolor{popblueL} Caractères & Pinyin & Nature & Traduction \\
\hline
加 & jiā & verbe & ajouter ; plus \\
\hline
减 & jiǎn & verbe & soustraire ; moins \\
\hline
等于 & děngyú & verbe & égaler ; être égal à \\
\hline
加法 & jiāfǎ & nom & l'addition \\
\hline
减法 & jiǎnfǎ & nom & la soustraction \\
\hline
计算 & jìsuàn & verbe / nom & calculer ; le calcul \\
\hline
\end{tabularx}
\end{center}

Remarque la formation des mots : \zh{加} + \zh{法} (méthode) donne \zh{加法}, « la méthode d'ajouter », c'est-à-dire l'addition ; de même \zh{减法}, la soustraction. Le chinois construit son vocabulaire par composition, comme des briques.

\courssub{Question d'observation : où se placent les mots-outils ?}

Reprenons la première phrase et découpons-la :

\begin{center}
\begin{tabularx}{\linewidth}{|X|X|X|X|X|}
\hline
\rowcolor{popblueL} 三 & 加 & 五 & 等于 & 八 \\
\hline
sān (3) & jiā (plus) & wǔ (5) & děngyú (égale) & bā (8) \\
\hline
\end{tabularx}
\end{center}

\textbf{Question d'observation.} Où se placent \zh{加}/\zh{减} et \zh{等于} par rapport aux nombres ?

Réponse attendue des élèves, à formuler collectivement : le mot-outil \zh{加} ou \zh{减} se place \cle{entre les deux nombres} de l'opération, et \zh{等于} se place \cle{entre l'opération et le résultat}. L'ordre est exactement celui du français « trois plus cinq égale huit » : nombre 1 + opération + nombre 2 + égale + résultat.

\begin{popfigure}
\begin{tikzpicture}[node distance=1.6cm and 2.2cm, every node/.style={font=\small}]
\node[draw=popblue, fill=popblueL, rounded corners, minimum width=1.4cm, minimum height=0.9cm] (A) {A : 三 sān};
\node[draw=poporange, fill=poporangeL, rounded corners, minimum width=1.4cm, minimum height=0.9cm, right=of A] (B) {B : 五 wǔ};
\node[draw=popgreen, fill=popgreenL, rounded corners, minimum width=1.4cm, minimum height=0.9cm, right=of B] (C) {C : 八 bā};
\draw[-{Stealth}, thick, popdark] (A) -- node[above, align=center]{加 jiā\\ \emph{plus / moins}} (B);
\draw[-{Stealth}, thick, popdark] (B) -- node[above, align=center]{等于 děngyú\\ \emph{égale}} (C);
\end{tikzpicture}
\legende{Schéma de la structure : A + 加/减 + B + 等于 + C}
\end{popfigure}

\begin{popfigure}
\begin{tikzpicture}[mindmap, grow cyclic, every node/.style={concept}, root concept/.append style={concept color=popblue, font=\small}, level 1/.append style={level distance=2.6cm, sibling angle=120, concept color=poporangeL}]
\node[concept, align=center]{计算 jìsuàn\\ \emph{calculer}}
  child { node[concept, align=center]{加法 jiāfǎ\\ \emph{addition}} }
  child { node[concept, align=center]{减法 jiǎnfǎ\\ \emph{soustraction}} }
  child { node[concept, align=center]{等于 děngyú\\ \emph{égale}} };
\end{tikzpicture}
\legende{Carte mentale : le vocabulaire du calcul en chinois}
\end{popfigure}

\courssub{Première formulation collective de la règle}

\begin{propriete}[Règle dégagée par la classe — le calcul en chinois]
Structure : \cle{Nombre 1 + 加/减 + Nombre 2 + 等于 + Résultat}.
\begin{itemize}
\item \zh{加} (jiā) exprime l'addition, \zh{减} (jiǎn) la soustraction ; ils se placent entre les deux nombres.
\item \zh{等于} (děngyú) annonce le résultat ; il ne se déplace jamais et ne change jamais de forme.
\item L'ordre des mots est identique à celui du français : « 3 + 5 = 8 » se lit de gauche à droite \zh{三加五等于八}.
\end{itemize}
\end{propriete}

\begin{exemplebox}[Deux exemples supplémentaires]
\zh{二加三等于五。} \emph{Èr jiā sān děngyú wǔ.} « Deux plus trois égale cinq. » (2 + 3 = 5)\\[4pt]
\zh{九减七等于二。} \emph{Jiǔ jiǎn qī děngyú èr.} « Neuf moins sept égale deux. » (9 − 7 = 2)
\end{exemplebox}

\begin{exoresolu}[Traduire en chinois]
Énoncé : traduis en chinois « Sept plus deux égale neuf » et « Vingt moins huit égale douze ».
\tcblower
Solution détaillée. On applique la structure Nombre 1 + 加/减 + Nombre 2 + 等于 + Résultat.\\
1) « Sept plus deux égale neuf » : sept = \zh{七} qī ; plus = \zh{加} jiā ; deux = \zh{二} èr ; égale = \zh{等于} děngyú ; neuf = \zh{九} jiǔ. On obtient : \zh{七加二等于九。} \emph{Qī jiā èr děngyú jiǔ.}\\
2) « Vingt moins huit égale douze » : vingt = \zh{二十} èrshí ; moins = \zh{减} jiǎn ; huit = \zh{八} bā ; égale = \zh{等于} děngyú ; douze = \zh{十二} shí'èr. On obtient : \zh{二十减八等于十二。} \emph{Èrshí jiǎn bā děngyú shí'èr.}\\
Vérification : chaque mot-outil est bien entre deux nombres, et la phrase se termine par le point chinois \zh{。}.
\end{exoresolu}

\begin{attention}[Pièges fréquents des francophones]
\begin{itemize}
\item Ne traduis pas « égale » par \zh{是} (shì, être) : on dit \zh{三加五等于八}, jamais \zh{三加五是八}. En calcul, le verbe obligatoire est \zh{等于}.
\item N'ajoute pas de classificateur ni de \zh{个} après les nombres dans une opération : \zh{三加五}, et non \zh{三个加五个}. Les nombres y sont nus.
\item Attention au ton de \zh{减} jiǎn (3\ieme{} ton) : prononcé jiā, il inverse toute l'opération !
\end{itemize}
\end{attention}

\begin{savaistu}[Le boulier chinois]
En Chine, les enfants apprennent traditionnellement les tables d'addition avec le boulier, \zh{算盘} (suànpán), instrument de calcul à boules utilisé depuis des siècles. Il est encore enseigné dans certaines écoles primaires chinoises, et les concours de calcul mental au boulier restent très populaires. Au Cameroun, c'est souvent au marché, comme Amina à Mokolo, que l'on devient champion du calcul mental !
\end{savaistu}

\coursec{Structure 1 : l'addition 加法}

\courssub{Transition depuis l'observation}

Dans la section précédente, nous avons observé un dialogue au marché de Mfoundi à Yaoundé où une acheteuse additionnait le prix du ndolé (\zh{五块} wǔ kuài, 5 yuans) et celui du riz (\zh{三块} sān kuài, 3 yuans) pour obtenir le total. En français, nous dirions « 5 plus 3 font 8 ». En chinois, la structure est tout aussi logique et transparente : on utilise le verbe \zh{加} (jiā, « ajouter / plus ») et le verbe \zh{等于} (děngyú, « être égal à »). Cette section pose les bases rigoureuses de l'addition, indispensables pour manipuler les prix, les quantités et les données chiffrées dans tout contexte économique.

\courssub{Observation et découverte}

\begin{textebox}[Dialogue au marché : calculer le total]
\textbf{Contexte :} Un élève de Terminale (小李 Xiǎo Lǐ) achète des fournitures dans une papeterie de Douala. Il demande le prix total.
\\ \zh{小李：老板，这本笔记本五块，那支笔三块，一共多少钱？}
\\ \zh{老板：五加三等于八，一共八块钱。}
\\ \zh{小李：好的，给您八块。}
\\
\textbf{Pinyin :}
\\ Xiǎo Lǐ: Lǎobǎn, zhè běn bǐjìběn wǔ kuài, nà zhī bǐ sān kuài, yīgòng duōshao qián?
\\ Lǎobǎn: Wǔ jiā sān děngyú bā, yīgòng bā kuài qián.
\\ Xiǎo Lǐ: Hǎo de, gěi nín bā kuài.
\\
\textbf{Traduction :}
\\ Petit Li : Patron, ce cahier coûte 5 yuans, ce stylo 3 yuans, combien au total ?
\\ Le patron : 5 plus 3 font 8, total 8 yuans.
\\ Petit Li : D'accord, voilà 8 yuans.
\end{textebox}

\noindent
Remarquez la phrase du patron : \zh{五加三等于八} (Wǔ jiā sān děngyú bā). Elle suit exactement l'ordre logique de l'opération : Nombre 1 + \zh{加} + Nombre 2 + \zh{等于} + Résultat.

\courssub{Schéma et règle fondamentale}

\begin{propriete}[Structure canonique de l'addition]
La structure de base de l'addition en chinois est immuable :
\\ \textbf{Nombre 1 + 加 (jiā) + Nombre 2 + 等于 (děngyú) + Nombre 3}
\\
\begin{itemize}
    \item \zh{加} (jiā) : verbe « ajouter », équivalent au « plus » mathématique (+).
    \item \zh{等于} (děngyú) : verbe « être égal à », équivalent au « = » ou au « font » du français.
    \item L'ordre des termes est \textbf{identique au français} : pas d'inversion, pas de préposition.
\end{itemize}
\end{propriete}

\begin{popfigure}
\begin{tikzpicture}[
    node distance=0.8cm and 1.2cm,
    box/.style={draw=popblue, thick, rounded corners, fill=popblueL, minimum height=1.2cm, minimum width=2.2cm, align=center, font=\Large},
    arrow/.style={-Stealth, thick, popdark},
    labelnode/.style={font=\small\itshape, text=popmuted, anchor=north}
]
% Nœuds principaux
\node[box, align=center] (n1){名词1\\ (N1)};
\node[box, right=of n1, align=center] (jia){\zh{加}\\ jiā};
\node[box, right=of jia, align=center] (n2){名词2\\ (N2)};
\node[box, right=of n2, align=center] (dengyu){\zh{等于}\\ děngyú};
\node[box, right=of dengyu, align=center] (n3){结果\\ (N3)};

% Exemple rempli en dessous
\node[box, below=1.5cm of n1, fill=popgreenL, draw=popgreen, align=center] (ex1){\zh{三}\\ sān};
\node[box, right=of ex1, fill=popgreenL, draw=popgreen, align=center] (exjia){\zh{加}\\ jiā};
\node[box, right=of exjia, fill=popgreenL, draw=popgreen, align=center] (ex2){\zh{五}\\ wǔ};
\node[box, right=of ex2, fill=popgreenL, draw=popgreen, align=center] (exdeng){\zh{等于}\\ děngyú};
\node[box, right=of exdeng, fill=popgreenL, draw=popgreen, align=center] (ex3){\zh{八}\\ bā};

% Flèches
\draw[arrow] (n1) -- (jia);
\draw[arrow] (jia) -- (n2);
\draw[arrow] (n2) -- (dengyu);
\draw[arrow] (dengyu) -- (n3);

% Étiquettes
\node[labelnode] at (n1.south) {Premier terme};
\node[labelnode] at (jia.south) {Opérateur};
\node[labelnode] at (n2.south) {Second terme};
\node[labelnode] at (dengyu.south) {Lien};
\node[labelnode] at (n3.south) {Résultat};

% Accolade exemple
\node[below=0.3cm of ex3, font=\small, text=popdark] {Exemple : 3 + 5 = 8};
\end{tikzpicture}
\legende{Schéma de la structure additive : N1 加 N2 等于 N3}
\end{popfigure}

\courssub{Emplois, variantes de registre et forme interrogative}

\begin{exemplebox}[Addition standard (registre neutre / écrit)]
\begin{itemize}
    \item \zh{二十加三十等于五十。} (Èrshí jiā sānshí děngyú wǔshí.) — 20 + 30 = 50.
    \item \zh{一百加两百等于三百。} (Yībǎi jiā liǎngbǎi děngyú sānbǎi.) — 100 + 200 = 300.
    \item \zh{零点五加零点五等于一。} (Líng diǎn wǔ jiā líng diǎn wǔ děngyú yī.) — 0,5 + 0,5 = 1.
\end{itemize}
\end{exemplebox}

\begin{exemplebox}[Variante familière orale : \zh{是} (shì) au lieu de \zh{等于}]
À l'oral, dans la vie quotidienne (marché, cours de récréation, calcul mental rapide), \zh{等于} est souvent remplacé par le verbe \zh{是} (shì, « être »).
\\ \zh{三加五是八。} (Sān jiā wǔ shì bā.) — 3 + 5 font 8.
\\ \zh{六加四是十。} (Liù jiā sì shì shí.) — 6 + 4 font 10.
\\
\cle{Repère} : Réservez \zh{等于} pour l'écrit, les examens, les explications formelles. Utilisez \zh{是} à l'oral uniquement.
\end{exemplebox}

\begin{exemplebox}[Forme interrogative : \zh{等于几} / \zh{等于多少}]
Pour demander le résultat d'une addition, on remplace le nombre final par un mot interrogatif.
\begin{itemize}
    \item \zh{几} (jǐ) : « combien » (résultat attendu \textbf{inférieur à 10}).
    \item \zh{多少} (duōshao) : « combien » (résultat \textbf{supérieur ou égal à 10}, ou inconnu grand).
\end{itemize}
\textbf{Exemples :}
\\ \zh{二加三等于几？} (Èr jiā sān děngyú jǐ?) — 2 + 3 = ? (Réponse attendue : 5)
\\ \zh{二十加三十等于多少？} (Èrshí jiā sānshí děngyú duōshao?) — 20 + 30 = ? (Réponse : 50)
\\ \zh{一百加两百等于多少？} (Yībǎi jiā liǎngbǎi děngyú duōshao?) — 100 + 200 = ?
\end{exemplebox}

\begin{exemplebox}[Réponse courte]
On peut omettre la répétition de l'opération.
\\ \zh{— 五加五等于几？} — \zh{等于十。} / \zh{十。}
\\ \zh{— 八加七等于多少？} — \zh{等于十五。} / \zh{十五。}
\end{exemplebox}

\courssub{Le cas des sommes d'argent : l'exception des classificateurs}

\begin{attention}[Piège majeur : Classificateur \textbf{interdit} dans le calcul abstrait, \textbf{obligatoire} pour l'argent]
C'est l'erreur n°1 des francophones qui traduisent mot à mot « 5 francs + 3 francs ».
\begin{itemize}
    \item \textbf{Calcul abstrait (mathématiques pures) :} \textbf{Pas de classificateur}.
    \\ \zh{五加三等于八。} (Correct) \hspace{2cm} \zh{五块加三块等于八块。} (Incorrect pour un calcul pur)
    \item \textbf{Sommes d'argent (prix, monnaie) :} \textbf{Classificateur obligatoire} après \textbf{chaque} nombre (y compris le résultat).
    \\ \zh{五块加三块等于八块。} (Wǔ kuài jiā sān kuài děngyú bā kuài.) — 5 yuans + 3 yuans = 8 yuans.
    \\ \zh{十元加五元等于十五元。} (Shí yuán jiā wǔ yuán děngyú shíwǔ yuán.) — 10 yuans + 5 yuans = 15 yuans.
\end{itemize}
\cle{Règle d'or} : Dans la structure \zh{N1 加 N2 等于 N3}, si N1 a un classificateur, N2 et N3 \textbf{doivent} avoir le même classificateur.
\end{attention}

\begin{exemplebox}[Application aux monnaies Cameroun / Chine]
\begin{itemize}
    \item \zh{五百弗朗加五百弗朗等于一千弗朗。} (Wǔbǎi Fǎláng jiā wǔbǎi Fǎláng děngyú yīqiān Fǎláng.) — 500 FCFA + 500 FCFA = 1000 FCFA. (Note : \zh{弗朗} Fǎláng = Franc, souvent utilisé sans classificateur spécifique ou avec \zh{元} yuán par analogie).
    \item \zh{十元加二十元等于三十元。} (Shí yuán jiā èrshí yuán děngyú sānshí yuán.) — 10 ¥ + 20 ¥ = 30 ¥.
    \item \zh{一块五加两块五等于四块。} (Yī kuài wǔ jiā liǎng kuài wǔ děngyú sì kuài.) — 1,50 ¥ + 2,50 ¥ = 4 ¥.
\end{itemize}
\end{exemplebox}

\courssub{Vocabulaire clé de l'addition}

\begin{center}
\begin{tabularx}{\linewidth}{|X|X|X|X|}
\hline
\rowcolor{popblueL} \textbf{Caractères} & \textbf{Pinyin} & \textbf{Nature} & \textbf{Traduction} \\ \hline
\zh{加} & jiā & verbe & ajouter, plus (+) \\ \hline
\zh{等于} & děngyú & verbe & être égal à (=) \\ \hline
\zh{是} & shì & verbe & être (variante orale de 等于) \\ \hline
\zh{几} & jǐ & pr. interrog. & combien (<10) \\ \hline
\zh{多少} & duōshao & pr. interrog. & combien (≥10) \\ \hline
\zh{一共} & yīgòng & adv. & au total, en tout \\ \hline
\zh{块} & kuài & classif. & yuan (口语, pièce/monnaie) \\ \hline
\zh{元} & yuán & classif. & yuan (formel, unité monétaire) \\ \hline
\zh{角} & jiǎo & classif. & jiao (0,1 yuan) \\ \hline
\zh{分} & fēn & classif. & fen (0,01 yuan) \\ \hline
\zh{弗朗} & Fǎláng & nom & Franc (CFA) \\ \hline
\end{tabularx}
\end{center}

\courssub{Méthode : Poser et répondre à une question de calcul}

\begin{methode}[Étapes pour formuler une addition en chinois]
\begin{enumerate}
    \item \textbf{Identifier les deux termes} (N1, N2) et leur nature (nombres purs ou sommes d'argent).
    \item \textbf{Choisir le registre} : \zh{等于} (écrit/formel) ou \zh{是} (oral/familier).
    \item \textbf{Appliquer le schéma} : N1 + \zh{加} + N2 + \zh{等于/是} + Résultat.
    \item \textbf{Gérer les classificateurs} : Si c'est de l'argent, ajouter le même classificateur (\zh{块}, \zh{元}, \zh{弗朗}) après \textbf{chaque} nombre (N1, N2, Résultat).
    \item \textbf{Pour une question} : Remplacer le Résultat par \zh{几} (résultat < 10) ou \zh{多少} (résultat ≥ 10).
\end{enumerate}
\end{methode}

\courssub{Comparaison avec le français et synthèse des pièges}

\begin{attention}[Erreurs fréquentes des francophones — Liste de contrôle]
\begin{enumerate}
    \item \textbf{Traduire « font » par \zh{做} (zuò, faire).}
    \\ \textcolor{red}{FAUX} : \zh{二加三做五。}
    \\ \textcolor{green}{JUSTE} : \zh{二加三等于五。} / \zh{二加三是五。}
    \item \textbf{Oublier le classificateur sur le résultat (somme d'argent).}
    \\ \textcolor{red}{FAUX} : \zh{五块加三块等于八。}
    \\ \textcolor{green}{JUSTE} : \zh{五块加三块等于八块。}
    \item \textbf{Mélanger les classificateurs.}
    \\ \textcolor{red}{FAUX} : \zh{五块加三元等于八块。}
    \\ \textcolor{green}{JUSTE} : Unifier l'unité (\zh{块} ou \zh{元}).
    \item \textbf{Inverser l'ordre (calque de l'anglais "5 plus 3") :} L'ordre chinois est \textbf{strictement} N1 加 N2. Ne jamais dire \zh{加五三}.
    \item \textbf{Utiliser \zh{和} (hé, « et ») au lieu de \zh{加}.} \zh{和} lie des noms, pas des nombres dans un calcul. \textcolor{red}{FAUX} : \zh{五和三等于八。}
\end{enumerate}
\end{attention}

\begin{definition}[\zh{等于} (děngyú)]
Verbe d'état / relation mathématique signifiant « être égal à ».
\\ \textbf{Caractéristiques :} Invariable (pas de conjugaison, pas de particule d'aspect comme \zh{了} ici). Structure fixe : [Opération] \zh{等于} [Résultat].
\\ \textbf{Ne pas confondre avec :} \zh{等} (děng, attendre / classe / grade) ; \zh{于} (yú, préposition littéraire « en / à / par »).
\end{definition}

\courssub{Exercices d'application résolus}

\begin{exoresolu}[Exercice 1 : Transformer en phrase chinoise]
\textbf{Énoncé :} Écrivez en caractères, pinyin et français les additions suivantes.
\begin{enumerate}
    \item 7 + 3 = 10 (calcul pur)
    \item 5 yuans + 5 yuans = 10 yuans (argent, registre formel)
    \item 100 + 200 = ? (question, calcul pur)
    \item 8 + 2 = ? (question, calcul pur, registre oral)
\end{enumerate}
\tcblower
\textbf{Solution détaillée :}
\begin{enumerate}
    \item \zh{七加三等于十。} (Qī jiā sān děngyú shí.) — 7 + 3 = 10.
    \item \zh{五元加五元等于十元。} (Wǔ yuán jiā wǔ yuán děngyú shí yuán.) — 5 ¥ + 5 ¥ = 10 ¥. (Utilisation de \zh{元} formel et répétition sur les 3 termes).
    \item \zh{一百加两百等于多少？} (Yībǎi jiā liǎngbǎi děngyú duōshao?) — 100 + 200 = combien ? (Résultat > 10 $\rightarrow$ \zh{多少}).
    \item \zh{八加二是几？} (Bā jiā èr shì jǐ?) — 8 + 2 = combien ? (Registre oral $\rightarrow$ \zh{是} ; Résultat < 10 $\rightarrow$ \zh{几}).
\end{enumerate}
\end{exoresolu}

\begin{exoresolu}[Exercice 2 : Corriger les erreurs]
\textbf{Énoncé :} Chaque phrase contient une erreur (grammaire, classificateur, registre). Corrigez-la.
\begin{enumerate}
    \item \zh{十加五做十五。}
    \item \zh{三块加四块等于七。}
    \item \zh{二十和三十等于五十。}
    \item \zh{六加四等于十块。}
\end{enumerate}
\tcblower
\textbf{Corrections et explications :}
\begin{enumerate}
    \item \zh{十加五等于十五。} (ou \zh{是十五}). \cle{Erreur} : \zh{做} ne s'utilise pas pour « font » en maths.
    \item \zh{三块加四块等于七块。} \cle{Erreur} : Manque le classificateur \zh{块} sur le résultat.
    \item \zh{二十加三十等于五十。} \cle{Erreur} : \zh{和} lie des noms, pas des termes d'addition. Opérateur = \zh{加}.
    \item \zh{六加四等于十。} (calcul pur) \textbf{OU} \zh{六块加四块等于十块。} (argent). \cle{Erreur} : Incohérence. Si pas de classificateur au début, pas à la fin. Si argent, classificateur partout.
\end{enumerate}
\end{exoresolu}

\courssub{Exercices d'entraînement (à faire en classe ou à la maison)}

\begin{exercice}[Entraînement : L'addition en situation]
\begin{enumerate}
    \item Traduisez en chinois (caractères + pinyin) :
    \begin{itemize}
        \item[a)] 12 + 8 = 20.
        \item[b)] 500 FCFA + 500 FCFA = 1000 FCFA.
        \item[c)] 3,50 ¥ + 1,50 ¥ = 5 ¥ (utilisez \zh{块} et \zh{角} ou \zh{元}).
    \end{itemize}
    \item Posez la question correspondante pour obtenir la réponse indiquée :
    \begin{itemize}
        \item[a)] Réponse : \zh{等于十五。} (Question avec \zh{多少})
        \item[b)] Réponse : \zh{是九。} (Question avec \zh{几}, registre oral)
    \end{itemize}
    \item Complétez avec \zh{加}, \zh{等于}, \zh{是}, \zh{几}, \zh{多少} ou le bon classificateur (\zh{块}, \zh{元}) :
    \begin{itemize}
        \item[a)] 百 \_\_\_\_ 五十 \_\_\_\_ 一百五十。
        \item[b)] 十 \_\_\_\_ \_\_\_\_ \_\_\_\_ 十块。
        \item[c)] 八 \_\_\_\_ 二 \_\_\_\_ \_\_\_\_ ？
        \item[d)] 三 \_\_\_\_ 七 \_\_\_\_ 十。
    \end{itemize}
\end{enumerate}
\end{exercice}

\begin{corrige}[Exercice : Entraînement]
\begin{enumerate}
    \item \textbf{Traduction :}
    \begin{itemize}
        \item[a)] \zh{十二加八等于二十。} (Shí'èr jiā bā děngyú èrshí.)
        \item[b)] \zh{五百弗朗加五百弗朗等于一千弗朗。} (Wǔbǎi Fǎláng jiā wǔbǎi Fǎláng děngyú yīqiān Fǎláng.)
        \item[c)] \zh{三块五加一块五等于五块。} (Sān kuài wǔ jiā yī kuài wǔ děngyú wǔ kuài.) \textit{Note : 3,50 = 三块五 ; 1,50 = 一块五.}
    \end{itemize}
    \item \textbf{Questions :}
    \begin{itemize}
        \item[a)] \zh{十加五等于多少？} (Shí jiā wǔ děngyú duōshao?)
        \item[b)] \zh{八加二是几？} (Bā jiā èr shì jǐ?)
    \end{itemize}
    \item \textbf{Complétion :}
    \begin{itemize}
        \item[a)] 百 \zh{加} 五十 \zh{等于} 一百五十。
        \item[b)] 十 \zh{块} \zh{加} 五 \zh{块} \zh{等于} 十块。 (ou \zh{十块加五块等于十五块} si le premier terme est 10).
        \item[c)] 八 \zh{加} 二 \zh{等于几} / \zh{是几} ？
        \item[d)] 三 \zh{加} 七 \zh{等于} / \zh{是} 十。
    \end{itemize}
\end{enumerate}
\end{corrige}

\coursec{Structure 2 : la soustraction 减法}

Dans la section précédente, nous avons appris à additionner avec \zh{加} (jiā) : \zh{三加五等于八。} (Sān jiā wǔ děngyú bā. « 3 + 5 = 8 »). Mais au marché de Mokolo ou de Douala, on ne fait pas qu'additionner : quand on tend un billet de 100 yuans pour un article à 80 yuans, il faut \emph{retirer}, calculer la monnaie à rendre. C'est exactement le rôle de la soustraction, \zh{减法} (jiǎnfǎ). Bonne nouvelle : la structure est identique à celle de l'addition ; seul le mot-outil change.

\courssub{Observation : rendre la monnaie au marché}

\begin{textebox}[Dialogue : au marché de Yaoundé]
顾客：老板，这个包多少钱？\\
Gùkè: Lǎobǎn, zhège bāo duōshao qián?\\
« Client : Patron, combien coûte ce sac ? »\\
老板：八十块。\\
Lǎobǎn: Bāshí kuài.\\
« Vendeur : 80 yuans. »\\
顾客：太贵了！便宜一点吧！\\
Gùkè: Tài guì le! Piányi yīdiǎn ba!\\
« Client : Trop cher ! Un peu moins cher, s'il vous plaît ! »\\
老板：好，打九折，七十二块。\\
Lǎobǎn: Hǎo, dǎ jiǔ zhé, qīshí'èr kuài.\\
« Vendeur : D'accord, réduction de 10 \%, 72 yuans. »\\
顾客：我给你一百块。\\
Gùkè: Wǒ gěi nǐ yībǎi kuài.\\
« Client : Je te donne 100 yuans. »\\
老板：一百减七十二等于二十八，找你二十八块。\\
Lǎobǎn: Yībǎi jiǎn qīshí'èr děngyú èrshíbā, zhǎo nǐ èrshíbā kuài.\\
« Vendeur : 100 moins 72 égale 28, je te rends 28 yuans. »\\
顾客：谢谢！\\
Gùkè: Xièxie!\\
« Client : Merci ! »
\end{textebox}

Observons la phrase clé : \zh{一百减七十二等于二十八。} (Yībǎi jiǎn qīshí'èr děngyú èrshíbā. « 100 − 72 = 28 »). On reconnaît le squelette de l'addition, mais le verbe \zh{加} a été remplacé par \zh{减}.

\begin{definition}[Le mot-outil 减]
\zh{减} (jiǎn), verbe, signifie « soustraire, retirer, diminuer ». Il est composé de la clé de la glace \zh{冫} à gauche et de \zh{咸} (xián) à droite. On le retrouve dans \zh{减法} (jiǎnfǎ, la soustraction), \zh{减少} (jiǎnshǎo, diminuer) et \zh{减肥} (jiǎnféi, maigrir, littéralement « réduire la graisse »).
\end{definition}

\courssub{Schéma et emploi de la soustraction}

\begin{propriete}[Structure de la soustraction]
La soustraction suit \textbf{exactement le même schéma que l'addition} ; seul le mot-outil change : \zh{加} → \zh{减}.
\begin{center}
\textbf{Nombre1 + 减 + Nombre2 + 等于 + Nombre3}\\
(N1 − N2 = N3)
\end{center}
\begin{itemize}
\item \textbf{Nombre1} : le nombre de départ, celui duquel on retire (le \emph{grand} nombre en général) ;
\item \zh{减} (jiǎn) : « moins » ;
\item \textbf{Nombre2} : la quantité retirée ;
\item \zh{等于} (děngyú) : « égale, font » ;
\item \textbf{Nombre3} : le résultat, c'est-à-dire le \emph{reste}.
\end{itemize}
\end{propriete}

\begin{popfigure}
\begin{tikzpicture}[node distance=0.5cm]
\node[draw=popblue, fill=popblueL, rounded corners, minimum width=1.6cm, minimum height=0.9cm] (n1) {N1};
\node[draw=popdark, fill=popcream, rounded corners, right=of n1, minimum width=1.1cm, minimum height=0.9cm] (jian) {减};
\node[draw=poporange, fill=poporangeL, rounded corners, right=of jian, minimum width=1.6cm, minimum height=0.9cm] (n2) {N2};
\node[draw=popdark, fill=popcream, rounded corners, right=of n2, minimum width=1.3cm, minimum height=0.9cm] (deng) {等于};
\node[draw=popgreen, fill=popgreenL, rounded corners, right=of deng, minimum width=1.6cm, minimum height=0.9cm] (n3) {N3};
\draw[-{Stealth}, very thick, popgreen] (deng.south) to[bend left=25] node[below, popgreen, font=\small]{reste 剩余} (n3.south);
\node[below=0.15cm of n1, font=\small, popblue]{départ};
\node[below=0.15cm of n2, font=\small, poporange]{retiré};
\end{tikzpicture}
\legende{Schéma de la soustraction : N1 减 N2 等于 N3 ; N3 est le reste.}
\end{popfigure}

\textbf{Emplois principaux :}
\begin{enumerate}
\item \textbf{Calculer un reste} : \zh{十减六等于四。} (Shí jiǎn liù děngyú sì. « 10 − 6 = 4 »).
\item \textbf{Calculer la monnaie à rendre} : argent donné − prix = monnaie.
\item \textbf{Calculer une réduction de prix} : ancien prix − réduction = nouveau prix.
\end{enumerate}

\begin{exemplebox}[Exemples traduits]
\begin{itemize}
\item \zh{十减六等于四。} \emph{Shí jiǎn liù děngyú sì.} « 10 − 6 = 4. »
\item \zh{一百减二十五等于七十五。} \emph{Yībǎi jiǎn èrshíwǔ děngyú qīshíwǔ.} « 100 − 25 = 75. »
\item \zh{五十块减十块等于四十块。} \emph{Wǔshí kuài jiǎn shí kuài děngyú sìshí kuài.} « 50 yuans − 10 yuans = 40 yuans. »
\item \zh{我给你一百块，东西八十块，一百减八十等于二十，找你二十块。} \emph{Wǒ gěi nǐ yībǎi kuài, dōngxi bāshí kuài, yībǎi jiǎn bāshí děngyú èrshí, zhǎo nǐ èrshí kuài.} « Je te donne 100 yuans, l'article coûte 80 yuans, 100 − 80 = 20, je te rends 20 yuans. »
\item \zh{这件衣服一百二十块，打八折，一百二十减二十四等于九十六。} \emph{Zhè jiàn yīfu yībǎi'èrshí kuài, dǎ bā zhé, yībǎi'èrshí jiǎn èrshísì děngyú jiǔshíliù.} « Ce vêtement coûte 120 yuans, réduction de 20 \%, 120 − 24 = 96. »
\end{itemize}
\end{exemplebox}

\courssub{La forme interrogative : poser la question}

Pour demander le résultat d'une soustraction, on remplace N3 par le mot interrogatif \zh{几} (jǐ, « combien ») — ou \zh{多少} (duōshao) pour les grands nombres :

\begin{center}
\textbf{Nombre1 + 减 + Nombre2 + 等于 + 几 / 多少 ?}
\end{center}

\begin{exemplebox}[Questions et réponses]
\begin{itemize}
\item \zh{十减四等于几？} \emph{Shí jiǎn sì děngyú jǐ?} « 10 − 4 font combien ? » → \zh{十减四等于六。} (Shí jiǎn sì děngyú liù. « 10 − 4 = 6. »)
\item \zh{一百减八十等于多少？} \emph{Yībǎi jiǎn bāshí děngyú duōshao?} « 100 − 80 font combien ? » → \zh{等于二十。} (Děngyú èrshí. « Ça fait 20. »)
\item \zh{三十减十五等于几？} \emph{Sānshí jiǎn shíwǔ děngyú jǐ?} « 30 − 15 font combien ? » → \zh{等于十五。} (Děngyú shíwǔ. « 15. »)
\end{itemize}
\end{exemplebox}

\courssub{Application : calculer la monnaie et les réductions}

\begin{methode}[Calculer la monnaie à rendre en trois étapes]
\begin{enumerate}
\item Identifier l'\textbf{argent donné} (\zh{给的钱} gěi de qián) et le \textbf{prix} (\zh{价格} jiàgé).
\item Appliquer la formule : \textbf{给的钱 减 价格 等于 找的钱} (argent donné − prix = monnaie rendue).
\item Formuler la réponse complète avec \zh{找你……块} (zhǎo nǐ ... kuài, « je te rends ... yuans »).
\end{enumerate}
Exemple : tu donnes 50 yuans pour un cahier à 35 yuans → \zh{五十减三十五等于十五，找你十五块。} (Wǔshí jiǎn sānshíwǔ děngyú shíwǔ, zhǎo nǐ shíwǔ kuài. « 50 − 35 = 15, je te rends 15 yuans. »)
\end{methode}

\begin{popfigure}
\begin{tikzpicture}[node distance=0.9cm and 1.4cm]
\node[draw=popblue, fill=popblueL, rounded corners, align=center, minimum width=3.4cm, minimum height=1cm] (a) {Donner le billet\\给钱 gěi qián};
\node[draw=poporange, fill=poporangeL, rounded corners, align=center, minimum width=3.4cm, minimum height=1cm, below=of a] (b) {Calculer\\减 jiǎn : donné − prix};
\node[draw=popgreen, fill=popgreenL, rounded corners, align=center, minimum width=3.4cm, minimum height=1cm, below=of b] (c) {Rendre la monnaie\\找钱 zhǎo qián};
\draw[-{Stealth}, very thick, popdark] (a) -- (b);
\draw[-{Stealth}, very thick, popdark] (b) -- (c);
\end{tikzpicture}
\legende{Organigramme du dialogue d'achat : donner → soustraire → rendre.}
\end{popfigure}

\begin{definition}[打折 : la réduction, vocabulaire d'examen]
\zh{打折} (dǎzhé), « faire une réduction ». Attention, la logique chinoise est inverse de la française : \zh{打九折} (dǎ jiǔ zhé) signifie « on paie 90 \% du prix », donc une réduction de 10 \% ; \zh{打八折} = on paie 80 \% (réduction de 20 \%). Utile dans tout dialogue d'achat à l'examen.
\end{definition}

\begin{center}
\begin{tabularx}{\linewidth}{|X|X|X|X|}
\hline
\rowcolor{popblueL} Caractères & Pinyin & Nature & Traduction \\
\hline
减 & jiǎn & verbe & soustraire, moins \\
\hline
减法 & jiǎnfǎ & nom & la soustraction \\
\hline
等于 & děngyú & verbe & égaler, faire \\
\hline
找钱 & zhǎo qián & verbe + nom & rendre la monnaie \\
\hline
价格 & jiàgé & nom & le prix \\
\hline
打折 & dǎzhé & verbe & faire une réduction \\
\hline
便宜 & piányi & adjectif & bon marché \\
\hline
剩余 & shèngyú & nom/verbe & le reste, rester \\
\hline
\end{tabularx}
\end{center}

\courssub{Comparaison avec le français et pièges}

\begin{attention}[Ne jamais inverser l'ordre des nombres]
En français, on peut dire « 4 soustrait de 10 » ou « retirer 4 de 10 », ce qui place le petit nombre en premier. Les francophones transposent parfois cet ordre en chinois : erreur ! En chinois, on dit \textbf{toujours} le nombre de départ d'abord : \zh{十减四} (shí jiǎn sì), jamais \zh{四减十} pour « 10 − 4 ». Le premier nombre est celui \emph{duquel} on soustrait, exactement comme dans l'opération écrite « 10 − 4 ».
\end{attention}

\begin{center}
\begin{tabularx}{\linewidth}{|X|X|X|}
\hline
\rowcolor{popblueL} Français & Chinois + pinyin & Remarque \\
\hline
10 − 6 = 4 & 十减六等于四。Shí jiǎn liù děngyú sì. & ordre identique \\
\hline
« 4 retiré de 10 » & 十减四 shí jiǎn sì & le chinois garde 10 en tête \\
\hline
Combien font 10 − 4 ? & 十减四等于几？Shí jiǎn sì děngyú jǐ? & 几 remplace le résultat \\
\hline
Je te rends 20 yuans. & 找你二十块。Zhǎo nǐ èrshí kuài. & 找 + personne + somme \\
\hline
\end{tabularx}
\end{center}

\begin{exoresolu}[Calculer la monnaie]
Énoncé : Marie achète un tissu à 65 yuans. Elle donne un billet de 100 yuans. Écrivez en chinois la phrase complète du vendeur (calcul + monnaie rendue).\tcblower
Solution : argent donné = 100, prix = 65. On applique la formule 给的钱 减 价格 等于 找的钱 : 100 − 65 = 35.\\
\zh{一百减六十五等于三十五，找你三十五块。}\\
\emph{Yībǎi jiǎn liùshíwǔ děngyú sānshíwǔ, zhǎo nǐ sānshíwǔ kuài.}\\
« 100 − 65 = 35, je te rends 35 yuans. »\\
Vérification : le grand nombre (100) est bien en première position, \zh{减} relie les deux nombres, \zh{等于} introduit le reste, et \zh{找} + personne + somme conclut la transaction.
\end{exoresolu}

\begin{exercice}[À vous de calculer]
Traduisez en chinois, avec pinyin :
\begin{enumerate}
\item 20 − 8 = 12.
\item 90 − 45 = 45.
\item Combien font 50 − 15 ?
\item Je te donne 200 yuans, le pantalon coûte 150 yuans ; calcule et rends la monnaie.
\end{enumerate}
\end{exercice}

\begin{corrige}[Exercice : à vous de calculer]
\begin{enumerate}
\item \zh{二十减八等于十二。} \emph{Èrshí jiǎn bā děngyú shí'èr.}
\item \zh{九十减四十五等于四十五。} \emph{Jiǔshí jiǎn sìshíwǔ děngyú sìshíwǔ.}
\item \zh{五十减十五等于几？} \emph{Wǔshí jiǎn shíwǔ děngyú jǐ?} (Réponse : \zh{等于三十五。} \emph{Děngyú sānshíwǔ.})
\item \zh{二百减一百五十等于五十，找你五十块。} \emph{Èrbǎi jiǎn yībǎi wǔshí děngyú wǔshí, zhǎo nǐ wǔshí kuài.}
\end{enumerate}
\end{corrige}

\begin{savaistu}[Négocier les prix : Cameroun et Chine]
Au Cameroun, marchander au marché est un art de vivre : on discute, on sourit, on propose un prix. En Chine, c'est pareil dans les marchés et les petites boutiques ! La phrase magique à retenir : \zh{便宜一点吧！} (Piányi yīdiǎn ba! « Un peu moins cher, s'il vous plaît ! »). Le vendeur répond souvent par un \zh{打折} : \zh{打九折} signifie que vous payez 90 \% du prix affiché. Attention donc : plus le chiffre devant \zh{折} est petit, plus la réduction est grande — \zh{打五折} (dǎ wǔ zhé), c'est moitié prix !
\end{savaistu}

\begin{aretenir}[L'essentiel de la soustraction]
\begin{itemize}
\item Schéma : \textbf{N1 + 减 + N2 + 等于 + N3} — même squelette que l'addition, seul \zh{加} devient \zh{减}.
\item Le premier nombre est toujours celui duquel on retire : pas d'inversion comme dans « 4 retiré de 10 ».
\item Question : remplacer le résultat par \zh{几} ou \zh{多少}.
\item Monnaie : \textbf{给的钱 减 价格 等于 找的钱}, puis \zh{找你……块}.
\item Réductions : \zh{打九折} = payer 90 \% du prix.
\end{itemize}
\end{aretenir}

\coursec{Les monnaies et les nombres}

Dans les sections précédentes, nous avons appris à calculer : \zh{三加五等于八。} (\emph{Sān jiā wǔ děngyú bā.} « Trois plus cinq égale huit. ») et \zh{十减四等于六。} (\emph{Shí jiǎn sì děngyú liù.} « Dix moins quatre égale six. »). Mais au marché de Mokolo à Yaoundé ou dans un magasin de Pékin, on ne calcule pas dans le vide : on calcule des \cle{prix}, avec une \cle{monnaie}. Cette section vous donne tous les outils pour parler d'argent en chinois : les noms des monnaies, la façon de dire un prix, et surtout la grande différence entre le français et le chinois dans la lecture des grands nombres.

\courssub{Observer : un dialogue au magasin}

\begin{textebox}[Au magasin de Douala — 在杜阿拉的商店]
顾客：这个书包多少钱？\\
\emph{Gùkè: Zhège shūbāo duōshao qián?}\\
« Client : Combien coûte ce cartable ? »\\
售货员：两百五十块。\\
\emph{Shòuhuòyuán: Liǎngbǎi wǔshí kuài.}\\
« Vendeuse : Deux cent cinquante yuans. »\\
顾客：太贵了！那支笔呢？\\
\emph{Gùkè: Tài guì le! Nà zhī bǐ ne?}\\
« Client : C'est trop cher ! Et ce stylo ? »\\
售货员：笔十五块五毛。\\
\emph{Shòuhuòyuán: Bǐ shíwǔ kuài wǔ máo.}\\
« Vendeuse : Le stylo fait quinze yuans cinquante. »\\
顾客：好，我买笔。给您钱。\\
\emph{Gùkè: Hǎo, wǒ mǎi bǐ. Gěi nín qián.}\\
« Client : Bien, j'achète le stylo. Voici l'argent (politesse). »
\end{textebox}

Observons. Pour demander un prix, on dit \zh{多少钱？} (\emph{duōshao qián?} « combien d'argent ? »). Pour répondre, le vendeur dit \zh{两百五十块} : nombre + \zh{块}. Et pour les petites sommes, on ajoute \zh{毛} : \zh{十五块五毛} = 15,50 yuans. Nous allons maintenant étudier chaque élément.

\courssub{La monnaie chinoise : 人民币}

\begin{definition}[人民币 : la monnaie du peuple]
La \cle{monnaie officielle} de la République populaire de Chine s'appelle \zh{人民币} (\emph{rénmínbì}, littéralement « monnaie du peuple »). Son symbole est \textbf{¥} (code international : CNY). Son unité de base est \zh{元} (\emph{yuán}) ; à l'oral, on dit presque toujours \zh{块} (\emph{kuài}). Un yuan se divise en dix \zh{毛} (\emph{máo}, à l'écrit : \zh{角} \emph{jiǎo}) : \textbf{1 块 = 10 毛}.
\end{definition}

\begin{propriete}[Dire un prix : nombre + 块 (+ nombre + 毛)]
Structure : \textbf{Sujet + nombre + 块 (+ nombre + 毛)}.\par
\zh{这个多少钱？} \emph{Zhège duōshao qián?} « Combien coûte ceci ? »\\
\zh{五十块。} \emph{Wǔshí kuài.} « Cinquante yuans. »\\
\zh{三块五毛。} \emph{Sān kuài wǔ máo.} « Trois yuans cinquante (3,50 ¥). »\par
À l'écrit officiel (étiquettes, banques, chèques), on utilise \zh{元} et \zh{角} : \zh{三元五角}. À l'oral quotidien : \zh{三块五毛}. Les deux sont corrects ; retenez la paire \zh{块/毛} pour parler, \zh{元/角} pour lire les prix affichés.
\end{propriete}

\begin{exemplebox}[Prix en yuans]
\zh{这本书二十块。} \emph{Zhè běn shū èrshí kuài.} « Ce livre coûte vingt yuans. »\\
\zh{一瓶水两块。} \emph{Yì píng shuǐ liǎng kuài.} « Une bouteille d'eau coûte deux yuans. »\\
\zh{这张票一百八十块。} \emph{Zhè zhāng piào yìbǎi bāshí kuài.} « Ce billet coûte cent quatre-vingts yuans. »\\
\zh{一部手机三万块？太贵了！} \emph{Yì bù shǒujī sānwàn kuài? Tài guì le!} « Un téléphone portable à trente mille yuans ? C'est trop cher ! »
\end{exemplebox}

\courssub{Les monnaies du monde}

Le mot « monnaie » ou « argent » se dit \zh{钱} (\emph{qián}). Chaque monnaie étrangère a un nom composé, souvent formé du nom du pays ou d'une transcription phonétique.

\begin{center}
\begin{tabularx}{\linewidth}{|X|X|X|X|}
\hline
\rowcolor{popblueL} Caractères & Pinyin & Nature & Traduction \\
\hline
钱 & qián & nom & argent, monnaie \\
\hline
人民币 & rénmínbì & nom & renminbi (monnaie chinoise) \\
\hline
元 / 块 & yuán / kuài & classif. & yuan (écrit / oral) \\
\hline
毛 / 角 & máo / jiǎo & classif. & dixième de yuan (oral / écrit) \\
\hline
美元 & měiyuán & nom & dollar américain \\
\hline
欧元 & ōuyuán & nom & euro \\
\hline
英镑 & yīngbàng & nom & livre sterling \\
\hline
日元 & rìyuán & nom & yen japonais \\
\hline
法郎 & fǎláng & nom & franc (FCFA) \\
\hline
奈拉 & nàilā & nom & naira (Nigeria) \\
\hline
塞地 & sàidì & nom & cedi (Ghana) \\
\hline
\end{tabularx}
\end{center}

Observez la logique de formation : \zh{美元} = \zh{美} (Amérique) + \zh{元} ; \zh{欧元} = \zh{欧} (Europe) + \zh{元} ; \zh{日元} = \zh{日} (Japon) + \zh{元}. En revanche, \zh{法郎}, \zh{奈拉} et \zh{塞地} sont des \cle{transcriptions phonétiques} : les caractères imitent le son (« fa-lang », « nai-la », « sai-di ») sans lien avec leur sens propre.

\begin{exemplebox}[Avec les monnaies étrangères et africaines]
\zh{我有一千美元。} \emph{Wǒ yǒu yìqiān měiyuán.} « J'ai mille dollars. »\\
\zh{这台电脑五百欧元。} \emph{Zhè tái diànnǎo wǔbǎi ōuyuán.} « Cet ordinateur coûte cinq cents euros. »\\
\zh{这个包五万法郎。} \emph{Zhège bāo wǔwàn fǎláng.} « Ce sac coûte cinquante mille francs. »\\
\zh{他去尼日利亚，换了两万奈拉。} \emph{Tā qù Nírìlìyà, huàn le liǎngwàn nàilā.} « Il est allé au Nigeria et a changé vingt mille nairas. »
\end{exemplebox}

Notez bien : avec les monnaies étrangères, on dit directement \textbf{nombre + monnaie} (\zh{一千美元}, \zh{五万法郎}), sans \zh{块}. Le mot \zh{块} ne s'emploie qu'avec le yuan chinois à l'oral.

\begin{popfigure}
\begin{tikzpicture}[mindmap, grow cyclic, every node/.style={concept}, root concept/.append style={concept color=popblue, font=\Large\bfseries}, level 1/.append style={level distance=3.2cm, sibling angle=51.5}, concept color=poporangeL]
\node[root concept, align=center]{钱 qián\\ argent}
  child { node[align=center]{人民币 rénmínbì\\ (Chine)} }
  child { node[align=center]{美元 měiyuán\\ (États-Unis)} }
  child { node[align=center]{欧元 ōuyuán\\ (Europe)} }
  child { node[align=center]{英镑 yīngbàng\\ (Royaume-Uni)} }
  child { node[align=center]{法郎 fǎláng\\ (FCFA)} }
  child { node[align=center]{奈拉 nàilā\\ (Nigeria)} }
  child { node[align=center]{塞地 sàidì\\ (Ghana)} };
\end{tikzpicture}
\legende{Carte mentale : les monnaies autour du mot 钱 (qián).}
\end{popfigure}

\begin{savaistu}[Le franc CFA en chinois]
Le Cameroun utilise le franc CFA d'Afrique centrale. En chinois, on dit \zh{中非法郎} (\emph{Zhōngfēi fǎláng}, « franc d'Afrique centrale » : \zh{中非} = Afrique centrale, \zh{法郎} = franc). Le franc CFA d'Afrique de l'ouest se dit \zh{西非法郎} (\emph{Xīfēi fǎláng}). Un Chinois qui visite le marché de Mvog-Mbi à Yaoundé paiera donc en \zh{中非法郎} !
\end{savaistu}

\courssub{Révision active des nombres : le système des 万}

Pour dire les prix, il faut maîtriser les grands nombres. Rappel des unités : \zh{一} 1, \zh{二} 2, \zh{三} 3, \zh{四} 4, \zh{五} 5, \zh{六} 6, \zh{七} 7, \zh{八} 8, \zh{九} 9, \zh{十} 10. Puis viennent les puissances : \zh{百} (\emph{bǎi}, 100), \zh{千} (\emph{qiān}, 1 000), \zh{万} (\emph{wàn}, 10 000).

\begin{propriete}[Le chinois compte par 万, pas par milliers]
Le français groupe les grands nombres par \textbf{milliers} (mille, million, milliard). Le chinois les groupe par \textbf{dix mille} : \zh{万}. Il n'existe pas de mot chinois pour « million » : on dit \zh{一百万} (« cent fois dix mille »).\par
\begin{center}
\begin{tabularx}{\linewidth}{|X|X|X|}
\hline
\rowcolor{popblueL} Nombre & Chinois + pinyin & Lecture \\
\hline
10 & 十 shí & dix \\
\hline
100 & 一百 yìbǎi & cent \\
\hline
1 000 & 一千 yìqiān & mille \\
\hline
10 000 & 一万 yíwàn & dix mille \\
\hline
100 000 & 十万 shíwàn & cent mille (dix 万) \\
\hline
1 000 000 & 一百万 yìbǎi wàn & un million (cent 万) \\
\hline
350 000 & 三十五万 sānshíwǔ wàn & trois cent cinquante mille \\
\hline
\end{tabularx}
\end{center}
\end{propriete}

\begin{popfigure}
\begin{tikzpicture}[scale=1, every node/.style={font=\small}]
\draw[-{Stealth}, thick, popline] (0,0) -- (13.5,0);
\foreach \x/\zhn/\py/\val in {0.4/十/shí/10, 2.6/百/bǎi/100, 4.8/千/qiān/1\,000, 7.0/万/wàn/10\,000, 9.4/十万/shíwàn/100\,000, 12.2/百万/bǎiwàn/1\,000\,000}{
  \fill[poporange] (\x,0) circle (2.5pt);
  \node[above, align=center, popdark] at (\x,0.15) {\textbf{\zhn}\\ \emph{\py}};
  \node[below, popmuted] at (\x,-0.15) {\val};
}
\end{tikzpicture}
\legende{Frise des puissances : après 千 (1 000), tout se compte en 万 (10 000).}
\end{popfigure}

\begin{methode}[Lire un grand nombre : découper par tranches de 万]
\begin{enumerate}
\item Repérez le chiffre des \textbf{dizaines de mille} : tout ce qui est à gauche du chiffre des mille forme le nombre de \zh{万}.
\item Lisez ce bloc comme un nombre simple, puis ajoutez \zh{万}.
\item Lisez le reste (milliers, centaines, dizaines) normalement.
\end{enumerate}
Exemple : 350 000 → 35 devant le bloc des mille → \zh{三十五万} (\emph{sānshíwǔ wàn}). Autre exemple : 48 500 → 4 万 + 8 500 → \zh{四万八千五百} (\emph{sìwàn bāqiān wǔbǎi}).
\end{methode}

\begin{exoresolu}[Traduire des prix en chinois]
Traduisez : a) « Ce sac coûte 50 000 francs. » b) « J'ai 1 000 dollars. » c) « Ce livre coûte 20 yuans. » d) « La voiture coûte 1 200 000 francs. »
\tcblower
a) 50 000 = 5 × 10 000 → \zh{这个包五万法郎。} \emph{Zhège bāo wǔwàn fǎláng.} (et non \zh{五十千} : le chinois ne dit jamais « cinquante mille » avec \zh{千}).\\
b) \zh{我有一千美元。} \emph{Wǒ yǒu yìqiān měiyuán.} (nombre + monnaie, sans \zh{块}).\\
c) \zh{这本书二十块。} \emph{Zhè běn shū èrshí kuài.} (classificateur \zh{本} pour les livres ; \zh{块} à l'oral).\\
d) 1 200 000 = 120 × 10 000 → \zh{这辆车一百二十万法郎。} \emph{Zhè liàng chē yìbǎi èrshí wàn fǎláng.}
\end{exoresolu}

\courssub{二 ou 两 ? La distinction essentielle}

Le chinois a deux mots pour « deux » : \zh{二} (\emph{èr}) et \zh{两} (\emph{liǎng}).

\begin{propriete}[Emplois de 二 et de 两]
\begin{itemize}
\item \zh{二} : dans les \textbf{calculs}, les numéros, les ordinaux et à l'intérieur des nombres composés : \zh{二加二等于四。} (\emph{Èr jiā èr děngyú sì.} « Deux plus deux égale quatre. »), \zh{十二} (12), \zh{第二} (deuxième).
\item \zh{两} : \textbf{devant un classificateur} : \zh{两个人} (\emph{liǎng ge rén}, « deux personnes »), \zh{两本书} (« deux livres »), \zh{两万块} (« vingt mille yuans »).
\item Devant \zh{百} et \zh{千}, les deux sont acceptés : \zh{两百块} = \zh{二百块} (200 yuans), \zh{两千块} = \zh{二千块} (2000 yuans). Devant \zh{万}, on préfère \zh{两} : \zh{两万块}.
\end{itemize}
\end{propriete}

\begin{attention}[Erreurs fréquentes des francophones]
\begin{itemize}
\item Ne dites jamais \zh{二个人} : devant un classificateur, c'est \zh{两个人}. En revanche, \zh{二万} n'est pas faux, mais \zh{两万} est plus naturel.
\item Le français compte en milliers, le chinois en \zh{万} : « 50 000 FCFA » = \zh{五万法郎}, jamais \zh{五十千法郎} (le mot \zh{千} ne peut pas dépasser neuf mille).
\item N'ajoutez pas \zh{块} aux monnaies étrangères : on dit \zh{一千美元}, pas \zh{一千块美元}.
\item Attention à l'ordre : le nombre précède toujours la monnaie, comme en français, mais sans « de » : \zh{五万法郎} et non \zh{法郎五万}.
\end{itemize}
\end{attention}

\begin{experience}[Bureau de change en classe]
Par groupes de trois : un élève est le changeur, deux autres sont des clients. Les clients veulent changer des \zh{美元}, des \zh{欧元} ou des \zh{法郎} contre des \zh{人民币}. Le changeur annonce : \zh{一美元换七块人民币。} (\emph{Yì měiyuán huàn qī kuài rénmínbì.} « Un dollar s'échange contre sept yuans. »). Les clients calculent : \zh{一百美元换多少人民币？} (« Combien de yuans pour cent dollars ? »). Cherchez d'abord sur Internet ou dans un journal le taux réel du jour : 1 美元 ≈ 多少人民币 ? Puis jouez la scène en chinois.
\end{experience}

\begin{exercice}[À vous de jouer]
\begin{enumerate}
\item Écrivez en chinois (caractères + pinyin) : 80 000 ; 250 000 ; 2 000 000.
\item Traduisez en chinois : « Ces deux livres coûtent cent vingt yuans. » ; « Il a trente mille francs. »
\item Corrigez la phrase erronée : \zh{五十千块太贵了。}
\end{enumerate}
\end{exercice}

\begin{corrige}[Exercice « À vous de jouer »]
1) \zh{八万} (\emph{bāwàn}) ; \zh{二十五万} (\emph{èrshíwǔ wàn}) ; \zh{两百万} (\emph{liǎngbǎi wàn}).\\
2) \zh{这两本书一百二十块。} (\emph{Zhè liǎng běn shū yìbǎi èrshí kuài.}) ; \zh{他有三万法郎。} (\emph{Tā yǒu sānwàn fǎláng.}).\\
3) \zh{五十千} est impossible : au-delà de neuf mille, on passe aux \zh{万}. Correction : \zh{五万块太贵了。} (\emph{Wǔwàn kuài tài guì le.} « Cinquante mille yuans, c'est trop cher. »).
\end{corrige}

\begin{aretenir}[L'essentiel de la section]
\begin{itemize}
\item Monnaie chinoise : \zh{人民币}, unité \zh{元}/\zh{块} ; 1 \zh{块} = 10 \zh{毛}.
\item Demander un prix : \zh{这个多少钱？} ; répondre : nombre + \zh{块} (+ \zh{毛}).
\item Monnaies étrangères : nombre + monnaie, sans \zh{块} : \zh{一千美元}, \zh{五万法郎}.
\item Le chinois compte par \zh{万} (10 000) : 100 000 = \zh{十万}, 1 000 000 = \zh{一百万}.
\item \zh{二} pour calculer, \zh{两} devant un classificateur : \zh{二加二} mais \zh{两个人}, \zh{两万块}.
\end{itemize}
\end{aretenir}

\coursec{Structure 3 : 除了……以外，……}

Dans les sections précédentes, nous avons appris à compter, à additionner (\zh{加法} jiāfǎ) et à soustraire (\zh{减法} jiǎnfǎ), ainsi que les principales monnaies : \zh{人民币} (rénmínbì, yuan chinois), \zh{美元} (měiyuán, dollar US), \zh{欧元} (ōuyuán, euro), \zh{法郎} (fǎláng, franc CFA). Or, dans la vie quotidienne — au marché, à la banque, en planifiant sa semaine —, on a constamment besoin d'exprimer l'inclusion (« en plus de X, aussi Y ») ou l'exclusion (« sauf X, tout le reste »). La structure \zh{除了……以外，……} (chúle … yǐwài, …) est l'outil grammatical standard pour cela. Elle est incontournable au niveau HSK 3-4 et très fréquente aux épreuves du Baccalauréat camerounais.

\courssub{Lexique : les monnaies et verbes clés}

\begin{center}
\begin{tabularx}{\linewidth}{|X|X|X|X|}
\hline
\rowcolor{popblueL} \textbf{Caractères} & \textbf{Pinyin} & \textbf{Nature} & \textbf{Traduction} \\
\hline
人民币 & rénmínbì & nom & yuan chinois (monnaie officielle) \\
\hline
美元 & měiyuán & nom & dollar des États-Unis \\
\hline
欧元 & ōuyuán & nom & euro \\
\hline
法郎 & fǎláng & nom & franc (CFA au Cameroun) \\
\hline
货币 & huòbì & nom & monnaie, devise \\
\hline
算账 & suànzhàng & verbe & faire les comptes, calculer \\
\hline
带 & dài & verbe & apporter (sur soi), emporter \\
\hline
收 & shōu & verbe & accepter (paiement), recevoir \\
\hline
开门 & kāimén & verbe & ouvrir (magasin), être ouvert \\
\hline
除了 & chúle & prép. & à part, excepté, en plus de \\
\hline
以外 & yǐwài & nom de loc. & en dehors de, à l'extérieur de \\
\hline
还 & hái & adv. & aussi, en plus (addition) \\
\hline
也 & yě & adv. & aussi, également (addition) \\
\hline
都 & dōu & adv. & tous, toutes (exclusion : « tous sauf… ») \\
\hline
只 & zhǐ & adv. & seulement, ne … que (exclusion) \\
\hline
\end{tabularx}
\end{center}

\courssub{Observation : au marché de Mokolo}

\begin{textebox}[Dialogue au marché de Mokolo — Zài Mòkēluó shìchǎng]
星期六早上，玛丽去市场买东西。除了蔬菜以外，她还买了水果和鱼。\\
Xīngqīliù zǎoshang, Mǎlì qù shìchǎng mǎi dōngxi. Chúle shūcài yǐwài, tā hái mǎile shuǐguǒ hé yú.\\
« Samedi matin, Marie va au marché faire des courses. En plus des légumes, elle a aussi acheté des fruits et du poisson. »\\[4pt]
她的朋友们都来了，除了保罗以外，大家都带了钱。\\
Tā de péngyoumen dōu láile, chúle Bǎoluó yǐwài, dàjiā dōu dàile qián.\\
« Ses amis sont tous venus ; sauf Paul, tout le monde a apporté de l'argent. »\\[4pt]
除了人民币以外，玛丽还会用美元算账。\\
Chúle rénmínbì yǐwài, Mǎlì hái huì yòng měiyuán suànzhàng.\\
« En plus des yuans, Marie sait aussi calculer en dollars. »
\end{textebox}

Observons les trois phrases soulignées. Dans la première et la troisième, le mot-outil après la virgule est \zh{还} (hái, « aussi ») : les légumes et le yuan sont \emph{compris} dans l'ensemble. Dans la deuxième, le mot-outil est \zh{都} (dōu, « tous ») : Paul est \emph{exclu} du groupe « tout le monde ». C'est ce mot-outil qui décide du sens global de la phrase.

\courssub{Schéma général de la structure}

\begin{propriete}[Schéma de la structure 除了……以外，……]
\cle{除了 + X + 以外，+ (sujet) + 也/还/都/只 + prédicat}
\begin{itemize}
\item \zh{除了} (chúle) : préposition marquant l'élément de référence X (« à part », « excepté », « en plus de ») ;
\item X : nom, pronom, verbe ou proposition (l'élément mis en avant) ;
\item \zh{以外} (yǐwài, « en dehors de ») : ferme le cadre prépositionnel — \textbf{peut être omis} à l'oral et à l'écrit courant ;
\item \textbf{Mot-outil obligatoire après la virgule} :
  \begin{itemize}
  \item \zh{还} (hái) / \zh{也} (yě) $\rightarrow$ \textbf{Addition (inclusion)} : X fait partie de l'ensemble.
  \item \zh{都} (dōu) / \zh{只} (zhǐ) $\rightarrow$ \textbf{Exclusion} : X est mis à part, retiré de l'ensemble.
  \end{itemize}
\end{itemize}
\end{propriete}

\begin{popfigure}
\begin{tikzpicture}[node distance=0.7cm and 1.8cm, every node/.style={align=center, font=\small}]
\node[draw=popdark, fill=popblueL, rounded corners, font=\bfseries\small] (start) {除了 + X + 以外，};
\node[draw=popgreen, fill=popgreenL, rounded corners, below left=1.0cm and 1.5cm of start, align=center] (add){\textbf{还 / 也}\\ ADDITION\\ X est \emph{inclus}};
\node[draw=poporange, fill=poporangeL, rounded corners, below right=1.0cm and 1.5cm of start, align=center] (exc){\textbf{都 / 只}\\ EXCLUSION\\ X est \emph{exclu}};
\draw[-{Stealth}, very thick, popgreen] (start) -- (add);
\draw[-{Stealth}, very thick, poporange] (start) -- (exc);
\node[below=0.3cm of add, font=\small] {\zh{除了汉语以外，他还学英语。}};
\node[below=0.3cm of exc, font=\small] {\zh{除了小王以外，大家都来了。}};
\end{tikzpicture}
\legende{Les deux branches de la structure : le mot-outil après la virgule tranche le sens.}
\end{popfigure}

\courssub{Emploi 1 : l'addition (sens inclusif)}

Quand le mot-outil est \zh{还} (hái) ou \zh{也} (yě), la structure signifie « en plus de X, aussi Y ». X est \textbf{compris} dans la réalisation de l'action ou dans l'ensemble décrit.

\begin{exemplebox}[Addition : en plus de X, aussi Y]
\zh{除了人民币以外，我还带了美元。}\\
\emph{Chúle rénmínbì yǐwài, wǒ hái dàile měiyuán.}\\
« En plus des yuans, j'ai aussi apporté des dollars. » (j'ai les deux monnaies)\\[4pt]
\zh{除了汉语以外，他也学英语。}\\
\emph{Chúle Hànyǔ yǐwài, tā yě xué Yīngyǔ.}\\
« En plus du chinois, il apprend aussi l'anglais. »\\[4pt]
\zh{除了英语以外，我们还会说法语。}\\
\emph{Chúle Yīngyǔ yǐwài, wǒmen hái huì shuō Fǎyǔ.}\\
« En plus de l'anglais, nous savons aussi parler français. »\\[4pt]
\zh{除了恩多雷以外，她还做了米饭。}\\
\emph{Chúle Ēnduōléi yǐwài, tā hái zuòle mǐfàn.}\\
« En plus du ndolé, elle a aussi préparé du riz. »
\end{exemplebox}

\courssub{Emploi 2 : l'exclusion}

Quand le mot-outil est \zh{都} (dōu, « tous ») ou \zh{只} (zhǐ, « seulement »), la structure signifie « sauf X », « à part X ». X est \textbf{exclu} de l'ensemble ou de l'action.

\begin{exemplebox}[Exclusion : sauf X, à part X]
\zh{除了小王以外，大家都来了。}\\
\emph{Chúle Xiǎo Wáng yǐwài, dàjiā dōu láile.}\\
« Sauf Xiao Wang, tout le monde est venu. » (lui seul est absent)\\[4pt]
\zh{除了数学以外，我只喜欢汉语。}\\
\emph{Chúle shùxué yǐwài, wǒ zhǐ xǐhuan Hànyǔ.}\\
« À part les maths, je n'aime que le chinois. »\\[4pt]
\zh{除了星期天以外，我每天都上课。}\\
\emph{Chúle xīngqītiān yǐwài, wǒ měitiān dōu shàngkè.}\\
« Sauf le dimanche, j'ai cours tous les jours. »\\[4pt]
\zh{除了足球以外，他什么运动都不喜欢。}\\
\emph{Chúle zúqiú yǐwài, tā shénme yùndòng dōu bù xǐhuan.}\\
« À part le football, il n'aime aucun sport. » (négation + 都 = « aucun »)
\end{exemplebox}

\begin{popfigure}
\begin{tikzpicture}
\draw[thick, popblue, fill=popblueL] (0,0) circle (2.2cm);
\draw[thick, poporange, fill=poporangeL, pattern=north east lines, pattern color=poporange] (-1.0,0.6) circle (0.6cm);
\node[popblue, font=\bfseries\small] at (1.2,-1.0) {都：tout le groupe};
\node[poporange, font=\bfseries\small, align=center] at (-1.0,0.6) {X\\(exclu)};
\node[align=center, font=\small] at (0,-3.0) {除了X以外，都…… : X est retiré du grand ensemble};
\end{tikzpicture}
\legende{Visualisation de l'exclusion : le grand cercle = l'ensemble (大家), le cercle hachuré = X (小王) qui en est sorti.}
\end{popfigure}

\courssub{Variantes syntaxiques : omission de 以外 et position du sujet}

\begin{propriete}[Variantes courantes]
\begin{enumerate}
\item \textbf{Omission de \zh{以外}} : très fréquent à l'oral et dans les textes courts, sans changement de sens.\\
\zh{除了汉语，他还学法语。} \emph{Chúle Hànyǔ, tā hái xué Fǎyǔ.} « En plus du chinois, il apprend aussi le français. »
\item \textbf{Sujet placé avant \zh{除了}} : met l'accent sur le sujet (topicalisation).\\
\zh{我除了人民币以外，还带了美元。} = \zh{除了人民币以外，我还带了美元。}\\
\emph{Wǒ chúle rénmínbì yǐwài, hái dàile měiyuán.} « Moi, en plus des yuans, j'ai aussi apporté des dollars. »
\end{enumerate}
\end{propriete}

\begin{center}
\begin{tabularx}{\linewidth}{|X|X|X|}
\hline
\rowcolor{popblueL} \textbf{Structure} & \textbf{Exemple (caractères + pinyin)} & \textbf{Traduction} \\
\hline
除了X以外，还/也… (addition) & 除了茶以外，我也喝咖啡。Chúle chá yǐwài, wǒ yě hē kāfēi. & En plus du thé, je bois aussi du café. \\
\hline
除了X以外，都… (exclusion) & 除了阿里以外，我们都去杜阿拉。Chúle Ālǐ yǐwài, wǒmen dōu qù Dù'ālā. & Sauf Ali, nous allons tous à Douala. \\
\hline
除了X以外，只… (exclusion) & 除了水以外，他只喝果汁。Chúle shuǐ yǐwài, tā zhǐ hē guǒzhī. & À part l'eau, il ne boit que du jus. \\
\hline
以外 omis (addition) & 除了法语，她还会说英语。Chúle Fǎyǔ, tā hái huì shuō Yīngyǔ. & En plus du français, elle sait aussi parler anglais. \\
\hline
Sujet avant 除了 & 他除了篮球以外，还喜欢排球。Tā chúle lánqiú yǐwài, hái xǐhuan páiqiú. & Lui, en plus du basket, il aime aussi le volley. \\
\hline
\end{tabularx}
\end{center}

\begin{attention}[« À part » est ambigu en français, pas en chinois]
En français, « À part le chinois, j'aime les langues » peut signifier « y compris le chinois » (addition) ou « sauf le chinois » (exclusion). En chinois, \textbf{l'ambiguïté n'existe pas} : le choix du mot-outil (\zh{还}/\zh{也} vs \zh{都}/\zh{只}) tranche systématiquement. \cle{Avant de traduire, identifiez toujours le mot-outil après la virgule.}
\end{attention}

\begin{attention}[Le mot-outil est obligatoire après la virgule]
La phrase \zh{除了汉语以外，我学英语。} est \textbf{incorrecte} (incomplète). Il faut obligatoirement \zh{还}, \zh{也}, \zh{都} ou \zh{只}.\\
Correct (addition) : \zh{除了汉语以外，我\bfseries 还} 学英语。\\
Correct (exclusion) : \zh{除了汉语以外，我\bfseries 什么都不} 学。 / \zh{除了汉语以外，我\bfseries 只} 学英语。
\end{attention}

\begin{methode}[Comment choisir le bon mot-outil : 还/也 ou 都/只 ?]
\begin{enumerate}
\item Repérez X (l'élément après \zh{除了}).
\item Posez la question : \textbf{X fait-il partie de l'ensemble décrit par le prédicat ?}
\item \textbf{OUI} (X est compris) $\rightarrow$ \textbf{Addition} : utilisez \zh{也} ou \zh{还} (« en plus de X, aussi… »).
\item \textbf{NON} (X est mis à part) $\rightarrow$ \textbf{Exclusion} : utilisez \zh{都} (« tous… sauf X ») ou \zh{只} (« seulement… »).
\item \textbf{Vérification} : \zh{都} appelle souvent un sujet pluriel (\zh{大家}, \zh{我们}) ou un adverbe d'étendue (\zh{每天}, \zh{什么}) ; \zh{只} précède le verbe ou l'objet unique.
\end{enumerate}
\end{methode}

\begin{exoresolu}[Traduire en chinois (Thème)]
Énoncé : Traduisez en chinois : a) « En plus du yuan, elle accepte aussi l'euro. » b) « Sauf le samedi, la boutique est ouverte tous les jours. » c) « Moi, à part le français, je ne parle que l'anglais. »
\tcblower
Solution détaillée :
a) Le yuan est \emph{compris} dans les monnaies acceptées $\rightarrow$ Addition, mot-outil \zh{也} ou \zh{还} :\\
\zh{除了人民币以外，她也收欧元。} \emph{Chúle rénmínbì yǐwài, tā yě shōu ōuyuán.}\\[4pt]
b) Le samedi est \emph{exclu} des jours d'ouverture $\rightarrow$ Exclusion, mot-outil \zh{都} (avec \zh{每天}) :\\
\zh{除了星期六以外，商店每天都开门。} \emph{Chúle xīngqīliù yǐwài, shāngdiàn měitiān dōu kāimén.}\\[4pt]
c) Le français est \emph{exclu} des langues parlées (je ne parle que l'anglais) $\rightarrow$ Exclusion, mot-outil \zh{只}, sujet topicalisé :\\
\zh{我除了法语以外，只会说英语。} \emph{Wǒ chúle Fǎyǔ yǐwài, zhǐ huì shuō Yīngyǔ.}
\end{exoresolu}

\begin{exercice}[Entraînement : choisir le mot-outil]
Complétez chaque phrase avec \zh{还}, \zh{也}, \zh{都} ou \zh{只} (une seule réponse possible par phrase selon le sens logique).
\begin{enumerate}
\item 除了汉语以外，我们班的同学 \underline{\hspace{2cm}} 学英语。 (Tous les élèves apprennent l'anglais \emph{en plus} du chinois.)
\item 除了妈妈以外，家里的人 \underline{\hspace{2cm}} 来了。 (Maman n'est pas venue, les autres si.)
\item 除了茶以外，他 \underline{\hspace{2cm}} 喝咖啡，不喝别的。 (Il ne boit que du café, pas de thé, pas d'autre chose.)
\item 除了数学以外，我 \underline{\hspace{2cm}} 喜欢汉语和历史。 (Les maths sont ma matière principale, le chinois et l'histoire viennent \emph{s'ajouter}.)
\item 除了足球以外，他 \underline{\hspace{2cm}} 什么运动 \underline{\hspace{2cm}} 不喜欢。 (Double trou : il n'aime \emph{aucun} sport sauf le foot.)
\end{enumerate}
\end{exercice}

\begin{corrige}[Exercice : choisir le mot-outil]
\begin{enumerate}
\item \zh{也} ou \zh{还} (Addition : le chinois est inclus dans le cursus). Phrase : \zh{除了汉语以外，我们班的同学\bfseries 也} 学英语。
\item \zh{都} (Exclusion : maman est la seule absente). Phrase : \zh{除了妈妈以外，家里的人\bfseries 都} 来了。
\item \zh{只} (Exclusion restrictive : le café est l'unique autre boisson). Phrase : \zh{除了茶以外，他\bfseries 只} 喝咖啡，不喝别的。
\item \zh{也} ou \zh{还} (Addition : j'aime les maths \emph{et aussi} le chinois/l'histoire). Phrase : \zh{除了数学以外，我\bfseries 还} 喜欢汉语和历史。
\item \zh{都} … \zh{不} (Exclusion totale : « aucun sport » = « tous les sports ne sont pas aimés »). Phrase : \zh{除了足球以外，他\bfseries 什么运动都不} 喜欢。
\end{enumerate}
\end{corrige}

\begin{savaistu}[La famille des mots de localisation en \zh{以} (yǐ)]
\zh{以外} (yǐwài, « en dehors de ») appartient à une famille productive de noms de localisation formés avec \zh{以} (yǐ, « à partir de », « selon ») + un mot de position. Ils servent à délimiter une frontière dans l'espace, le temps ou la quantité :
\begin{center}
\begin{tabularx}{\linewidth}{|X|X|X|}
\hline
\rowcolor{poppurpleL} \textbf{Mot} & \textbf{Sens littéral} & \textbf{Exemple d'usage} \\
\hline
\zh{以内} (yǐnèi) & à l'intérieur de & \zh{三天以内} (sān tiān yǐnèi) : dans les trois jours / en moins de trois jours \\
\hline
\zh{以前} (yǐqián) & devant (dans le temps) & \zh{八点以前} (bā diǎn yǐqián) : avant huit heures \\
\hline
\zh{以后} (yǐhòu) & derrière (dans le temps) & \zh{下课以后} (xiàkè yǐhòu) : après la classe \\
\hline
\zh{以上} (yǐshàng) & au-dessus de & \zh{一百元以上} (yībǎi yuán yǐshàng) : cent yuans ou plus / au-dessus de 100 \\
\hline
\zh{以下} (yǐxià) & en dessous de & \zh{十八岁以下} (shíbā suì yǐxià) : moins de 18 ans / mineur \\
\hline
\end{tabularx}
\end{center}
\cle{Mécanisme} : \zh{以} + [Point de repère] = Limite. Retenez ce schéma pour construire vous-même des expressions de durée, quantité ou délai.
\end{savaistu}

\coursec{Comparaison avec le français et synthèse}

\courssub{Transition : du calcul à la comparaison}

Dans les sections précédentes, nous avons appris à poser une addition avec \zh{加} (jiā), une soustraction avec \zh{减} (jiǎn), à dire un prix avec les monnaies chinoise et camerounaise, et à exprimer l'addition ou l'exclusion avec le cadre \zh{除了……以外，……} (chúle……yǐwài, ……). Il est temps de prendre du recul : comment le chinois dit-il ce que le français dit autrement ? Cette dernière section compare systématiquement les deux langues, récapitule les quatre savoirs de la leçon, puis te propose une méthode de vérification et une auto-évaluation. C'est le moment de consolider avant l'évaluation.

\courssub{Comparer les deux langues : le calcul}

\begin{exemplebox}[Observer avant de comparer]
Observons deux calculs, d'abord en français, puis en chinois :\\[4pt]
\zh{三加二等于五。} \emph{Sān jiā èr děngyú wǔ.} « Trois plus deux égale cinq. » (3 + 2 = 5)\\[4pt]
\zh{十减四等于六。} \emph{Shí jiǎn sì děngyú liù.} « Dix moins quatre égale six. » (10 − 4 = 6)
\end{exemplebox}

Que remarque-t-on ? L'\cle{ordre des mots} est exactement le même dans les deux langues : premier nombre, opération, deuxième nombre, résultat. C'est une excellente nouvelle pour les francophones : aucun déplacement à effectuer. Mais les \cle{mots-outils} sont de nature différente.

\begin{propriete}[Le chinois calcule avec des verbes]
En français, « plus », « moins » et « égal » sont des mots invariables, presque des symboles lus à voix haute. En chinois, ce sont de vrais \cle{verbes} :
\begin{itemize}
\item \zh{加} (jiā) : verbe « ajouter » → \zh{三加二} littéralement « trois ajoute deux » ;
\item \zh{减} (jiǎn) : verbe « retrancher » → \zh{十减四} littéralement « dix retranche quatre » ;
\item \zh{等于} (děngyú) : verbe « être égal à » → il est \textbf{obligatoire} et ne peut jamais être omis.
\end{itemize}
Structure générale : \textbf{N1 + 加/减 + N2 + 等于 + N3}.
\end{propriete}

\begin{propriete}[等于 est un verbe : on peut le nier]
Comme tout verbe, \zh{等于} se nie avec \zh{不} : \zh{不等于} (bù děngyú) signifie « n'est pas égal à », ce qui correspond au symbole français « ≠ ».\\[4pt]
\zh{二加二不等于五。} \emph{Èr jiā èr bù děngyú wǔ.} « Deux plus deux n'est pas égal à cinq. » (2 + 2 ≠ 5)\\[4pt]
\zh{一百美元不等于一百人民币。} \emph{Yìbǎi Měiyuán bù děngyú yìbǎi rénmínbì.} « Cent dollars ne valent pas cent yuans. »
\end{propriete}

\courssub{Comparer les deux langues : sauf, à part, en plus de}

Le français possède trois expressions distinctes : « sauf » (exclusion), « à part » (exclusion ou addition), « en plus de » (addition). Le chinois n'utilise qu'\cle{un seul cadre} : \zh{除了……以外，……}. C'est le \cle{mot-outil de la deuxième partie de la phrase} qui décide du sens :

\begin{center}
\begin{tabularx}{\linewidth}{|X|X|X|}
\hline
\rowcolor{popblueL} Structure & Exemple en caractères + pinyin & Traduction \\
\hline
除了 A 以外，B 都…… & 除了我以外，大家都喜欢足球。Chúle wǒ yǐwài, dàjiā dōu xǐhuan zúqiú. & Sauf moi, tout le monde aime le football. (exclusion, avec 都) \\
\hline
除了 A 以外，B 还…… & 除了英语以外，我还学汉语。Chúle Yīngyǔ yǐwài, wǒ hái xué Hànyǔ. & En plus de l'anglais, j'apprends aussi le chinois. (addition, avec 还) \\
\hline
除了 A 以外，B 也…… & 除了人民币以外，我也有美元。Chúle rénmínbì yǐwài, wǒ yě yǒu Měiyuán. & À part les yuans, j'ai aussi des dollars. (addition, avec 也) \\
\hline
\end{tabularx}
\end{center}

Règle d'or : \zh{都} (dōu) marque l'\cle{exclusion} (« sauf »), \zh{还} (hái) et \zh{也} (yě) marquent l'\cle{addition} (« en plus de », « à part »). Sans mot-outil dans la seconde partie, la phrase est fautive ou ambiguë. Nuance utile : \zh{还} introduit souvent une information \emph{nouvelle} qui s'ajoute, tandis que \zh{也} signale simplement une similitude (« aussi »).

\begin{popfigure}
\begin{tikzpicture}[font=\small]
\draw[thick, popline] (0,0) rectangle (7.2,4.6);
\node[fill=popblue, text=white, font=\bfseries, minimum width=6.8cm] at (3.6,4.2) {Français};
\node[align=left] at (3.6,3.5) {3 + 2 = 5 : trois \textbf{plus} deux \textbf{égale} cinq};
\node[align=left] at (3.6,2.9) {10 − 4 = 6 : dix \textbf{moins} quatre \textbf{égale} six};
\node[align=left] at (3.6,2.3) {\textbf{sauf} moi, tout le monde…};
\node[align=left] at (3.6,1.7) {\textbf{en plus de} l'anglais, j'apprends aussi…};
\node[align=left] at (3.6,1.1) {cinquante mille : 50 000};
\node[align=left] at (3.6,0.5) {100 dollars \textbf{valent} 700 yuans};
\draw[thick, popline] (7.4,0) rectangle (14.6,4.6);
\node[fill=poporange, text=white, font=\bfseries, minimum width=6.8cm] at (11,4.2) {中文 Zhōngwén};
\node[align=left] at (11,3.5) {三加二等于五 (verbe 加 + verbe 等于)};
\node[align=left] at (11,2.9) {十减四等于六 (verbe 减 + verbe 等于)};
\node[align=left] at (11,2.3) {除了我以外，大家都… (都 = exclusion)};
\node[align=left] at (11,1.7) {除了英语以外，我还学… (还 = addition)};
\node[align=left] at (11,1.1) {五万 (5 × 万, pas « cinquante mille »)};
\node[align=left] at (11,0.5) {一百美元等于七百多人民币 (多 = « plus de »)};
\end{tikzpicture}
\legende{Tableau comparatif Français / 中文 : six équivalences à mémoriser}
\end{popfigure}

\courssub{Les quatre erreurs-types des francophones}

\begin{attention}[Récapitulatif des erreurs à ne plus commettre]
\begin{enumerate}
\item \zh{三加五八。} ✗ → \zh{三加五等于八。} ✓ : le verbe \zh{等于} est \textbf{obligatoire} ; on ne peut pas juxtaposer le calcul et le résultat comme en français parlé (« trois et cinq, huit »).
\item \zh{二加三做五。} ✗ → \zh{二加三等于五。} ✓ : le français dit « deux plus trois \emph{font} cinq », mais « faire » ne se traduit pas par \zh{做} (zuò) ici ; il faut le verbe \zh{等于}.
\item \zh{除了英语，我学汉语。} ✗ → \zh{除了英语以外，我还学汉语。} ✓ : après le cadre \zh{除了……以外}, la seconde partie exige un mot-outil (\zh{还}, \zh{也} ou \zh{都}) ; « sauf » ne se traduit jamais par une préposition seule.
\item \zh{五十千} ✗ → \zh{五万} ✓ : le chinois compte par \zh{万} (wàn, dix mille), pas par milliers ; 50 000 se dit « cinq dix-mille », \zh{五万} (wǔ wàn). Attention aux prix en francs CFA : 50 000 FCFA = \zh{五万非郎}.
\end{enumerate}
\end{attention}

\courssub{Synthèse des quatre savoirs de la leçon}

\begin{aretenir}[Bilan de la leçon 26]
\begin{enumerate}
\item \textbf{L'addition 加法} : N1 + \zh{加} + N2 + \zh{等于} + N3. Ex. : \zh{七加八等于十五。} (Qī jiā bā děngyú shíwǔ.) « 7 + 8 = 15. »
\item \textbf{La soustraction 减法} : N1 + \zh{减} + N2 + \zh{等于} + N3. Ex. : \zh{二十减五等于十五。} (Èrshí jiǎn wǔ děngyú shíwǔ.) « 20 − 5 = 15. »
\item \textbf{Les monnaies et les nombres} : \zh{块} (kuài) pour les yuans, \zh{美元} (Měiyuán) pour les dollars, \zh{非郎} (fēiláng) pour les francs CFA ; grands nombres avec \zh{万} ; \zh{多} (duō) après un nombre rond = « plus de » : \zh{七百多} « plus de 700 ».
\item \textbf{除了……以外，……} : avec \zh{都} = exclusion (« sauf »), avec \zh{还}/\zh{也} = addition (« en plus de »).
\end{enumerate}
\end{aretenir}

\begin{exemplebox}[Phrase-bilan de la leçon]
\zh{除了人民币以外，我还有美元，一百美元等于七百多人民币。}\\[2pt]
\emph{Chúle rénmínbì yǐwài, wǒ hái yǒu Měiyuán, yìbǎi Měiyuán děngyú qībǎi duō rénmínbì.}\\[2pt]
« En plus des yuans, j'ai aussi des dollars ; cent dollars valent plus de sept cents yuans. »\\[4pt]
Cette phrase réunit trois savoirs de la leçon : le cadre \zh{除了……以外} avec \zh{还} (addition), le vocabulaire des monnaies (\zh{人民币}, \zh{美元}) et le verbe \zh{等于} avec \zh{多} pour « plus de ».
\end{exemplebox}

\begin{popfigure}
\begin{tikzpicture}[mindmap, grow cyclic, every node/.style={concept, font=\small, align=center}, concept color=popblue!60, level 1/.append style={level distance=3.4cm, sibling angle=90}]
\node[concept color=poporange!80, font=\bfseries, align=center]{第26课\\ 加法和减法}
  child[concept color=popblue!50] { node[align=center]{加法\\ N1+加+N2\\ +等于+N3} }
  child[concept color=popgreen!50] { node[align=center]{减法\\ N1+减+N2\\ +等于+N3} }
  child[concept color=popgold!60] { node[align=center]{钱和数字\\ 块、美元、非郎\\ 万、多} }
  child[concept color=poppink!60] { node[align=center]{除了……以外\\ 都 = sauf\\ 还/也 = en plus de} };
\end{tikzpicture}
\legende{Carte mentale finale : les quatre branches de la leçon 26}
\end{popfigure}

\courssub{Méthode : vérifier sa phrase en trois questions}

\begin{methode}[Le réflexe du bon élève avant de rendre sa copie]
Avant de valider une phrase de calcul, de prix ou d'exclusion, pose-toi trois questions :
\begin{enumerate}
\item \textbf{Ordre du calcul} : ma phrase suit-elle bien le schéma N1 + \zh{加}/\zh{减} + N2 + \zh{等于} + N3 ? Le verbe \zh{等于} est-il présent (ou sa négation \zh{不等于}) ?
\item \textbf{Argent} : ai-je mis le bon mot de mesure (\zh{块} pour les yuans) et le bon nom de monnaie (\zh{人民币}, \zh{美元}, \zh{非郎}) ? Mes grands nombres passent-ils par \zh{万} et non par « mille » ?
\item \textbf{Mot-outil} : après \zh{除了……以外}, ma seconde partie contient-elle \zh{都} (exclusion) ou \zh{还}/\zh{也} (addition) ?
\end{enumerate}
Si les trois réponses sont « oui », la phrase est correcte.
\end{methode}

\begin{exoresolu}[Exercice de synthèse : corriger et traduire]
Énoncé. Voici quatre phrases d'élèves de Terminale A à Yaoundé. Trois contiennent une erreur ; une seule est correcte. Trouve la phrase correcte, corrige les trois autres, puis traduis les quatre phrases.\\
a) \zh{五加六做十一。} b) \zh{一百减四十等于六十。} c) \zh{除了汉语，我学英语。} d) 50 000 FCFA : \zh{五十千非郎}.
\tcblower
Solution détaillée.\\
a) Faute : \zh{做} ne peut pas exprimer « font » dans un calcul. Correction : \zh{五加六等于十一。} \emph{Wǔ jiā liù děngyú shíyī.} « 5 + 6 = 11. »\\
b) Phrase correcte : ordre N1 + \zh{减} + N2 + \zh{等于} + N3 respecté. \emph{Yìbǎi jiǎn sìshí děngyú liùshí.} « 100 − 40 = 60. »\\
c) Faute : il manque le mot-outil après \zh{除了……以外} (et \zh{以外} est omis). Correction : \zh{除了汉语以外，我还学英语。} \emph{Chúle Hànyǔ yǐwài, wǒ hái xué Yīngyǔ.} « En plus du chinois, j'apprends aussi l'anglais. »\\
d) Faute : le chinois ne compte pas par milliers. Correction : \zh{五万非郎} \emph{wǔ wàn fēiláng} « 50 000 francs CFA ».
\end{exoresolu}

\begin{exercice}[À toi de jouer]
1) Traduis en chinois : « 8 + 9 = 17 » ; « 50 − 20 = 30 » ; « 3 + 4 ≠ 8 ».\\
2) Dis en chinois : 30 000 FCFA ; 200 dollars ; « plus de 500 yuans ».\\
3) Complète avec \zh{都}, \zh{还} ou \zh{也} : \zh{除了足球以外，我……喜欢篮球。} (j'aime aussi le basket) ; \zh{除了王明以外，我们……去市场了。} (sauf Wang Ming).\\
4) Rédige une phrase-bilan personnelle sur le modèle de la phrase-bilan de la leçon, avec une monnaie et un calcul.
\end{exercice}

\begin{corrige}[Exercice « À toi de jouer »]
1) \zh{八加九等于十七。} (Bā jiā jiǔ děngyú shíqī.) ; \zh{五十减二十等于三十。} (Wǔshí jiǎn èrshí děngyú sānshí.) ; \zh{三加四不等于八。} (Sān jiā sì bù děngyú bā.)\\
2) \zh{三万非郎} (sān wàn fēiláng) ; \zh{二百美元} (èrbǎi Měiyuán) ; \zh{五百多块} (wǔbǎi duō kuài).\\
3) \zh{还} (addition : « j'aime aussi ») ; \zh{都} (exclusion : « tous sauf lui »).\\
4) Modèle possible : \zh{除了非郎以外，我还有人民币，一百块人民币等于八千多非郎。} « En plus des francs CFA, j'ai aussi des yuans ; cent yuans valent plus de huit mille francs CFA. »
\end{corrige}

\courssub{Auto-évaluation de fin de leçon}

Coche mentalement chaque compétence que tu maîtrises. Si une case reste vide, retourne à la section correspondante.

\begin{center}
\begin{tabularx}{\linewidth}{|X|X|}
\hline
\rowcolor{popblueL} Je sais… & Exemple que je peux produire \\
\hline
poser une addition en chinois & 三加二等于五。 \\
\hline
poser une soustraction et la nier & 十减四等于六。 / 二加二不等于五。 \\
\hline
dire un prix et rendre la monnaie & 一共三十五块。找你十五块。 \\
\hline
nommer les monnaies et les grands nombres & 人民币、美元、非郎 ; 五万 \\
\hline
exprimer l'exclusion avec 都 & 除了我以外，大家都喜欢足球。 \\
\hline
exprimer l'addition avec 还/也 & 除了英语以外，我还学汉语。 \\
\hline
\end{tabularx}
\end{center}

\begin{savaistu}[Compter en Chine et au Cameroun]
En Chine, on négocie les prix au marché exactement comme à Mokolo ou à Sandaga : le calcul mental rapide est une compétence sociale admirée. Les Chinois utilisent souvent les gestes de la main pour montrer les chiffres de un à dix d'une seule main, très utile dans un marché bruyant. Autre point commun : comme le franc CFA, le yuan possède un nom officiel (\zh{人民币}, « monnaie du peuple ») et un nom courant (\zh{块}, « morceau »), tout comme on dit « mille francs » pour « mille francs CFA » au quotidien.
\end{savaistu}

Tu es maintenant capable de calculer, d'acheter, de rendre la monnaie et d'exprimer l'exclusion ou l'addition en chinois : quatre compétences concrètes du monde économique, au cœur du module 4. La prochaine leçon te permettra de réutiliser ces structures dans des situations de travail et de commerce plus complexes.

\newpage

\coursec{Activité d'intégration}

\courssub{Objectifs de l'activité}

À l'issue de cette activité d'intégration, tu seras capable de :
\begin{itemize}
\item mobiliser dans une situation réelle de communication les structures de calcul étudiées : l'addition avec \zh{加} (jiā), la soustraction avec \zh{减} (jiǎn) et l'égalité avec \zh{等于} (děngyú) ;
\item utiliser correctement les nombres et les mots de mesure monétaire (\zh{块}, \zh{毛}, \zh{分}) ainsi que le lexique des monnaies (\zh{法郎}, \zh{人民币}, \zh{美元}) ;
\item employer la structure \zh{除了……以外，……还……} pour exprimer l'addition d'un élément à un ensemble, et la distinguer de \zh{除了……以外，……都……} qui exprime l'exclusion ;
\item comprendre et produire oralement des échanges commerciaux authentiques (demander un prix, négocier, payer, rendre la monnaie) ;
\item rédiger un court texte écrit fonctionnel : un reçu de transaction en quatre phrases chinoises.
\end{itemize}

\courssub{Situation}

C'est samedi matin au marché Mokolo de Yaoundé. Ton groupe anime un stand de fruits et de tissus. Deux clients arrivent : Amina, une cliente camerounaise qui paie en francs CFA, et M. Wang, un touriste chinois en visite au Cameroun, qui a dans son portefeuille des yuans et des dollars américains. Le vendeur doit annoncer les prix, les clients doivent calculer le total de leurs achats à voix haute en chinois, vérifier la monnaie rendue et expliquer avec quelles monnaies ils peuvent payer.

Voici l'affiche des prix du stand :

\begin{textebox}[L'affiche du stand — 摊位的价格表]
今日价格 (Jīnrì jiàgé) — Prix du jour\\
菠萝：五百法郎一个。(Bōluó: wǔ bǎi fǎláng yí ge.) — Ananas : 500 FCFA l'unité.\\
芒果：三百法郎一公斤。(Mángguǒ: sān bǎi fǎláng yì gōngjīn.) — Mangues : 300 FCFA le kilo.\\
香蕉：两百法郎一串。(Xiāngjiāo: liǎng bǎi fǎláng yí chuàn.) — Bananes : 200 FCFA le régime.\\
布料：五千法郎一块。(Bùliào: wǔ qiān fǎláng yí kuài.) — Tissu : 5 000 FCFA la pièce.\\
除了法郎以外，我们也收美元和人民币。(Chúle fǎláng yǐwài, wǒmen yě shōu Měiyuán hé Rénmínbì.) — Outre les francs CFA, nous acceptons aussi les dollars et les yuans.
\end{textebox}

\begin{savaistu}[Les monnaies au quotidien]
Au Cameroun, on paie en francs CFA (\zh{法郎}, fǎláng, littéralement « franc »). En Chine, la monnaie est le yuan, dont le nom officiel est \zh{人民币} (Rénmínbì, « monnaie du peuple ») : à l'écrit on compte en \zh{元} (yuán), mais à l'oral on dit \zh{块} (kuài), avec \zh{毛} (máo) pour les dixièmes et \zh{分} (fēn) pour les centièmes — par exemple \zh{三块五} (sān kuài wǔ), « 3,50 yuans ». Le dollar américain se dit \zh{美元} (Měiyuán, « monnaie américaine »). Dans les marchés internationaux et les zones touristiques, il n'est pas rare qu'un vendeur accepte plusieurs monnaies : c'est exactement la situation où la structure \zh{除了……以外，……还……} devient indispensable.
\end{savaistu}

\courssub{Consignes}

Travail en groupes de trois : un vendeur, deux clients (un client camerounais, un client chinois).

\begin{enumerate}
\item \textbf{Préparation du stand (compréhension et lexique).} Lisez l'affiche des prix à voix haute. Soulignez les nombres et les classificateurs (\zh{个}, \zh{公斤}, \zh{串}, \zh{块}). Chaque client choisit au moins deux articles à acheter.
\item \textbf{Demander les prix (oral).} Les clients interrogent le vendeur avec la question \zh{多少钱？} (Duōshao qián ?) ou \zh{这个怎么卖？} (Zhège zěnme mài ?). Le vendeur répond avec une phrase complète : \zh{菠萝五百法郎一个。} (Bōluó wǔ bǎi fǎláng yí ge.) — « L'ananas est à 500 FCFA l'unité. »
\item \textbf{Calculer le total (grammaire : l'addition).} Chaque client calcule son total à voix haute, en chinois, avec la structure \zh{A 加 B 等于 C}. Exemple : \zh{五百加三百等于八百。} (Wǔ bǎi jiā sān bǎi děngyú bā bǎi.) — « 500 + 300 = 800. »
\item \textbf{Payer et rendre la monnaie (grammaire : la soustraction).} Le client donne un billet plus grand que le total ; le vendeur calcule la monnaie à rendre avec \zh{A 减 B 等于 C}. Exemple : \zh{一千减八百等于两百。} (Yì qiān jiǎn bā bǎi děngyú liǎng bǎi.) — « 1 000 − 800 = 200. »
\item \textbf{Dire avec quelles monnaies on paie (grammaire : 除了……以外).} Chaque client précise les monnaies dont il dispose avec la structure \zh{除了……以外，……还……}. Exemple : \zh{除了法郎以外，我还有美元。} (Chúle fǎláng yǐwài, wǒ hái yǒu Měiyuán.) — « Outre les francs CFA, j'ai aussi des dollars. »
\item \textbf{Représentation devant la classe.} Chaque groupe joue sa scène. Les autres groupes remplissent une fiche d'écoute : articles achetés, total annoncé, monnaie rendue, monnaies utilisées. Ils vérifient les calculs et signalent les erreurs éventuelles.
\item \textbf{Production écrite de clôture.} Individuellement, rédigez le reçu de la transaction en \textbf{quatre phrases chinoises} : (1) les articles achetés, (2) le calcul du total avec \zh{加}, (3) le calcul de la monnaie rendue avec \zh{减}, (4) une phrase avec \zh{除了……以外，……还……} sur les monnaies acceptées.
\end{enumerate}

\begin{attention}[Pièges à éviter pendant le jeu de rôle]
\begin{itemize}
\item Ne dis jamais \zh{等于了} : \zh{等于} est un verbe d'état, il ne prend pas \zh{了}.
\item Dans \zh{除了……以外}, le mot \zh{以外} ne s'omet pas à l'écrit ; à l'oral on peut dire simplement \zh{除了……}, mais dans ton reçu écrit, écris la forme complète.
\item Attention à l'ordre : on dit \zh{五百加三百} (500 + 300), jamais l'inverse si le résultat attendu suit cet ordre ; et pour la soustraction, le plus grand nombre vient d'abord : \zh{一千减八百}, pas \zh{八百减一千}.
\item Deux cents se dit \zh{两百} (liǎng bǎi) avec \zh{两}, et non \zh{二百} dans la langue courante du marché.
\item Ne confonds pas les deux sens de \zh{除了……以外} : avec \zh{还}/\zh{也}, on \emph{ajoute} un élément (« outre…, aussi… ») ; avec \zh{都}, on \emph{exclut} un élément (« sauf…, tous… »). Compare : \zh{除了法郎以外，我们还收美元。} (« Outre les francs, nous acceptons aussi les dollars. ») et \zh{除了法郎以外，我们什么都不收。} (« Sauf les francs, nous n'acceptons rien. »).
\end{itemize}
\end{attention}

\courssub{Ressources à mobiliser}

\begin{propriete}[Rappel des structures de calcul]
\begin{itemize}
\item \textbf{Addition} : A + \zh{加} + B + \zh{等于} + C. \zh{三加五等于八。} (Sān jiā wǔ děngyú bā.) — « 3 + 5 = 8. »
\item \textbf{Soustraction} : A + \zh{减} + B + \zh{等于} + C. \zh{十减四等于六。} (Shí jiǎn sì děngyú liù.) — « 10 − 4 = 6. »
\item \textbf{Addition d'un élément (inclusion)} : \zh{除了} + X + \zh{以外}，sujet + \zh{还}/\zh{也} + verbe. \zh{除了法郎以外，我们还收美元。} (Chúle fǎláng yǐwài, wǒmen hái shōu Měiyuán.) — « Outre les francs, nous acceptons aussi les dollars. »
\item \textbf{Exclusion d'un élément} : \zh{除了} + X + \zh{以外}，sujet + \zh{都} + verbe (souvent à la négative). \zh{除了王同学以外，大家都来了。} (Chúle Wáng tóngxué yǐwài, dàjiā dōu lái le.) — « Sauf le camarade Wang, tout le monde est venu. »
\end{itemize}
\end{propriete}

\begin{center}
\begin{tabularx}{\linewidth}{|X|X|X|X|}
\hline
\rowcolor{popblueL} Caractères & Pinyin & Nature & Traduction \\
\hline
加法 & jiāfǎ & n. & addition \\
\hline
减法 & jiǎnfǎ & n. & soustraction \\
\hline
等于 & děngyú & v. & être égal à \\
\hline
法郎 & fǎláng & n. & franc (CFA) \\
\hline
人民币 & Rénmínbì & n. & yuan (monnaie chinoise) \\
\hline
美元 & Měiyuán & n. & dollar américain \\
\hline
价格 & jiàgé & n. & prix \\
\hline
找钱 & zhǎo qián & v. & rendre la monnaie \\
\hline
一共 & yígòng & adv. & au total \\
\hline
收 & shōu & v. & accepter, recevoir (argent) \\
\hline
\end{tabularx}
\end{center}

\courssub{Proposition de production et corrigé commenté}

\begin{exoresolu}[Le reçu de la transaction — modèle de production]
Rédige le reçu de l'achat d'Amina (un ananas à 500 FCFA et un kilo de mangues à 300 FCFA ; elle paie avec un billet de 1 000 FCFA) en quatre phrases chinoises.
\tcblower
\textbf{Reçu modèle :}\\
\zh{阿米娜买了一个菠萝和一公斤芒果。} (Āmínà mǎi le yí ge bōluó hé yì gōngjīn mángguǒ.) — « Amina a acheté un ananas et un kilo de mangues. »\\
\zh{五百加三百等于八百，一共八百法郎。} (Wǔ bǎi jiā sān bǎi děngyú bā bǎi, yígòng bā bǎi fǎláng.) — « 500 + 300 = 800, soit 800 FCFA au total. »\\
\zh{她给了一千法郎，一千减八百等于两百，我找给她两百法郎。} (Tā gěi le yì qiān fǎláng, yì qiān jiǎn bā bǎi děngyú liǎng bǎi, wǒ zhǎo gěi tā liǎng bǎi fǎláng.) — « Elle a donné 1 000 FCFA ; 1 000 − 800 = 200, je lui ai rendu 200 FCFA. »\\
\zh{除了法郎以外，我们还收美元和人民币。} (Chúle fǎláng yǐwài, wǒmen hái shōu Měiyuán hé Rénmínbì.) — « Outre les francs CFA, nous acceptons aussi les dollars et les yuans. »

\textbf{Commentaire des choix de langue :}
\begin{itemize}
\item Phrase 1 : le classificateur \zh{个} accompagne \zh{菠萝} et \zh{公斤} sert de mot de mesure pour les mangues ; \zh{了} marque l'action accomplie de l'achat. Le prénom Amina se transcrit \zh{阿米娜} (Āmínà).
\item Phrase 2 : l'addition suit strictement le schéma A + \zh{加} + B + \zh{等于} + C ; l'adverbe \zh{一共} (« au total ») résume naturellement le calcul, comme sur un vrai reçu.
\item Phrase 3 : la soustraction respecte l'ordre grand nombre − petit nombre ; \zh{找给她} illustre le verbe \zh{找} (« rendre la monnaie ») suivi de \zh{给} qui introduit le bénéficiaire.
\item Phrase 4 : la structure \zh{除了……以外} est complète, avec \zh{还} qui introduit l'élément ajouté ; les noms de monnaies étrangères prennent la majuscule en pinyin (\zh{Měiyuán}, \zh{Rénmínbì}) car ce sont des noms propres.
\end{itemize}
\end{exoresolu}

\courssub{Bilan de l'activité}

Cette activité t'a permis de réinvestir l'ensemble des savoir-faire de la leçon dans une situation de communication authentique et ancrée dans la réalité camerounaise : le marché. Tu as compris un document réel (une affiche de prix), utilisé le lexique des monnaies et des nombres, appliqué oralement les structures de calcul \zh{加} et \zh{减} avec \zh{等于}, exprimé l'addition d'un élément avec \zh{除了……以外，……还……} en le distinguant de l'exclusion avec \zh{都}, puis consolidé le tout par une production écrite fonctionnelle. La vérification croisée des calculs par les autres groupes a en outre développé ta compréhension orale et ton esprit critique. Ces compétences te seront directement utiles pour toute situation d'achat, de change ou de négociation impliquant le chinois, au Cameroun comme en Chine.

\newpage

\coursec{Méthodes et exercices résolus}

\coursec{Leçon 26 — Grammaire : L'addition et la soustraction ; les monnaies ; 除了……以外……}

\begin{textebox}[Observation : Calculer au marché de Mokolo]
\zh{客户：老板，西红柿多少钱一斤？} \\
\zh{老板：五块钱一斤。} \\
\zh{客户：给我两斤。再来一斤苦菜。} \\
\zh{老板：西红柿两斤十块，苦菜一斤五块。一共十五块。} \\
\zh{客户：二十块给你，找我五块。} \\
\zh{老板：好，找您五块。谢谢！} \\

\textbf{Pinyin :} \\
Kèhù: Lǎobǎn, xīhóngshì duōshao qián yì jīn? \\
Lǎobǎn: Wǔ kuài qián yì jīn. \\
Kèhù: Gěi wǒ liǎng jīn. Zài lái yì jīn kǔcài. \\
Lǎobǎn: Xīhóngshì liǎng jīn shí kuài, kǔcài yì jīn wǔ kuài. Yígòng shíwǔ kuài. \\
Kèhù: Èrshí kuài gěi nǐ, zhǎo wǒ wǔ kuài. \\
Lǎobǎn: Hǎo, zhǎo nín wǔ kuài. Xièxie! \\

\textbf{Traduction :} \\
Client : Patron, combien le demi-kilo de tomates ? \\
Vendeur : 5 yuans le demi-kilo. \\
Client : Donnez-m'en deux. Et un demi-kilo de ndolé (苦菜). \\
Vendeur : Deux demi-kilos de tomates = 10 yuans, un demi-kilo de ndolé = 5 yuans. Total 15 yuans. \\
Client : Voici 20 yuans, rendez-moi 5 yuans. \\
Vendeur : D'accord, je vous rends 5 yuans. Merci !
\end{textebox}

\courssub{L'addition : 加法}

\begin{methode}[Dire et écrire une addition : 加法]
Pour poser ou résoudre une addition en chinois, suis ces étapes :
\begin{enumerate}
\item Identifie les deux nombres (chiffres ou caractères) : \zh{三} (sān, 3), \zh{五} (wǔ, 5).
\item Utilise la structure fixe : \cle{A 加 B 等于 C} (A jiā B děngyú C), littéralement « A plus B égale C ».
\item Le verbe \cle{加} (jiā, ajouter/plus) se place \textbf{toujours entre} les deux nombres. Jamais « 加三加五 ».
\item Pour poser la question, remplace le résultat par \cle{多少} (duōshao) : \zh{A 加 B 等于多少？}
\item Au marché, pour demander le total : \cle{一共多少钱？} (Yígòng duōshao qián ?). \cle{一共} (yígòng, en tout) se place avant le verbe ou le nombre.
\item Réponse courte orale : \zh{是 + nombre} (Shì + nombre, « C'est + nombre »).
\end{enumerate}
\end{methode}

\begin{exemplebox}[Exemples d'additions]
\begin{center}
\begin{tabularx}{\linewidth}{|X|X|X|}
\hline
\rowcolor{popblueL} \textbf{Chinois (Caractères + Pinyin)} & \textbf{Structure} & \textbf{Traduction} \\
\hline
\zh{三加五等于八。} (Sān jiā wǔ děngyú bā.) & 3 + 5 = 8 & Trois plus cinq font huit. \\
\zh{一百加二百等于三百。} (Yìbǎi jiā èrbǎi děngyú sānbǎi.) & 100 + 200 = 300 & Cent plus deux cents font trois cents. \\
\zh{五十加三十等于八十。} (Wǔshí jiā sānshí děngyú bāshí.) & 50 + 30 = 80 & Cinquante plus trente font quatre-vingts. \\
\zh{这个加那个一共多少钱？} (Zhège jiā nàge yígòng duōshao qián?) & Sujet + 加 + Sujet + 一共 & Celui-ci plus celui-là, combien en tout ? \\
\hline
\end{tabularx}
\end{center}
\end{exemplebox}

\begin{exoresolu}[Au marché de Mokolo : Addition des achats]
Énoncé. Amina achète des tomates (\zh{西红柿} xīhóngshì) pour 500 F et du ndolé (\zh{苦菜} kǔcài) pour 300 F.\\
1) Écrire l'addition en chinois : « 500 plus 300 égale 800 ».\\
2) Demander le total : « Combien en tout ? »\\
3) Répondre en phrase complète avec la monnaie (franc = \zh{法郎} fǎláng).
\tcblower
Solution détaillée.\\
1) Les centaines s'écrivent avec \cle{百} (bǎi) : 500 = \zh{五百} (wǔbǎi), 300 = \zh{三百} (sānbǎi), 800 = \zh{八百} (bābǎi).\\
\zh{五百加三百等于八百。} (Wǔbǎi jiā sānbǎi děngyú bābǎi.)\\
2) \zh{一共多少钱？} (Yígòng duōshao qián ?) — \cle{一共} résume le total.\\
3) Phrase complète : \zh{一共八百法郎。} (Yígòng bābǎi fǎláng.) Ou avec les articles : \zh{西红柿和苦菜一共八百法郎。} (Xīhóngshì hé kǔcài yígòng bābǎi fǎláng.)\\
\cle{Attention :} Pas de classificateur entre le nombre et l'unité monétaire (pas de « 个 »).
\end{exoresolu}

\courssub{La soustraction : 减法}

\begin{methode}[Dire et écrire une soustraction : 减法]
La structure miroir de l'addition est :
\begin{enumerate}
\item \cle{A 减 B 等于 C} (A jiǎn B děngyú C) : « A moins B égale C ».
\item Ordre strict : le nombre de départ (le plus grand) \textbf{avant} \cle{减} (jiǎn), le nombre retranché \textbf{après}.
\item Question : \zh{A 减 B 等于多少？}
\item Dans la vie courante (rendu de monnaie), on utilise souvent le verbe \cle{剩} (shèng, rester) : \zh{还剩多少？} (Hái shèng duōshao ? « Combien reste-t-il encore ? »).
\item \cle{给} (gěi, donner) + \cle{您} (nín, vous poli) pour rendre la monnaie : \zh{找您五块} (zhǎo nín wǔ kuài, « je vous rends 5 yuans »).
\end{enumerate}
\end{methode}

\begin{exemplebox}[Exemples de soustractions]
\begin{center}
\begin{tabularx}{\linewidth}{|X|X|X|}
\hline
\rowcolor{poporangeL} \textbf{Chinois} & \textbf{Structure} & \textbf{Traduction} \\
\hline
\zh{十减三等于七。} (Shí jiǎn sān děngyú qī.) & 10 - 3 = 7 & Dix moins trois font sept. \\
\zh{五千减两千五等于两千五。} (Wǔqiān jiǎn liǎngqiānwǔ děngyú liǎngqiānwǔ.) & 5000 - 2500 = 2500 & 5000 moins 2500 font 2500. \\
\zh{还剩两千五百法郎。} (Hái shèng liǎngqiānwǔbǎi fǎláng.) & 还 + 剩 + Nombre + Monnaie & Il reste encore 2500 francs. \\
\zh{找您两千法郎。} (Zhǎo nín liǎngqiān fǎláng.) & 找 + Personne + Montant & Je vous rends 2000 francs. \\
\hline
\end{tabularx}
\end{center}
\end{exemplebox}

\begin{exoresolu}[La monnaie à rendre : Soustraction réelle]
Énoncé. Un repas coûte 2 500 F. Le client paie avec un billet de 5 000 F.\\
1) Écrire l'opération formelle avec \cle{减}.\\
2) Demander : « Combien me reste-t-il ? » (client) / « Combien je lui rends ? » (vendeur).\\
3) Réponse du vendeur.
\tcblower
Solution détaillée.\\
1) 5 000 = \zh{五千} (wǔqiān). 2 500 = \zh{两千五} (liǎngqiānwǔ, forme courante) ou \zh{二千五百} (èrqiān wǔbǎi).\\
\zh{五千减两千五等于两千五。} (Wǔqiān jiǎn liǎngqiānwǔ děngyú liǎngqiānwǔ.)\\
2) Client : \zh{还剩多少钱？} (Hái shèng duōshao qián ?) — Vendeur (calcul) : \zh{五千减两千五……} \\
3) Vendeur : \zh{还剩两千五百法郎，找您。} (Hái shèng liǎngqiānwǔbǎi fǎláng, zhǎo nín.)\\
\cle{Distinction :} \zh{减} = opération mathématique écrite ; \zh{剩} = résultat concret (reste) ; \zh{找} = action de rendre la monnaie (verbe transactionnel).
\end{exoresolu}

\courssub{Les monnaies et les grands nombres}

\begin{definition}[Vocabulaire : Monnaies et unités]
\begin{center}
\begin{tabularx}{\linewidth}{|X|X|X|X|}
\hline
\rowcolor{popgreenL} \textbf{Caractères} & \textbf{Pinyin} & \textbf{Nature} & \textbf{Traduction} \\
\hline
人民币 & rénmínbì & nom & Monnaie chinoise (RMB) \\
元 / 块 & yuán / kuài & nom / unité & Yuan (officiel / oral) \\
角 / 毛 & jiǎo / máo & nom / unité & Jiao (0,1 yuan) / Mao (oral) \\
分 & fēn & nom / unité & Fen (0,01 yuan) \\
法郎 & fǎláng & nom & Franc (CFA) \\
美元 & měiyuán & nom & Dollar US \\
欧元 & ōuyuán & nom & Euro \\
\hline
\end{tabularx}
\end{center}
\end{definition}

\begin{propriete}[Règles des grands nombres : la base 万 (wàn)]
\begin{enumerate}
\item Le chinois groupe par \textbf{quatre zéros} (dizaines de mille), pas par trois (mille).
\item \cle{百} (bǎi) = 100 ; \cle{千} (qiān) = 1 000 ; \cle{万} (wàn) = 10 000.
\item \textbf{Conversion clé :} 10 000 = \zh{一万} (yí wàn) — \textbf{jamais} « 十千 ».
\item 36 000 = 3,6 × 10 000 = \zh{三万六千} (sān wàn liù qiān). \textbf{Interdit :} « 三十六千 ».
\item 100 000 = \zh{十万} (shí wàn) ; 1 000 000 = \zh{一百万} (yì bǎi wàn).
\item Zéros intermédiaires : 10 001 = \zh{一万零一} (yí wàn líng yī) ; 10 100 = \zh{一万零一百} (yí wàn líng yì bǎi).
\end{enumerate}
\end{propriete}

\begin{exemplebox}[Prix et conversion : Structure « Sujet + Nombre + Monnaie »]
\begin{center}
\begin{tabularx}{\linewidth}{|X|X|}
\hline
\rowcolor{poppurpleL} \textbf{Chinois} & \textbf{Traduction} \\
\hline
\zh{这个手机八百块。} (Zhège shǒujī bābǎi kuài.) & Ce téléphone coûte 800 yuans. (Pas de verbe « coûter »). \\
\zh{一块钱大约八十五法郎。} (Yí kuài qián dàyuē bāshíwǔ fǎláng.) & 1 yuan ≈ 85 francs. (\cle{大约} dàyuē = environ). \\
\zh{五万法郎等于多少块？} (Wǔwàn fǎláng děngyú duōshao kuài?) & 50 000 francs = combien de yuans ? \\
\hline
\end{tabularx}
\end{center}
\end{exemplebox}

\begin{exoresolu}[Convertir et manipuler les prix]
Énoncé. 1) Écrire en chiffres et en caractères : « dix mille », « trente-six mille », « cent mille ».\\
2) Traduire : « Ce vêtement coûte 360 yuans. » (\zh{衣服} yīfu).\\
3) Traduire : « 10 000 francs CFA font environ 118 yuans. » (Taux fictif pour l'exercice).
\tcblower
Solution détaillée.\\
1) 10 000 = \zh{一万} (yí wàn) ; 36 000 = \zh{三万六千} (sān wàn liù qiān) ; 100 000 = \zh{十万} (shí wàn).\\
2) \zh{这件衣服三百六十块。} (Zhè jiàn yīfu sānbǎi liùshí kuài.) — Classificateur de \zh{衣服} : \cle{件} (jiàn).\\
3) \zh{一万法郎大约是一百一十八块。} (Yí wàn fǎláng dàyuē shì yìbǎi yīshíbā kuài.) — \cle{是} (shì) pour l'équivalence.
\end{exoresolu}

\courssub{La structure 除了……以外，……}

\begin{methode}[Maîtriser 除了……以外，…… (Chúle A yǐwài, …)]
Cette structure introduit un élément \cle{A} (nom, pronom, verbe) pour l'exclure ou l'ajouter. \cle{以外} (yǐwài) = « en dehors de ». La forme complète est exigée au Bac écrit.
\end{methode}

\begin{propriete}[Deux emplois opposés : Le mot-outil décide]
\begin{center}
\begin{tabularx}{\linewidth}{|X|X|X|X|}
\hline
\rowcolor{poppinkL} \textbf{Emploi} & \textbf{Sens} & \textbf{Mot-outil 2e partie} & \textbf{Exemple} \\
\hline
\textbf{Exclusion} & « Sauf / À part » (les autres \textbf{ne font pas}) & \cle{都} (dōu) + V / \cle{没/不} + V & \zh{除了他以外，没人来。} (Sauf lui, personne n'est venu.) \\
\textbf{Addition} & « En plus de / Outre » (lui \textbf{aussi} fait) & \cle{还} (hái) / \cle{也} (yě) + V & \zh{除了汉语以外，我还学法语。} (En plus du chinois, j'apprends le français.) \\
\hline
\end{tabularx}
\end{center}
\cle{Règle d'or :} Sans \cle{还/也}, la phrase bascule vers l'exclusion. \zh{除了足球，他喜欢篮球} $\rightarrow$ ambigu (sauf le foot, il aime le basket ?). Correct : \zh{除了足球以外，他\cle{还}喜欢篮球。}
\end{propriete}

\begin{attention}[Pièges francophones sur 除了……以外]
\begin{enumerate}
\item \textbf{Oublier 还/也} en sens addition : « 除了中文，我学英文 » $\times$ $\rightarrow$ \zh{除了中文以外，我\cle{还}学英文。} $\checkmark$
\item \textbf{Mettre une phrase complète après 除了} : « 除了他来了以外 » $\times$. Après 除了 : Nom / Pronom / Verbe seul (\zh{除了迟到}, \zh{除了他}).
\item \textbf{Confondre 都 (exclusion) et 还 (addition)} : \zh{除了周一，我都上班} (exclusion : je travaille tous les autres jours) vs \zh{除了周一，我还上班} (addition : en plus des autres jours, je travaille aussi le lundi — rare).
\item \textbf{Omettre 以外 à l'écrit} : Toléré à l'oral (\zh{除了他，没人来}), \textbf{interdit au Bac écrit}.
\end{enumerate}
\end{attention}

\begin{exoresolu}[Exclusion vs Addition : Transformer]
Énoncé. 1) « À part le lundi, je travaille tous les jours. » (Exclusion)\\
2) « En plus du football, il aime aussi le basketball. » (Addition)\\
3) Compléter : \zh{除了米饭以外，我们还吃\_\_\_\_。} (ndolé = \zh{苦菜} kǔcài)
\tcblower
Solution détaillée.\\
1) Exclusion $\rightarrow$ \cle{都} + verbe affirmatif (ou \cle{没/不} + V). \zh{除了星期一以外，我每天\cle{都}工作。} (Chúle xīngqīyī yǐwài, wǒ měitiān \cle{dōu} gōngzuò.)\\
2) Addition $\rightarrow$ \cle{还/也} + verbe. \zh{除了足球以外，他\cle{还}喜欢篮球。} (Chúle zúqiú yǐwài, tā \cle{hái} xǐhuan lánqiú.)\\
3) Addition (ndolé s'ajoute au riz) $\rightarrow$ \cle{还} avant \zh{吃}. \zh{除了米饭以外，我们\cle{还}吃\cle{苦菜}。} (Chúle mǐfàn yǐwài, wǒmen \cle{hái} chī \cle{kǔcài}.)
\end{exoresolu}

\courssub{Synthèse comparative : Les trois structures du calcul}

\begin{propriete}[Tableau récapitulatif pour le Bac]
\begin{center}
\begin{tabularx}{\linewidth}{|X|X|X|X|}
\hline
\rowcolor{popgoldL} \textbf{Notion} & \textbf{Structure type} & \textbf{Exemple (Car. + Pin.)} & \textbf{Traduction} \\
\hline
Addition & A \cle{加} B \cle{等于} C & \zh{五百加三百等于八百。} (Wǔbǎi jiā sānbǎi děngyú bābǎi.) & 500 + 300 = 800 \\
Soustraction (maths) & A \cle{减} B \cle{等于} C & \zh{五千减两千五等于两千五。} (Wǔqiān jiǎn liǎngqiānwǔ děngyú liǎngqiānwǔ.) & 5000 - 2500 = 2500 \\
Reste (vie courante) & \cle{还剩} + Nombre + Monnaie & \zh{还剩两千五百法郎。} (Hái shèng liǎngqiānwǔbǎi fǎláng.) & Il reste 2500 F \\
Rendre monnaie & \cle{找} + Personne + Montant & \zh{找您两千五。} (Zhǎo nín liǎngqiānwǔ.) & Je vous rends 2500 \\
Total & \cle{一共} + Nombre + Monnaie & \zh{一共八百块。} (Yígòng bābǎi kuài.) & Total 800 yuans \\
\hline
\end{tabularx}
\end{center}
\end{propriete}

\courssub{Expression écrite : Rédiger un dialogue d'achat}

\begin{methode}[Composer un dialogue de marché (6 répliques = 3 tours)]
\begin{enumerate}
\item \textbf{Salutation} : \zh{你好！} / \zh{老板好！}
\item \textbf{Question prix} : \zh{……多少钱一斤？} (Classificateur \cle{斤} jīn = 500g pour fruits/légumes).
\item \textbf{Réponse prix} : Nombre + Monnaie (pas de classificateur).
\item \textbf{Calcul total} : Utiliser \cle{加……等于……} ou \cle{一共……}.
\item \textbf{Paiement / Monnaie} : \cle{给您……} (client) / \cle{找您……} (vendeur).
\item \textbf{Politesse} : \zh{谢谢，再见！}
\end{enumerate}
\end{methode}

\begin{exoresolu}[Modèle de dialogue : 6 répliques]
Énoncé. Client achète bananes (\zh{香蕉} xiāngjiāo) 2 000 F et arachides (\zh{花生} huāshēng) 1 000 F. Paie 5 000 F. Demande la monnaie.
\tcblower
Solution — Modèle corrigé (6 répliques).\\
\zh{A：你好！香蕉多少钱一斤？} (Nǐ hǎo! Xiāngjiāo duōshao qián yì jīn?) \\
\zh{B：两千法郎。花生一千法郎。} (Liǎngqiān fǎláng. Huāshēng yìqiān fǎláng.) \\
\zh{A：两千加一千等于三千，一共三千法郎。} (Liǎngqiān jiā yìqiān děngyú sānqiān, yígòng sānqiān fǎláng.) \\
\zh{B：对，一共三千。} (Duì, yígòng sānqiān.) \\
\zh{A：给你五千，请找我两千。} (Gěi nǐ wǔqiān, qǐng zhǎo wǒ liǎngqiān.) \\
\zh{B：好的，找您两千。谢谢！} (Hǎo de, zhǎo nín liǎngqiān. Xièxie!) \\
\cle{Justification :} Réplique 3 = structure 加法 + 一共. Réplique 5 = 给 + 找. Vocabulaire précis (斤, 法郎).
\end{exoresolu}

\begin{exercice}[Entraînement : À vous de jouer]
\begin{enumerate}
\item \textbf{Calcul écrit :} Écrire en chinois (caractères + pinyin) : 8 000 + 2 000 = 10 000 ; 50 000 - 30 000 = 20 000.
\item \textbf{Conversion :} Traduire : « 20 000 francs CFA font environ 300 yuans. »
\item \textbf{Structure 除了 :} Traduire les deux sens :\\
a) « À part le riz, nous ne mangeons rien d'autre. » (Exclusion)\\
b) « En plus du chinois, j'étudie aussi l'anglais. » (Addition)
\item \textbf{Dialogue :} Rédiger un dialogue de 6 répliques : Achat de 3 kg d'oranges (\zh{橙子} chéngzi) à 1 500 F/kg et 2 kg de mangues (\zh{芒果} mángguǒ) à 2 000 F/kg. Paiement 10 000 F. Calcul total + monnaie.
\end{enumerate}
\end{exercice}

\begin{corrige}[Exercice Leçon 26]
\begin{enumerate}
\item \zh{八千加两千等于一万。} (Bāqiān jiā liǎngqiān děngyú yí wàn.) ; \zh{五万减三万等于两万。} (Wǔwàn jiǎn sānwàn děngyú liǎng wàn.)
\item \zh{两万法郎大约是三百块。} (Liǎng wàn fǎláng dàyuē shì sānbǎi kuài.)
\item a) Exclusion : \zh{除了米饭以外，我们什么都不吃。} (Chúle mǐfàn yǐwài, wǒmen shénme dōu bù chī.)\\
b) Addition : \zh{除了汉语以外，我还学英语。} (Chúle Hànyǔ yǐwài, wǒ hái xué Yīngyǔ.)
\item \textbf{Dialogue type :}\\
\zh{A：橙子多少钱一斤？} \\
\zh{B：一千五百法郎。芒果两千法郎。} \\
\zh{A：三斤橙子四千五，两斤芒果四千。四千五加四千等于八千五，一共八千五百法郎。} \\
\zh{B：对，一共八千五。} \\
\zh{A：给你一万，找我一千五。} \\
\zh{B：好，找您一千五百。再见！}
\end{enumerate}
\end{corrige}

\begin{savaistu}[Le Franc CFA et le Yuan : Realités économiques]
Le \cle{Franc CFA} (XAF/ XOF) est arrimé à l'\cle{Euro} (1 € = 655,957 F CFA). Le \cle{Yuan} (CNY / RMB) flotte dans une bande gérée par la Banque Populaire de Chine.
\begin{itemize}
\item Taux indicatif 2024 : 1 Yuan ≈ 80--85 F CFA (variable).
\item Au Cameroun, les paiements mobiles (Mobile Money : MTN MoMo, Orange Money) dominent le quotidien. En Chine, \cle{Alipay} (支付宝 Zhīfùbǎo) et \cle{WeChat Pay} (微信支付 Wēixìn Zhīfù) sont omniprésents (QR code), l'espèce devient rare.
\item Négociation : Au marché camerounais (\zh{叫价} jiàojià), on négocie souvent. En Chine, prix fixes en supermarché, négociation possible sur les marchés de gros (\zh{批发市场} pīfā shìchǎng) ou pour gros volumes.
\item Vocabulaire utile : \zh{零钱} (língqián, monnaie / petite coupure), \zh{扫码} (sǎomǎ, scanner le code QR), \zh{转账} (zhuǎnzhàng, virement).
\end{itemize}

\begin{center}
\begin{tabularx}{\linewidth}{|X|X|X|X|}
\hline
\rowcolor{popblueL} \textbf{Caractères} & \textbf{Pinyin} & \textbf{Nature} & \textbf{Traduction} \\
\hline
加 & jiā & verbe & Ajouter, plus (+) \\
减 & jiǎn & verbe & Soustraire, moins (-) \\
等于 & děngyú & verbe & Égaler (=) \\
一共 & yígòng & adv. & En tout, au total \\
剩 & shèng & verbe & Rester (reste) \\
找 & zhǎo & verbe & Rendre (monnaie), chercher \\
给 & gěi & verbe/prép. & Donner, à (préposition) \\
除了 & chúle & prép. & Excepté, à part \\
以外 & yǐwài & nom/loc. & En dehors de, à l'extérieur \\
还 & hái & adv. & Encore, aussi (addition) \\
也 & yě & adv. & Aussi (addition) \\
都 & dōu & adv. & Tous (exclusion/distributif) \\
万 & wàn & nombre & Dix mille (10 000) \\
千 & qiān & nombre & Mille (1 000) \\
百 & bǎi & nombre & Cent (100) \\
块 & kuài & classif/unite & Yuan (oral) \\
元 & yuán & classif/unite & Yuan (écrit/officiel) \\
法郎 & fǎláng & nom & Franc CFA \\
美元 & měiyuán & nom & Dollar US \\
欧元 & ōuyuán & nom & Euro \\
大约 & dàyuē & adv. & Environ, à peu près \\
斤 & jīn & classif. & Demi-kilo (500 g) \\
西红柿 & xīhóngshì & nom & Tomate \\
苦菜 & kǔcài & nom & Ndolé (légume amer) \\
香蕉 & xiāngjiāo & nom & Banane \\
花生 & huāshēng & nom & Arachide \\
橙子 & chéngzi & nom & Orange \\
芒果 & mángguǒ & nom & Mangue \\
手机 & shǒujī & nom & Téléphone portable \\
衣服 & yīfu & nom & Vêtement \\
件 & jiàn & classif. & Classif. vêtements, affaires \\
\hline
\end{tabularx}
\end{center}
\end{savaistu}


\newpage

\coursec{Exercices}

\courssub{Je vérifie mes connaissances}

\begin{exercice}[QCM : la bonne structure]
Entoure la bonne réponse.
\begin{enumerate}
\item « 8 + 5 = 13 » se dit en chinois :
\begin{itemize}
\item[a)] 八减五等于十三。
\item[b)] 八加五等于十三。
\item[c)] 八乘五等于十三。
\end{itemize}
\item Pour dire « En plus du chinois, il apprend aussi l'anglais », on utilise :
\begin{itemize}
\item[a)] 除了汉语以外，他还学英语。
\item[b)] 除了汉语以外，他都学英语。
\item[c)] 除了汉语以外，他只学英语。
\end{itemize}
\item « 100 moins 40 égalent 60 » se dit :
\begin{itemize}
\item[a)] 一百加四十等于六十。
\item[b)] 一百减四十等于六十。
\item[c)] 一百等于四十减六十。
\end{itemize}
\item Le symbole monétaire du yuan chinois est :
\begin{itemize}
\item[a)] 元 (yuán)
\item[b)] 法郎 (fǎláng)
\item[c)] 美元 (měiyuán)
\end{itemize}
\end{enumerate}
\end{exercice}

\begin{exercice}[Vrai ou faux ? Justifie]
Dis si chaque phrase est vraie ou fausse, puis justifie ta réponse en une phrase en français.
\begin{enumerate}
\item Dans la structure \zh{除了……以外，……还……}, la deuxième partie s'ajoute à la première (sens d'addition).
\item \zh{除了……以外，……都……} exprime une addition.
\item En chinois, « 2 + 3 = 5 » se dit \zh{二加三等于五。}
\item Le franc CFA se dit \zh{美元} en chinois.
\end{enumerate}
\end{exercice}

\begin{exercice}[Vocabulaire : les monnaies]
Complète le tableau en donnant la traduction française de chaque monnaie.
\begin{center}
\begin{tabularx}{\linewidth}{|X|X|X|}\hline
\rowcolor{popblueL} Caractères & Pinyin & Traduction \\ \hline
人民币 & rénmínbì & \ldots \\ \hline
美元 & měiyuán & \ldots \\ \hline
欧元 & ōuyuán & \ldots \\ \hline
法郎 & fǎláng & \ldots \\ \hline
日元 & rìyuán & \ldots \\ \hline
\end{tabularx}
\end{center}
\end{exercice}

\begin{exercice}[Associe le pinyin aux caractères]
Relie chaque mot chinois à son pinyin.
\begin{enumerate}
\item 加 \hspace{2cm} a) děngyú
\item 减 \hspace{2cm} b) chúle
\item 等于 \hspace{2cm} c) jiā
\item 除了 \hspace{2cm} d) jiǎn
\end{enumerate}
\end{exercice}

\courssub{Je m'entraîne}

\begin{exercice}[Écris les calculs en caractères]
Écris chaque calcul en chinois, avec une phrase complète (caractères + pinyin).
\begin{enumerate}
\item 7 + 8 = 15
\item 25 + 75 = 100
\item 1000 $-$ 250 = 750
\item 5000 $-$ 1800 = 3200
\end{enumerate}
\end{exercice}

\begin{exercice}[Traduction français $\rightarrow$ chinois]
Traduis les phrases suivantes en chinois (caractères et pinyin).
\begin{enumerate}
\item En plus du chinois, elle apprend aussi l'anglais.
\item Sauf lundi, je travaille tous les jours.
\item 100 moins 35 égalent 65.
\item En plus de l'euro, il a aussi apporté des dollars.
\end{enumerate}
\end{exercice}

\begin{exercice}[Traduction chinois $\rightarrow$ français]
Traduis les phrases suivantes en français.
\begin{enumerate}
\item \zh{除了美元以外，他还带了欧元。} (Chúle měiyuán yǐwài, tā hái dài le ōuyuán.)
\item \zh{二百减八十等于一百二十。} (Èrbǎi jiǎn bāshí děngyú yībǎi èrshí.)
\item \zh{除了星期天以外，我每天都上课。} (Chúle xīngqītiān yǐwài, wǒ měitiān dōu shàngkè.)
\end{enumerate}
\end{exercice}

\begin{exercice}[Transformation et commentaire]
\begin{enumerate}
\item Transforme la phrase \zh{除了汉语以外，我也学英语。} (addition) en une phrase d'exclusion avec \zh{只}.
\item Commente en français le changement de sens entre les deux phrases.
\end{enumerate}
\end{exercice}

\begin{exercice}[Dialogue à trous : au marché de Yaoundé]
Complète le dialogue avec les mots suivants : \zh{加}, \zh{减}, \zh{等于}, \zh{块}, \zh{除了……以外，……还……}.
\begin{textebox}[Au marché]
A：这个多少钱？\hspace{1cm} (Zhège duōshao qián ?)\\
B：十五 \ldots\ldots 。\\
A：那个呢？\\
B：二十块。十五 \ldots\ldots 二十 \ldots\ldots 三十五。\\
A：我给你五十块。五十 \ldots\ldots 三十五 \ldots\ldots 十五。\\
B：对，我找您十五块。\\
A：\ldots\ldots 这个 \ldots\ldots ，我 \ldots\ldots 买那个。
\end{textebox}
\end{exercice}

\courssub{Je me prépare au Bac}

\begin{exercice}[Texte et questions de compréhension]
Lis le texte suivant, puis réponds aux questions en français.
\begin{textebox}[Les achats de Lili]
明天是莉莉的生日。妈妈给了她一百块钱。莉莉去商店买了一个本子，五块钱，又买了一支笔，三块钱。五加三等于八，她一共花了八块钱。一百减八等于九十二，她还剩九十二块钱。除了本子和笔以外，她还想买一个小蛋糕，二十块钱。妈妈说：「除了蛋糕以外，你还可以买一瓶水。」莉莉很高兴。\\
(Míngtiān shì Lìli de shēngrì. Māma gěi le tā yībǎi kuài qián. Lìli qù shāngdiàn mǎi le yí ge běnzi, wǔ kuài qián, yòu mǎi le yì zhī bǐ, sān kuài qián. Wǔ jiā sān děngyú bā, tā yígòng huā le bā kuài qián. Yībǎi jiǎn bā děngyú jiǔshí'èr, tā hái shèng jiǔshí'èr kuài qián. Chúle běnzi hé bǐ yǐwài, tā hái xiǎng mǎi yí ge xiǎo dàngāo, èrshí kuài qián. Māma shuō : « Chúle dàngāo yǐwài, nǐ hái kěyǐ mǎi yì píng shuǐ. » Lìli hěn gāoxìng.)
\end{textebox}
\begin{enumerate}
\item Combien d'argent la mère de Lili lui a-t-elle donné ?
\item Quels articles Lili a-t-elle déjà achetés et à quels prix ?
\item Écris en chinois le calcul du total de ses achats.
\item Combien lui reste-t-il ? Écris le calcul en chinois.
\item Question de langue : explique la différence entre \zh{除了……以外，……还……} et \zh{除了……以外，……都……} en t'appuyant sur deux exemples traduits.
\end{enumerate}
\end{exercice}

\begin{exercice}[Thème et version]
\textbf{Thème.} Traduis en chinois (caractères et pinyin) :
\begin{enumerate}
\item En plus du français, nous apprenons aussi le chinois et l'anglais.
\item Sauf le dimanche, mon père travaille tous les jours.
\item 50 plus 30 égalent 80 ; 80 moins 20 égalent 60.
\end{enumerate}
\textbf{Version.} Traduis en français :
\begin{enumerate}
\item \zh{除了人民币以外，很多中国人也用美元。} (Chúle rénmínbì yǐwài, hěn duō Zhōngguó rén yě yòng měiyuán.)
\item \zh{一千减二百五十等于七百五十。} (Yìqiān jiǎn èrbǎi wǔshí děngyú qībǎi wǔshí.)
\end{enumerate}
\end{exercice}

\begin{exercice}[Problème contextualisé et expression écrite]
Tu as 10 000 FCFA. Tu achètes un cahier à 1 500 FCFA et un sac à 6 000 FCFA.
\begin{enumerate}
\item Écris en chinois (caractères et pinyin) le calcul du total de tes achats.
\item Écris en chinois le calcul de la monnaie qu'on doit te rendre.
\item Dis en une phrase chinoise avec quelles monnaies on peut payer au Cameroun et en Chine.
\end{enumerate}
\textbf{Expression écrite.} Rédige un court texte en chinois (60 à 80 caractères) intitulé \zh{我去市场买东西} (Je vais au marché faire des achats) : raconte un achat au marché de ta ville en utilisant au moins une addition, une soustraction et la structure \zh{除了……以外，……还……}.
\end{exercice}

\newpage

\coursec{Corrigés des exercices}

\begin{corrige}[Exercice 1 — QCM : la bonne structure]
\begin{enumerate}
\item \textbf{b)} \zh{八加五等于十三。} (\emph{Bā jiā wǔ děngyú shísān.}) « 8 + 5 = 13 ». Le verbe \zh{加} signifie « ajouter » ; a) \zh{减} = soustraire ; c) \zh{乘} = multiplier.
\item \textbf{a)} \zh{除了汉语以外，他还学英语。} (\emph{Chúle Hànyǔ yǐwài, tā hái xué Yīngyǔ.}) Sens d'addition : \zh{还} « aussi ». b) \zh{都} exprime l'exclusion (« il n'apprend que l'anglais, pas le chinois ») ; c) \zh{只} « seulement » est contradictoire avec l'idée d'addition.
\item \textbf{b)} \zh{一百减四十等于六十。} (\emph{Yībǎi jiǎn sìshí děngyú liùshí.}) « 100 $-$ 40 = 60 ». a) est une addition ; c) inverse l'ordre des termes.
\item \textbf{a)} \zh{元} (\emph{yuán}) est l'unité monétaire chinoise. b) \zh{法郎} = franc ; c) \zh{美元} = dollar américain.
\end{enumerate}
\end{corrige}

\begin{corrige}[Exercice 2 — Vrai ou faux ? Justifie]
\begin{enumerate}
\item \textbf{Vrai.} Avec \zh{还} « aussi », l'élément de la deuxième partie s'ajoute à celui de la première : \zh{除了汉语以外，他还学英语。} = il apprend le chinois \emph{et} l'anglais.
\item \textbf{Faux.} \zh{除了……以外，……都……} exprime une \emph{exclusion} : \zh{除了星期一以外，我每天都上课。} « Sauf lundi, j'ai cours tous les jours » — lundi est exclu.
\item \textbf{Vrai.} \zh{二加三等于五。} (\emph{Èr jiā sān děngyú wǔ.}) « 2 + 3 = 5 » : nombre + \zh{加} + nombre + \zh{等于} + résultat.
\item \textbf{Faux.} \zh{美元} signifie « dollar américain » ; le franc se dit \zh{法郎} (\emph{fǎláng}), donc le franc CFA = \zh{非洲法郎} ou simplement \zh{法郎} dans le contexte camerounais.
\end{enumerate}
\end{corrige}

\begin{corrige}[Exercice 3 — Vocabulaire : les monnaies]
\begin{center}
\begin{tabularx}{\linewidth}{|X|X|X|}\hline
\rowcolor{popblueL} Caractères & Pinyin & Traduction \\ \hline
人民币 & rénmínbì & le renminbi (monnaie chinoise, « monnaie du peuple ») \\ \hline
美元 & měiyuán & le dollar américain \\ \hline
欧元 & ōuyuán & l'euro \\ \hline
法郎 & fǎláng & le franc (au Cameroun : le franc CFA) \\ \hline
日元 & rìyuán & le yen japonais \\ \hline
\end{tabularx}
\end{center}
\textbf{Point de vigilance :} ne pas confondre \zh{元} (unité du yuan) et \zh{日元} (yen) : le premier caractère indique le pays (\zh{日} = Japon, \zh{美} = États-Unis, \zh{欧} = Europe). À l'oral, on compte l'argent avec \zh{块} (\emph{kuài}) : \zh{十五块钱} « quinze yuans ».
\end{corrige}

\begin{corrige}[Exercice 4 — Associe le pinyin aux caractères]
1 — c) : \zh{加} = \emph{jiā} « ajouter ».\\
2 — d) : \zh{减} = \emph{jiǎn} « soustraire ».\\
3 — a) : \zh{等于} = \emph{děngyú} « égaler ».\\
4 — b) : \zh{除了} = \emph{chúle} « sauf, en plus de ».\\
\textbf{Astuce :} \zh{加} et \zh{减} partagent la clé \zh{力} « force » ; attention à ne pas confondre les tons : \emph{jiā} (premier ton) / \emph{jiǎn} (troisième ton).
\end{corrige}

\begin{corrige}[Exercice 5 — Écris les calculs en caractères]
\begin{enumerate}
\item \zh{七加八等于十五。} (\emph{Qī jiā bā děngyú shíwǔ.}) « 7 + 8 = 15 ».
\item \zh{二十五加七十五等于一百。} (\emph{Èrshíwǔ jiā qīshíwǔ děngyú yībǎi.}) « 25 + 75 = 100 ».
\item \zh{一千减二百五十等于七百五十。} (\emph{Yìqiān jiǎn èrbǎi wǔshí děngyú qībǎi wǔshí.}) « 1000 $-$ 250 = 750 ».
\item \zh{五千减一千八百等于三千二百。} (\emph{Wǔqiān jiǎn yìqiān bābǎi děngyú sānqiān èrbǎi.}) « 5000 $-$ 1800 = 3200 ».
\end{enumerate}
\textbf{Points de vigilance :} structure fixe \cle{nombre} + \zh{加}/\zh{减} + \cle{nombre} + \zh{等于} + \cle{résultat} ; ne pas oublier le point chinois \zh{。} ; 1800 se dit \zh{一千八百} (et non \zh{十八百}).
\end{corrige}

\begin{corrige}[Exercice 6 — Traduction français $\rightarrow$ chinois]
\begin{enumerate}
\item \zh{除了汉语以外，她还学英语。} (\emph{Chúle Hànyǔ yǐwài, tā hái xué Yīngyǔ.}) — addition : \zh{还}.
\item \zh{除了星期一以外，我每天都工作。} (\emph{Chúle xīngqīyī yǐwài, wǒ měitiān dōu gōngzuò.}) — exclusion : \zh{都} avec \zh{每天}.
\item \zh{一百减三十五等于六十五。} (\emph{Yībǎi jiǎn sānshíwǔ děngyú liùshíwǔ.})
\item \zh{除了欧元以外，他还带了美元。} (\emph{Chúle ōuyuán yǐwài, tā hái dài le měiyuán.})
\end{enumerate}
\textbf{Points de vigilance :} « en plus de » $\rightarrow$ \zh{除了……以外，……还……} ; « sauf » $\rightarrow$ \zh{除了……以外，……都……} ; ne jamais traduire « sauf » par \zh{还}.
\end{corrige}

\begin{corrige}[Exercice 7 — Traduction chinois $\rightarrow$ français]
\begin{enumerate}
\item « En plus des dollars, il a aussi apporté des euros. » (addition : \zh{还}).
\item « 200 moins 80 égalent 120. »
\item « Sauf le dimanche, je vais en cours tous les jours. » (exclusion : \zh{都}).
\end{enumerate}
\end{corrige}

\begin{corrige}[Exercice 8 — Transformation et commentaire]
\begin{enumerate}
\item \zh{除了汉语以外，我只学英语。} (\emph{Chúle Hànyǔ yǐwài, wǒ zhǐ xué Yīngyǔ.}) « Sauf le chinois, je n'apprends que l'anglais » — c'est-à-dire : je n'apprends \emph{pas} le chinois, j'apprends uniquement l'anglais.
\item \textbf{Commentaire.} La phrase d'origine avec \zh{也} (« aussi ») exprime une \emph{addition} : le chinois et l'anglais sont tous deux appris. Avec \zh{只} (« seulement »), la phrase devient une \emph{exclusion} : le chinois est exclu de l'ensemble des langues apprises. Le sens est donc radicalement inversé : on passe de « deux langues apprises » à « une seule langue apprise ». C'est le mot-outil de la deuxième partie (\zh{还}/\zh{也} contre \zh{都}/\zh{只}) qui détermine le sens de toute la structure.
\end{enumerate}
\end{corrige}

\begin{corrige}[Exercice 9 — Dialogue à trous : au marché de Yaoundé]
Dialogue complété :
\begin{textebox}[Au marché — corrigé]
A：这个多少钱？\hspace{1cm} (Zhège duōshao qián ?)\\
B：十五\textbf{块}。\\
A：那个呢？\\
B：二十块。十五\textbf{加}二十\textbf{等于}三十五。\\
A：我给你五十块。五十\textbf{减}三十五\textbf{等于}十五。\\
B：对，我找您十五块。\\
A：\textbf{除了}这个\textbf{以外}，我\textbf{还}买那个。
\end{textebox}
Traduction : « A : Combien coûte ceci ? B : Quinze yuans. A : Et cela ? B : Vingt yuans. Quinze plus vingt égalent trente-cinq. A : Je vous donne cinquante yuans. Cinquante moins trente-cinq égalent quinze. B : C'est juste, je vous rends quinze yuans. A : En plus de celui-ci, j'achète aussi celui-là. »\\
\textbf{Point de vigilance :} \zh{块} (\emph{kuài}) est la forme orale de \zh{元} pour compter l'argent ; la dernière réplique exige la structure complète d'addition \zh{除了……以外，……还……}.
\end{corrige}

\begin{corrige}[Exercice 10 — Texte et questions de compréhension]
\begin{enumerate}
\item Sa mère lui a donné \textbf{100 yuans} (\zh{一百块钱}).
\item Elle a acheté \textbf{un cahier à 5 yuans} (\zh{一个本子，五块钱}) et \textbf{un stylo à 3 yuans} (\zh{一支笔，三块钱}).
\item \zh{五加三等于八。} (\emph{Wǔ jiā sān děngyú bā.}) « 5 + 3 = 8 » — elle a dépensé 8 yuans en tout.
\item \zh{一百减八等于九十二。} (\emph{Yībǎi jiǎn bā děngyú jiǔshí'èr.}) « 100 $-$ 8 = 92 » — il lui reste 92 yuans.
\item \textbf{Question de langue.} \zh{除了……以外，……还……} exprime l'\emph{addition} : \zh{除了本子和笔以外，她还想买一个小蛋糕。} « En plus du cahier et du stylo, elle veut aussi acheter un petit gâteau » — le gâteau s'ajoute aux achats. \zh{除了……以外，……都……} exprime l'\emph{exclusion} : \zh{除了星期天以外，我每天都上课。} « Sauf le dimanche, je vais en cours tous les jours » — le dimanche est exclu. C'est l'adverbe de la seconde partie (\zh{还} contre \zh{都}) qui fixe le sens.
\end{enumerate}
\textbf{Barème indicatif (10 pts) :} Q1 : 1 pt ; Q2 : 2 pts ; Q3 : 2 pts ; Q4 : 2 pts ; Q5 : 3 pts (distinction 1 pt + deux exemples traduits 2 pts).
\end{corrige}

\begin{corrige}[Exercice 11 — Thème et version]
\textbf{Thème.}
\begin{enumerate}
\item \zh{除了法语以外，我们还学汉语和英语。} (\emph{Chúle Fǎyǔ yǐwài, wǒmen hái xué Hànyǔ hé Yīngyǔ.})
\item \zh{除了星期天以外，我爸爸每天都工作。} (\emph{Chúle xīngqītiān yǐwài, wǒ bàba měitiān dōu gōngzuò.})
\item \zh{五十加三十等于八十；八十减二十等于六十。} (\emph{Wǔshí jiā sānshí děngyú bāshí ; bāshí jiǎn èrshí děngyú liùshí.})
\end{enumerate}
\textbf{Version.}
\begin{enumerate}
\item « En plus du renminbi, beaucoup de Chinois utilisent aussi le dollar. »
\item « 1000 moins 250 égalent 750. »
\end{enumerate}
\textbf{Barème indicatif (10 pts) :} 2 pts par phrase (structure 1 pt, exactitude des caractères et du pinyin 1 pt).\\
\textbf{Points de vigilance :} « sauf » = \zh{除了……以外，……都……} ; « en plus de » = \zh{除了……以外，……还……} ; place de \zh{每天} avant \zh{都}.
\end{corrige}

\begin{corrige}[Exercice 12 — Problème contextualisé et expression écrite]
\begin{enumerate}
\item Total des achats : \zh{一千五百加六千等于七千五百。} (\emph{Yìqiān wǔbǎi jiā liùqiān děngyú qīqiān wǔbǎi.}) « 1500 + 6000 = 7500 ».
\item Monnaie à rendre : \zh{一万减七千五百等于二千五百。} (\emph{Yíwàn jiǎn qīqiān wǔbǎi děngyú èrqiān wǔbǎi.}) « 10 000 $-$ 7 500 = 2 500 ». On doit te rendre 2 500 FCFA.
\item \zh{在喀麦隆人们用法郎，在中国人们用人民币。} (\emph{Zài Kāmàilóng rénmen yòng fǎláng, zài Zhōngguó rénmen yòng rénmínbì.}) « Au Cameroun on paie en francs (CFA), en Chine on paie en renminbi. »
\end{enumerate}
\textbf{Proposition de rédaction modèle} (\zh{我去市场买东西}, environ 70 caractères) :
\begin{textebox}[Modèle de production écrite]
星期天，我去市场买东西。我买了一个本子，五百法郎，又买了一支笔，三百法郎。五百加三百等于八百，我一共花了八百法郎。我给老板一千法郎，一千减八百等于二百，他找我二百法郎。除了本子和笔以外，我还买了一些水果。今天我很高兴！\\
(Xīngqītiān, wǒ qù shìchǎng mǎi dōngxi. Wǒ mǎi le yí ge běnzi, wǔbǎi fǎláng, yòu mǎi le yì zhī bǐ, sānbǎi fǎláng. Wǔbǎi jiā sānbǎi děngyú bābǎi, wǒ yígòng huā le bābǎi fǎláng. Wǒ gěi lǎobǎn yìqiān fǎláng, yìqiān jiǎn bābǎi děngyú èrbǎi, tā zhǎo wǒ èrbǎi fǎláng. Chúle běnzi hé bǐ yǐwài, wǒ hái mǎi le yìxiē shuǐguǒ. Jīntiān wǒ hěn gāoxìng !)\\
« Dimanche, je vais au marché faire des achats. J'achète un cahier à 500 francs et un stylo à 300 francs. 500 + 300 = 800, je dépense 800 francs en tout. Je donne 1000 francs au vendeur ; 1000 $-$ 800 = 200, il me rend 200 francs. En plus du cahier et du stylo, j'achète aussi des fruits. Aujourd'hui, je suis très content ! »
\end{textebox}
\textbf{Barème indicatif (10 pts) :} calculs en chinois corrects : 3 pts ; phrase sur les monnaies : 2 pts ; texte : 5 pts (présence des trois structures exigées 2 pts, correction grammaticale 2 pts, respect de la longueur 1 pt).\\
\textbf{Points de vigilance :} 10 000 se dit \zh{一万} (le chinois compte par \emph{wàn}, pas par milliers) ; \zh{一共} « en tout » se place avant le verbe ; vérifier que la structure \zh{除了……以外，……还……} est complète (les trois éléments).
\end{corrige}

\newpage

\coursec{Fiche bilan}

\begin{popfigure}
\begin{tikzpicture}[mindmap, grow cyclic, every node/.style={concept}, concept color=popblue!60, text=white,
root concept/.append style={concept color=popdark, font=\bfseries},
level 1/.append style={level distance=3.2cm, sibling angle=60},
level 1 concept/.append style={font=\small\bfseries}]
\node[root concept, align=center]{Leçon 26\\ 加法和减法\\ Nombres, monnaies,\\ 除了……以外}
child[concept color=popblue]{node[align=center]{Addition 加法\\ A 加 B 等于 C}}
child[concept color=poporange]{node[align=center]{Soustraction 减法\\ A 减 B 等于 C}}
child[concept color=popgreen]{node[align=center]{Monnaies\\ 元/块 角/毛 分\\ 人民币 法郎}}
child[concept color=poppurple]{node[align=center]{Nombres\\ 零 一 十 百 千 万\\ 两 vs 二}}
child[concept color=poppink]{node[align=center]{除了……以外\\ ……都/还/也……}}
child[concept color=popgold]{node[align=center]{Erreurs fréquentes\\ ordre des mots,\\ tons, classificateurs}};
\end{tikzpicture}
\legende{Carte mentale de la leçon : calculer, nommer les monnaies et exclure avec 除了……以外}
\end{popfigure}

\begin{bilan}
\textbf{1. Lexique clé à connaître par cœur}
\begin{center}
\begin{tabularx}{\linewidth}{|X|X|X|X|}
\hline
\rowcolor{popblueL} Caractères & Pinyin & Nature & Traduction \\
\hline
加 & jiā & verbe & ajouter, plus \\
\hline
减 & jiǎn & verbe & soustraire, moins \\
\hline
等于 & děngyú & verbe & égaler, faire \\
\hline
得 & dé & verbe & donner (résultat) \\
\hline
人民币 & rénmínbì & nom & monnaie chinoise (RMB) \\
\hline
元 / 块 & yuán / kuài & nom / classif. & yuan (écrit / oral) \\
\hline
角 / 毛 & jiǎo / máo & nom & dixième de yuan \\
\hline
分 & fēn & nom & centième de yuan \\
\hline
法郎 & fǎláng & nom & franc (FCFA) \\
\hline
美元 / 欧元 & Měiyuán / Ōuyuán & nom & dollar / euro \\
\hline
零 & líng & nombre & zéro \\
\hline
两 & liǎng & nombre & deux (devant classificateur) \\
\hline
万 & wàn & nombre & dix mille \\
\hline
除了……以外 & chúle……yǐwài & structure & à part……, sauf…… \\
\hline
便宜 / 贵 & piányi / guì & adjectif & bon marché / cher \\
\hline
一共 & yígòng & adverbe & en tout, au total \\
\hline
\end{tabularx}
\end{center}

\textbf{2. Règles de grammaire à connaître par cœur}
\begin{center}
\begin{tabularx}{\linewidth}{|X|X|X|}
\hline
\rowcolor{popblueL} Structure & Exemple (caractères + pinyin) & Traduction \\
\hline
A + 加 + B + 等于 + C & 三加五等于八。Sān jiā wǔ děngyú bā. & 3 + 5 = 8 \\
\hline
A + 减 + B + 等于 + C & 十减四等于六。Shí jiǎn sì děngyú liù. & 10 − 4 = 6 \\
\hline
nombre + 块 + nombre + 毛 & 三块五 (sān kuài wǔ) & 3,50 yuans \\
\hline
这 + 多少 + 钱 ? & 这个多少钱？Zhège duōshao qián ? & Combien coûte ceci ? \\
\hline
除了 + X + 以外，……都…… & 除了苹果以外，我都买。Chúle píngguǒ yǐwài, wǒ dōu mǎi. & Sauf les pommes, j'achète tout. \\
\hline
除了 + X + 以外，……还/也…… & 除了汉语以外，我还学英语。Chúle Hànyǔ yǐwài, wǒ hái xué Yīngyǔ. & En plus du chinois, j'apprends aussi l'anglais. \\
\hline
\end{tabularx}
\end{center}

\textbf{3. Points critiques}
\begin{itemize}
\item \cle{等于} se place \textbf{avant} le résultat ; jamais d'inversion comme en français (« huit égale trois plus cinq » est impossible).
\item \cle{两} (et non 二) devant un classificateur : \zh{两个人} « deux personnes », mais \zh{二加二} dans les calculs.
\item \cle{除了……以外} : 以外 peut être omis, mais \cle{都} (exclusion) ou \cle{还/也} (addition) est \textbf{obligatoire} dans la seconde partie.
\item Sens double : 除了……都…… = « sauf » ; 除了……还…… = « en plus de ». Le verbe suivant décide du sens.
\item Les grands nombres se comptent par \cle{万} (10 000), pas par milliers : 20 000 = \zh{两万} (liǎng wàn).
\end{itemize}

\textbf{4. Méthode express pour traduire un calcul}
\begin{enumerate}
\item Repérer l'opération (+ ou −) : 加 ou 减.
\item Placer les nombres dans l'ordre : A + opérateur + B + 等于 + résultat.
\item Vérifier les tons : jiā (1\textsuperscript{er} ton), jiǎn (3\textsuperscript{e} ton), děngyú (3\textsuperscript{e}+2\textsuperscript{e}).
\item Pour un prix : nombre + 块/元 (+ nombre + 毛/角) ; à l'oral, on préfère 块 et 毛.
\end{enumerate}
\end{bilan}

\begin{aretenir}[Checklist avant le Bac]
\begin{itemize}
\item[$\square$] Je sais écrire et prononcer 加 (jiā), 减 (jiǎn), 等于 (děngyú) avec les bons tons.
\item[$\square$] Je construis sans hésiter : A + 加/减 + B + 等于 + C.
\item[$\square$] Je connais les unités de monnaie : 元/块, 角/毛, 分, et je dis un prix à l'oral (三块五).
\item[$\square$] Je sais nommer les monnaies : 人民币, 法郎, 美元, 欧元.
\item[$\square$] Je compte jusqu'à 10 000 et j'utilise 万 correctement (两万 = 20 000).
\item[$\square$] Je choisis correctement entre 二 (calculs, numéros) et 两 (devant classificateur).
\item[$\square$] Je maîtrise 除了……以外，……都…… (exclusion = « sauf »).
\item[$\square$] Je maîtrise 除了……以外，……还/也…… (addition = « en plus de »).
\item[$\square$] Je n'oublie jamais 都 ou 还/也 dans la seconde partie de la phrase.
\item[$\square$] Je sais poser la question du prix : 这个多少钱？ et y répondre.
\item[$\square$] Je peux rédiger un court dialogue au marché avec calcul et monnaie (contexte camerounais : FCFA, marché de Yaoundé).
\item[$\square$] Je vérifie la ponctuation chinoise pleine chasse (，。？) dans mes productions écrites.
\end{itemize}
\end{aretenir}

\findecours

\end{document}
